\documentclass[11pt,a4paper]{book}

\usepackage[margin=2.6cm]{geometry}
\usepackage{amsmath,amssymb,amsthm}
\usepackage{stmaryrd}
\usepackage{enumitem}
\usepackage{booktabs}
\usepackage{tikz}
\usepackage{tikz-cd}
\usepackage[most]{tcolorbox}
\usepackage{pifont}
\usepackage[colorlinks=true,linkcolor=blue,citecolor=blue,urlcolor=blue]{hyperref}

%% ---- Notation (kept consistent with the rest of the repo) ----
\newcommand{\Set}{\mathbf{Set}}
\newcommand{\Cat}{\mathbf{Cat}}
\newcommand{\Cof}{\mathbf{Cof}}
\newcommand{\Cont}{\mathbf{Cont}}
\newcommand{\DCont}{\mathbf{DCont}}
\newcommand{\Poly}{\mathbf{Poly}}
\newcommand{\Fam}{\mathbf{Fam}}
\newcommand{\y}{y}
\newcommand{\op}{^{\mathrm{op}}}
\newcommand{\cont}[2]{#1 \triangleleft #2}
\newcommand{\Ext}[1]{\llbracket #1 \rrbracket}
\newcommand{\id}{\mathrm{id}}
\newcommand{\cod}{\mathrm{cod}}
\newcommand{\dom}{\mathrm{dom}}
\newcommand{\ob}{\mathrm{ob}\,}
\newcommand{\sh}{\sharp}
\newcommand{\co}{\mathrel{\text{\reflectbox{$\rightharpoonup$}}}}
\DeclareMathOperator{\Hom}{Hom}
\DeclareMathOperator{\End}{End}
\newcommand{\Cathash}{\mathbb{C}\mathrm{at}^\sharp}
\newcommand{\Mon}{\mathrm{Mon}}
\newcommand{\Comon}{\mathrm{Comon}}
\newcommand{\str}{\mathrm{str}}
\newcommand{\Arr}{\mathsf{Arr}}
\newcommand{\Workers}{\textup{\textsc{Workers}}}

%% ---- Flagging environments (Neil's convention) ----
\theoremstyle{plain}
\newtheorem{Theorem}{Theorem}[chapter]
\newtheorem{Proposition}[Theorem]{Proposition}
\newtheorem{Lemma}[Theorem]{Lemma}
\newtheorem{Corollary}[Theorem]{Corollary}
\newtheorem{Conjecture}[Theorem]{Conjecture}
\theoremstyle{definition}
\newtheorem{Definition}[Theorem]{Definition}
\newtheorem{Example}[Theorem]{Example}
\theoremstyle{remark}
\newtheorem{Remark}[Theorem]{Remark}

%% A small macro for the mandatory provenance tag, so it is visually uniform.
\newcommand{\prov}[1]{\hfill\textsf{\footnotesize[#1]}}

%% Teaching / correction boxes (used from the characterisation chapter on).
\newtcolorbox{teachbox}[1][]{colback=black!4,colframe=black!55,
  fonttitle=\bfseries,title={#1},breakable,left=2mm,right=2mm,top=1mm,bottom=1mm}

\title{\textbf{The Category of Containers}\\[0.4em]
  \large a living reference}
\author{MacBeth}
\date{Draft, June 2026, v0.1}

\begin{document}
\maketitle

\frontmatter
\tableofcontents

\chapter{Preface}
\markboth{PREFACE}{PREFACE}

\section*{Purpose of this document}

This is a single, evolving reference centred on \emph{the category of
containers}. It is meant to grow: chapters are added and revised in place,
and the whole document compiles to one PDF that can be read end to end or
consulted by section. The organising thesis is the equivalence chain
\[
  \textbf{Containers} \ \simeq\ \textbf{Directed Containers}
  \ \simeq\ \textbf{Polynomial Comonoids} \ \simeq\ \textbf{Small Categories},
\]
which we take as the foundation for \emph{verified composition}: once a
system (a supply chain, an agent hierarchy, a state space, a tax computation)
is presented as a container, its composition is governed by the structure of
$\Cont$ and of the comonoids inside it, and correctness of composition becomes
a statement that can be proved --- and, increasingly, machine-checked. Object
level equivalence is settled groundwork; the live mathematics lives at the
\emph{morphism} level, and that is where this book spends most of its effort.

\section*{The flagging convention}

Mathematical honesty is the first principle of this document. Every
\textbf{Definition} and every result is flagged, and every
\textbf{Theorem}/\textbf{Proposition}/\textbf{Lemma}/\textbf{Corollary}
statement \emph{ends with a bracketed provenance tag}, one of:
\begin{itemize}[nosep]
  \item \textsf{[Cited: \textit{author year, ref}]} --- a known result from the
        literature, not claimed as new;
  \item \textsf{[MacBeth]} --- claimed and/or proved by MacBeth, informally
        (not machine-checked);
  \item \textsf{[MacBeth, Lean-verified: \textit{file}]} --- machine-checked in
        the repository's \texttt{lean/} tree, in the named source file;
  \item \textsf{[Conjecture, MacBeth]} or \textsf{[Open]} --- for conjectures.
\end{itemize}
A tag \textsf{[Cited:\,\ldots]} on a statement that MacBeth has also verified
independently will say so; when in doubt the weaker (cited, not-checked) tag is
used, and a footnote records the uncertainty. The bias is deliberately towards
\emph{under}-claiming. The reading and citation tracker in
Appendix~\ref{app:tracker} records, paper by paper, how thoroughly each source
has been engaged.

A note on the Lean tags. Only two Lean \emph{source} files are committed to the
repository at the time of this draft: \texttt{Basic.lean} (the container
representation theorem) and \texttt{Directed.lean} (the comonad laws from the
directed-container laws). Results checked in those files carry the
\textsf{Lean-verified} tag naming the committed file. Where MacBeth has a Lean
formalisation that is \emph{not} yet committed as source to this repository
(notably the $\DCont\cong\Cof$ and Zappa--Sz\'ep associativity developments),
the statement is tagged \textsf{[Cited]} for the underlying theorem and a
footnote flags that the Lean source is not in this tree --- so the reader is
never told something is machine-checked here when its source cannot be found
here.

\mainmatter

\chapter{Containers and their extension functor}
\label{ch:cont}

\section{Containers}

\begin{Definition}[Container]\label{def:container}
A \emph{container} is a pair $\cont SP$ where $S$ is a set of \emph{shapes} and
$P\colon S\to\Set$ is a family of \emph{positions}, one set $P(s)$ for each shape
$s$. Its \emph{extension} (or \emph{interpretation}) is the functor
$\Ext{\cont SP}\colon\Set\to\Set$ defined on objects by
\[
  \Ext{\cont SP}(X) \;=\; \coprod_{s\in S} X^{P(s)}
    \;=\; \bigl\{ (s,v) \mid s\in S,\ v\colon P(s)\to X \bigr\},
\]
and on a morphism $f\colon X\to Y$ by postcomposition,
$\Ext{\cont SP}(f)(s,v) = (s,\, f\circ v)$.
\end{Definition}

Unpacking: a shape $s$ is a structure with its positions $P(s)$; a labelling
$v\colon P(s)\to X$ fills every position with an element of $X$; the coproduct
over $s$ is the disjoint union ``pick exactly one shape''. The functorial action
relabels --- it changes the data at the positions but never the shape, so the
\emph{shape projection} $\pi(s,v)=s$ is a natural transformation to the constant
functor at $S$. A container thus separates \emph{structure} (the shape) from
\emph{content} (the labels), and its functorial action is exactly \texttt{map}.

\begin{Example}[The container menagerie]\label{ex:menagerie}
\textbf{Maybe}: two shapes, ``nothing'' ($P=\varnothing$, term $X^0=1$) and
``just'' ($P=\{*\}$, term $X$); extension $1+X$.
\textbf{Lists}: shape $n\in\mathbb N$, positions $\{0,\dots,n-1\}$; extension
$\sum_{n} X^n$.
\textbf{Binary trees}: a shape is a planar binary tree, positions are leaves;
the generating function satisfies $T(X)=1+X\cdot T(X)^2$, so the number of shapes
with $n$ positions is the Catalan number $C_n=\frac1{n+1}\binom{2n}{n}$.
\textbf{Matrices}: shape $(m,n)$, positions the $mn$ entries; extension
$\sum_{m,n} X^{m\times n}$. In each case the coefficient of $X^n$ in the
generating function counts \emph{how many} shapes have $n$ positions, while the
container remembers \emph{which} shapes and \emph{where} the positions are.
\end{Example}

A standing intuition (``decategorification''): writing $\sum_s X^{P(s)}$ as if it
were a polynomial in a formal indeterminate $X$ guides the right intuitions about
products, composition and derivatives. But $X$ is a \emph{set}, not an
indeterminate: $X^{P(s)}=\Hom(P(s),X)$ is a set of functions, and the extension
lives in $[\Set,\Set]$, not in a ring of power series.

\section{The representation theorem}

The combinatorial data of containers captures \emph{exactly} the natural
transformations between their extensions.

\begin{Theorem}[Representation / full faithfulness]\label{thm:rep}
The extension functor $\Ext{-}\colon\Cont\to[\Set,\Set]$ is full and faithful:
every natural transformation $\Ext{\cont{S_1}{P_1}}\Rightarrow\Ext{\cont{S_2}{P_2}}$
arises from a unique container morphism, and distinct morphisms induce distinct
natural transformations.
\prov{Cited: Abbott--Altenkirch--Ghani 2005}%
\footnote{Independently formalised by MacBeth in the committed Lean file
\texttt{lean/Containers/Containers/Basic.lean} (round-trip
\texttt{ContainerMorphism}\,$\leftrightarrow$\,\texttt{ExtNatTrans}, no
\texttt{sorry}).}
\end{Theorem}

\begin{proof}[Proof sketch]
By the Yoneda lemma a natural transformation $\alpha$ is determined by its action
on the ``generic element'' $(s,\id_{P_1(s)})\in\Ext{\cont{S_1}{P_1}}(P_1(s))$ of
each shape: $\alpha_{P_1(s)}(s,\id)=(u(s),f_s)$ reads off both a shape map
$u\colon S_1\to S_2$ and a position map $f_s\colon P_2(u(s))\to P_1(s)$, and
naturality forces every other component. Conversely every $(u,f)$ is natural
(Chapter~\ref{ch:contcat}). See \cite{AAG05} for the full argument; the
round-trip is machine-checked in \texttt{Basic.lean}.
\end{proof}

\chapter{The category \texorpdfstring{$\Cont$}{Cont}}
\label{ch:contcat}

\section{Container morphisms: the variance}

\begin{Definition}[Container morphism]\label{def:contmor}
A \emph{container morphism} $(u,f)\colon\cont{S_1}{P_1}\to\cont{S_2}{P_2}$ consists
of
\begin{itemize}[nosep]
  \item a \emph{shape map} $u\colon S_1\to S_2$ (\emph{forward}), and
  \item for each $s\in S_1$, a \emph{position map}
        $f_s\colon P_2(u(s))\to P_1(s)$ (\emph{backward}).
\end{itemize}
It induces the natural transformation $\alpha_X(s,v)=\bigl(u(s),\, v\circ f_s\bigr)$.
\end{Definition}

The variance is not a convention but a necessity. To build a new labelling of the
\emph{output} positions $P_2(u(s))$ out of an old labelling $v\colon P_1(s)\to X$
without inspecting labels, the only move is to precompose:
\[
  P_2(u(s)) \xrightarrow{\ f_s\ } P_1(s) \xrightarrow{\ v\ } X,
\]
which forces $f_s$ to run from output positions to input positions. Each output
position $p'$ asks ``which input position supplies my label?'' and the answer is
$f_s(p')$. Consequently a container morphism can only \emph{rearrange},
\emph{discard} or \emph{duplicate} labels --- never \emph{create} one. (Sorting,
filtering and thresholding are \emph{not} container morphisms: they inspect
labels, and so are not natural in $X$.)

\begin{Definition}[Identity and composition]\label{def:contcomp}
The identity on $\cont SP$ is $(\id_S,\id)$. The composite of
$(u,f)\colon\cont{S_1}{P_1}\to\cont{S_2}{P_2}$ and
$(u',f')\colon\cont{S_2}{P_2}\to\cont{S_3}{P_3}$ has shape map $u'\circ u$ and
position map running backward through both stages,
\[
  (u'\circ u)^\sharp_s \;=\; f_s\circ f'_{u(s)}\colon
   P_3\bigl(u'(u(s))\bigr) \xrightarrow{\ f'_{u(s)}\ } P_2(u(s))
   \xrightarrow{\ f_s\ } P_1(s).
\]
These satisfy the category axioms; the resulting category is $\Cont$.
\end{Definition}

\begin{Proposition}[Naturality and the correspondence]\label{prop:nat}
For every container morphism $(u,f)$, the family $\alpha_X(s,v)=(u(s),v\circ f_s)$
is natural in $X$; and (Theorem~\ref{thm:rep}) every natural transformation
between container extensions is $\alpha$ for a unique $(u,f)$. Hence
$\Ext{-}\colon\Cont\to[\Set,\Set]$ is an isomorphism onto its image.
\prov{Cited: AAG05}
\end{Proposition}

\begin{proof}
Naturality is associativity of composition: for $g\colon X\to Y$, both ways round
the naturality square send $(s,v)$ to $(u(s),\, g\circ v\circ f_s)$, because
$f_s$ does not look at labels. Full faithfulness is Theorem~\ref{thm:rep}.
\end{proof}

\section{Worked example: \texttt{tail}}

The function $\mathrm{tail}\colon\mathrm{List}(X)\to 1+\mathrm{List}(X)$ is a
container morphism $(u,f)$: the shape map sends length $n$ to length $n-1$ (and
$0$ to ``nothing''), and the (backward) position map sends position $i$ of the
output to position $i+1$ of the input. The same $(u,f)$ acts on lists of integers,
strings, or anything --- the labels are shifted one place left, never inspected.
This is the content of naturality: the morphism is parametric in the label type.

\section{An application vista: Weihrauch problems}
\label{sec:weihrauch}

Before leaving the bare category, it is worth flagging one place where these two
ingredients --- a set of shapes with a family of positions, and morphisms that run
forward on shapes and backward on positions --- have recently turned up far from
data structures. In computable analysis a \emph{Weihrauch problem} is a multi-valued
task: a set of \emph{instances}, and for each instance a set of admissible
\emph{solutions}. That is precisely a container --- instances are shapes, solutions
are positions --- and a \emph{Weihrauch reduction}, the standard notion of ``problem
$A$ is no harder than problem $B$'', is exactly a container morphism: forward on
instances (translate an $A$-instance to a $B$-instance), backward on solutions
(translate a $B$-solution back to an $A$-solution), the same variance for the same
reason. Pradic and Price make this identification precise \cite{PradicPrice}: their
category of problems is a category of containers over a base of realizability
assemblies, and --- pleasingly for us --- their sequential composition of problems is
written with the very symbol $\triangleleft$ we will give to the composition product
in Chapter~\ref{ch:algebra}, the Weihrauch star. Their base is only \emph{weakly}
locally cartesian closed, so their composition is a mere quasi-bifunctor; the theory
of this book lives over $\Set$, which is genuinely locally cartesian closed, and one
lesson of the comparison is to see exactly where that hypothesis earns its keep.
We record the connection as a vista, not a thread we pursue here: the container
calculus that follows is, unexpectedly, also a calculus of computational hardness.

\chapter{Which functors are containers?}
\label{ch:which}

Chapter~\ref{ch:cont} proved that the extension functor
$\Ext{-}\colon\Cont\to[\Set,\Set]$ is full and faithful
(Theorem~\ref{thm:rep}): the combinatorics of containers captures \emph{exactly}
the natural transformations between their extensions. Full and faithful means
$\Ext{-}$ is an isomorphism \emph{onto its image}. It does not tell us what that
image \emph{is}. This chapter closes the gap. It answers, exactly, the question a
reader should have been asking since the definition:

\begin{center}
\itshape Hand me a functor $F\colon\Set\to\Set$. Is it a container?
\end{center}

\paragraph{Why the question has teeth.} It is not idle taxonomy. Look ahead to the
cofree comonad of Chapter~\ref{ch:moncomon}: to build the behaviours of a system ---
its infinite unfoldings --- one forms the \emph{terminal sequence}
\[
  1 \;\longleftarrow\; F1 \;\longleftarrow\; FF1 \;\longleftarrow\; FFF1
    \;\longleftarrow\;\cdots
\]
and takes its limit, where the final coalgebra lives. That limit is a
\emph{cofiltered} one, and the whole construction works only if $F$
\emph{preserves} it. So before we can compose behaviours we are forced to ask which
functors preserve limits of this kind. The answer, we will see, is precisely: the
containers. The equivalence chain, the comonoids, the monoidal structures of the
later chapters --- the entire programme of this book lives \emph{inside} the image
of $\Ext{-}$, so knowing which functors lie in that image is knowing the exact
reach of the theory: what it covers, and where its boundary runs.

\paragraph{The promise.} By the end of the chapter you will be able to look at a
functor and decide whether it is a container; you will know how to read its shapes
and positions back off; and you will see why the first formula anyone guesses for
the positions is wrong. The tools are one structural identification and one theorem
--- the \emph{wide-pullback} (connected-limit) characterisation --- and both are
best met after a little computation.

\section{A container is a family of sets}
\label{sec:fam}

Before we test functors, we should see $\Cont$ for what it is. Strip the extension
away and look at the raw data of a container: a set $S$, and for each $s\in S$ a set
$P(s)$. That is an $S$-indexed \emph{family} of sets. Nothing more.

Families of objects of a category $\mathcal C$ are themselves the objects of a
category, the \emph{free coproduct completion} $\Fam(\mathcal C)$. Its objects are
pairs $(S,\,P\colon S\to\mathcal C)$; a morphism $(S,P)\to(T,Q)$ is a reindexing
function $u\colon S\to T$ together with, for each $s\in S$, a $\mathcal C$-morphism
comparing $P(s)$ with $Q(u(s))$. The completion turns coproducts that $\mathcal C$
may lack into formal ones: the family $(S,P)$ \emph{is} the formal coproduct
$\coprod_{s\in S}P(s)$.

Now watch the variance. A container morphism $\cont SP\to\cont TQ$, as
Definition~\ref{def:contmor} had it, is a shape map $u\colon S\to T$ \emph{forward}
together with position maps $f_s\colon Q(u(s))\to P(s)$ running \emph{backward}.
Backward is exactly a morphism in $\Set\op$. So the fibrewise comparison lives in
$\Set\op$, not $\Set$, and the category is identified on the nose.

\begin{Proposition}[$\Cont$ is the free coproduct completion of $\Set\op$]
\label{prop:fam}
There is an isomorphism of categories $\Cont\cong\Fam(\Set\op)$, the identity on
objects, sending a container morphism $(u,f)$ to the reindexing $u$ together with
the family $(f_s)_{s\in S}$ read as morphisms in $\Set\op$.
\prov{Folklore}
\end{Proposition}

This is not deep, but it is clarifying, and it pays for itself immediately. It is a
standard fact that the free coproduct completion of any category $\mathcal C$
embeds into $[\mathcal C\op,\Set]$ by the left Kan extension of the Yoneda embedding
along itself --- concretely, a family becomes a coproduct of representables. Feed it
$\mathcal C=\Set\op$ and the abstract statement becomes our extension functor:
\[
  (S,P) \;\longmapsto\; \coprod_{s\in S}\y^{P(s)}, \qquad
  \y^{A} = \Hom_{\Set}(A,-),
\]
because a representable on $\Set\op$ evaluated in $\Set$ is the covariant
hom-functor $\y^A=(-)^A$. So $\Ext{\cont SP}=\coprod_s\y^{P(s)}$ is not a definition
we imposed; it is what the free coproduct completion of $\Set\op$ \emph{does} when
you land it in functors. Containers are exactly \emph{coproducts of representable
endofunctors of $\Set$}. Hold that phrase --- it is the whole answer to our
question, once we understand which functors it describes.

\begin{teachbox}[Terminology hazard: positions versus directions]
The reader who also reads Spivak's \emph{Poly} must keep two dictionaries. There a
polynomial functor is written $p=\sum_{i\in p(1)}\y^{p[i]}$, the elements of $p(1)$
are called \emph{positions}, and the elements of each fibre $p[i]$ are called
\emph{directions}. In our language the elements of $p(1)=S$ are \emph{shapes} and
the elements of the fibres $P(s)$ are \emph{positions}. The naming is exactly
\emph{opposite}: Spivak's ``positions'' are our shapes, his ``directions'' our
positions. The objects are the same; only the words differ. We keep
shapes-and-positions throughout, because in the applications a shape is a structure
and a position is a slot inside it, and that is the intuition we want the words to
carry.
\end{teachbox}

\section{The test}
\label{sec:test}

We now know the target shape: a container extension is a coproduct of
representables, $F(X)\cong\coprod_{s\in S}X^{P(s)}$. So ``is $F$ a container?'' is
``can $F$ be written that way?'' We gather evidence by computing: the arithmetic of
finite sets decides several cases outright and --- better --- teaches the method.

Here is the fact that does the work. If $F=\Ext{\cont SP}$ and $X$ is finite with
$|X|=n$, then
\[
  |F(n)| \;=\; \sum_{s\in S}\, n^{\,|P(s)|}.
\]
The cardinality of a container extension at $n$ is a sum of powers of $n$, one power
per shape, the exponent being the number of positions. Small $n$ probe the sum. At
$n=1$ every power is $1$, so $|F(1)|=|S|$: the value at a one-element set counts the
shapes. At $n=0$ every positive power vanishes and $0^0=1$, so $|F(0)|$ counts the
shapes with \emph{no} positions --- the constants. These two readings alone are
often enough.

\begin{Example}[Friendly functors pass]\label{ex:pass}
$F(X)=1+X$: shapes $\{0,1\}$ with $|P|=0,1$; extension $\coprod\{X^0,X^1\}$. A
container. $F(X)=X^2$: one shape, two positions. A container. Lists,
$F(X)=\coprod_n X^n$: shape $n\in\mathbb N$, positions $\{0,\dots,n-1\}$. A
container. In each case the coproduct-of-powers form is visible on the page.
\end{Example}

\begin{Example}[The covariant powerset fails --- in one line]\label{ex:powerset}
Let $\mathcal P\colon\Set\to\Set$ be the covariant powerset functor, $\mathcal
P(X)=\{\text{subsets of }X\}$, acting on maps by direct image. Is it a container?

Count. $|\mathcal P(0)|=1$, $|\mathcal P(1)|=2$, $|\mathcal P(2)|=4$, in general
$|\mathcal P(n)|=2^n$. Suppose $\mathcal P=\Ext{\cont SP}$. From $n=1$ we read
$|S|=|\mathcal P(1)|=2$: two shapes. From $n=0$ we read that exactly one of them has
no positions; call the other shape $s_\star$, with $|P(s_\star)|=p$. Evaluate at
$n=2$:
\[
  4 \;=\; |\mathcal P(2)| \;=\; \sum_{s\in S} 2^{|P(s)|}
        \;=\; 2^{0}+2^{p} \;=\; 1+2^{p}.
\]
So $2^p=3$, which no natural number $p$ satisfies. The covariant powerset is
\emph{not} a container.
\end{Example}

The powerset is worth pausing on, because the reason it fails is the reason the
theory has a boundary at all. A subset of $X$ is a collection of elements with no
record of \emph{where} each element sits. A container labelling $v\colon P(s)\to X$,
by contrast, assigns each position its own element, and the position is a name that
persists. Powerset forgets the names; it is a functor built by \emph{quotienting},
and quotients are exactly what containers cannot do. We meet this boundary again in
\S\ref{sec:limits}, from the colimit side, and it is the same wall.

Counting settles individual cases, but we want a criterion that works for every
functor, finite or not. Here it is. The right notion is that of a \emph{connected
limit}: a limit whose indexing diagram is connected (any two objects joined by a
zig-zag of arrows). A pullback is connected --- its shape
$\bullet\to\bullet\leftarrow\bullet$ is joined up --- and so are \emph{wide
pullbacks} (many arrows into one target) and \emph{cofiltered limits}, the terminal
sequence of the chapter's opening among them. A product is \emph{not} connected: its
shape is a bare set of objects with no arrows between them, so it falls into
disconnected pieces. Connected versus disconnected is exactly the line the
characterisation runs along.

\begin{Theorem}[Characterisation of containers]\label{thm:char}
For a functor $F\colon\Set\to\Set$ the following are equivalent.
\begin{enumerate}[label=\textup{(\roman*)},nosep]
\item $F$ is a container extension: $F\cong\Ext{\cont SP}$ for some container.
\item $F$ preserves connected limits (equivalently: wide pullbacks and cofiltered
      limits).
\item $F$ is a coproduct of representable functors.
\end{enumerate}
\prov{Cited: Gambino--Kock 2009, \S1.18 \& Prop.~1.22; equivalence due to Diers
1977 and Carboni--Johnstone 1995}
\end{Theorem}

Why is preserving connected limits the same as being a coproduct of representables?
The intuition is short and worth carrying. A single representable $\y^A=(-)^A$
preserves \emph{all} limits: it is the right adjoint to $-\times A$, and right
adjoints preserve limits. The question is what survives when we glue representables
by a coproduct. In $\Set$, coproducts commute with connected limits but not with
arbitrary ones: a connected diagram cannot ``see'' the disjoint branches of a
coproduct separately, so the coproduct passes through the limit, whereas a product
(disconnected) can pair up elements across different branches and breaks the
interchange. Preserving connected limits is therefore precisely the trace left in
$[\Set,\Set]$ by ``being a coproduct of things that preserve everything.'' Condition
(ii) is the shadow of condition (iii).

\begin{teachbox}[Why ``connected'', and neither ``wide pullbacks'' nor ``all'']
Condition (ii) is easy to mis-state in either direction.

\emph{Not merely wide pullbacks.} Wide pullbacks are the connected limits one meets
most often, so one is tempted to shorten (ii) to ``preserves wide pullbacks''. That
is genuinely weaker. Cofiltered limits are connected too --- the terminal sequence
$1\leftarrow F1\leftarrow FF1\leftarrow\cdots$ that opened the chapter is one --- yet
they are not built from pullbacks, and must be demanded separately. This is exactly
why the cofree comonad of Chapter~\ref{ch:moncomon} needs the full strength of (ii):
its final-coalgebra limit is cofiltered, not a wide pullback. ``Preserves connected
limits'' unpacks to wide pullbacks \emph{and} cofiltered limits together, and
dropping the latter loses functors \prov{Cited: Gambino--Kock, \S1.18}.

\emph{Not all limits either.} A container need \emph{not} preserve the \emph{empty}
limit --- the terminal object $1$ --- because $\Ext{\cont SP}(1)=S$ can be any set,
whereas preserving the empty limit would force $F(1)\cong1$. The empty diagram is
\emph{disconnected} --- a connected category is required to be nonempty, so the
empty one fails the definition at the first clause --- and it sits rightly outside
(ii). Connectedness is exactly the dividing
line: it keeps the cofiltered limits, which containers respect, and discards the
products and the empty limit, which they do not. A one-line sanity check makes this
memorable: whenever you claim ``$F$ preserves connected limits'', test it against
$F(1)$; if $F(1)\neq1$ then the terminal object is \emph{not} among the preserved
limits, confirming that ``connected'' really must exclude it.
\end{teachbox}

\section{Reading the positions back off}
\label{sec:recover}

The test tells us \emph{whether} $F$ is a container. Suppose it passes. How do we
recover the data $\cont SP$?

The shapes are free: $S=F(1)$, exactly as the counting used. A one-element set has
one labelling per shape, so $F(1)=\coprod_s 1^{P(s)}=\coprod_s 1=S$. That half is
honest arithmetic. The positions are the subtle half, and here a natural first
guess goes wrong.

\begin{teachbox}[Correction: the fibre of $F(2)\to F(1)$ is \emph{not} $P(s)$]
The tempting formula runs: the unique map $2\to1$ induces $F(2)\to F(1)=S$; take the
fibre over a shape $s$ and call it the positions. Compute what this actually
returns. Since $F(2)=\coprod_{s'}2^{P(s')}$ and the map to $S$ forgets the labelling
$P(s')\to2$ and remembers only $s'$, the fibre over $s$ is
\[
  \bigl\{\,\varphi\colon P(s)\to 2\,\bigr\} \;=\; 2^{P(s)},
\]
the set of \emph{subsets} of $P(s)$ --- the \emph{powerset} of the positions, not
the positions. Evaluation-and-fibre overcounts by an exponential. No single
evaluation of $F$, followed by a fibre, can extract $P(s)$: the functor at a fixed
set only ever sees \emph{labellings} of the positions, never the positions bare.
\end{teachbox}

The correction is not a patch; it is a signpost pointing at the real mechanism. To
see $P(s)$ itself we must ask for the labelling that names each position \emph{by
itself} --- the identity labelling --- and that is a universal property, not an
evaluation.

\begin{Proposition}[Position recovery via the generic element]\label{prop:recover}
Let $F$ be a container. For each shape $s\in F(1)$ there is a set $P(s)$ and a
\emph{generic element} $g_s\in F(P(s))$ lying over $s$, universal in the sense that
for every set $X$ and every $x\in F(X)$ lying over $s$ there is a unique map
$\alpha\colon P(s)\to X$ with $F(\alpha)(g_s)=x$. This $P(s)$ is the object of
positions, and
\[
  F(X)\;\cong\;\coprod_{s\in F(1)} X^{P(s)}
\]
is the container form, recovered componentwise by the Yoneda lemma. The generic
element is the initial object of the category of elements of $F$ sitting over $s$;
its existence is exactly the connected-limit preservation of
Theorem~\ref{thm:char}.
\prov{Cited: Gambino--Kock 2009, Prop.~1.22}
\end{Proposition}

The generic element is the identity labelling $g_s=(s,\id_{P(s)})$ wearing abstract
clothing: it is the one element of $F(P(s))$ from which every other labelled
structure of shape $s$ is obtained by relabelling, and by exactly one relabelling.
``Exactly one'' is the universal property, and it is what pins $P(s)$ down to the
positions rather than to any of their exponential shadows.

\begin{Remark}[The derivative sees the total position set]
A second, tidier-looking phrasing is worth knowing --- and worth not overtrusting.
The total space of positions $\coprod_s P(s)$ is the value at $1$ of the
\emph{derivative} $\partial F$, the container of one-hole contexts sketched in the
Phase~2 outline (Chapter~\ref{ch:phase2}), and $P(s)$ is its fibre over $s$. This is
correct, but it is a restatement, not a shortcut: $\partial F$ is itself a derived
construction, not a single evaluation of $F$, so it does not escape the moral of the
correction box. If you want the positions, you must, one way or another, invoke the
universal property.
\end{Remark}

\section{Limits and colimits in $\Cont$}
\label{sec:limits}

We have been studying $\Ext{-}$ as a bridge from $\Cont$ into $[\Set,\Set]$. How
much categorical structure does the bridge carry across? The identification of
\S\ref{sec:fam} answers the first half at a stroke: free coproduct completions are
complete and cocomplete, so \emph{$\Cont$ is both}. What is delicate is which limits
and colimits $\Ext{-}$ \emph{preserves}, and the characterisation theorem has
already told us where the answer turns.

Start with the structure that is easy, concrete, and machine-checked. Binary
products and coproducts in $\Cont$ have a clean description, and the variance is the
whole lesson.

\begin{Proposition}[Products and coproducts of containers]\label{prop:prodcoprod}
In $\Cont$:
\begin{itemize}[nosep]
\item the \emph{coproduct} of $\cont SP$ and $\cont TQ$ has shapes $S+T$ and
      positions inherited case-wise ($P$ on the left summand, $Q$ on the right);
\item the \emph{product} of $\cont SP$ and $\cont TQ$ has shapes $S\times T$ and,
      over $(s,t)$, positions $P(s)+Q(t)$ --- a \emph{coproduct} of position sets.
\end{itemize}
On extensions, $\Ext{-}$ carries both to the pointwise coproduct and product of
functors. (These reappear, framed as the monoidal $+$ and $\times$, in
Chapter~\ref{ch:algebra}.)
\prov{MacBeth}%
\footnote{MacBeth has a Lean formalisation of both universal properties in the file
\texttt{Cont.lean} (built, zero \texttt{sorry}); that source is \emph{not yet
committed} to this repository's \texttt{lean/} tree, so the tag is \textsf{[MacBeth]}
here rather than \textsf{Lean-verified}, matching the convention of the Preface.}
\end{Proposition}

The product deserves a second look, because it inverts the naive expectation. The
positions of a product are a \emph{sum}, not a product. The moment you accept
$\Cont\cong\Fam(\Set\op)$ (\S\ref{sec:fam}) this is not a surprise but a
\emph{compulsion}: a product in $\Fam(\mathcal C)$ is computed fibrewise in
$\mathcal C$, and a product in $\Set\op$ is a coproduct in $\Set$, so the fibres of
a product of containers can only be coproducts. The same fact seen closer up is the
variance of \S\ref{sec:fam} playing out on universal properties: the projection
$\pi_1\colon\cont SP\times\cont TQ\to\cont SP$ has the forward shape map
$(s,t)\mapsto s$, and on positions it must run \emph{backward}, so it is the
coproduct injection $P(s)\hookrightarrow P(s)+Q(t)$. A morphism into the product is
a pair of morphisms, and pairing forward shape maps forces summing backward position
maps. Under $\Ext{-}$ this reads as the familiar polynomial identity
$\Ext{\cont SP\times\cont TQ}(X)=\coprod_{s,t}X^{P(s)+Q(t)}
=\coprod_{s,t}X^{P(s)}\times X^{Q(t)}$, the product of the two extensions --- so
$\Ext{-}$ preserves this product. It is worth being clear about why that is a
separate fact and not a corollary of the characterisation theorem.

\begin{teachbox}[What $\Ext{-}$ preserves, and where it stops]
Theorem~\ref{thm:char} says a container extension \emph{preserves connected limits}.
Products are the archetypal \emph{dis}connected limit, so connected-limit
preservation says nothing about them directly; nevertheless $\Ext{-}$ does preserve
products, by the explicit computation above (and coproducts, since a coproduct of
shapes maps to a coproduct of extensions). The clean summary:
\[
  \Ext{-}\ \text{preserves}\ \begin{cases}
    \text{connected limits} & \text{(Theorem~\ref{thm:char})}\\
    \text{products and coproducts} & \text{(Proposition~\ref{prop:prodcoprod})}
  \end{cases}
\]
but $\Ext{-}$ is \emph{not} cocontinuous: it does not preserve general colimits. The
first thing it fails to preserve is the \emph{coequaliser}.
\end{teachbox}

Why should coequalisers be where preservation breaks? Because a coequaliser is a
quotient, and we have already watched a quotient carry an object out of the
container world once. The covariant powerset of \S\ref{sec:test} is itself a
quotient of a coproduct of representables: it is the list functor with its
labellings identified whenever they differ by a rearrangement or a repetition. That
double identification is exactly what a coproduct of representables cannot express
--- a representable names its positions and keeps them distinct --- and it is why
the counting failed. Coequalisers reproduce the phenomenon in general: coequalise
two container morphisms and the resulting functor can acquire automorphisms of its
positions. The mildest such quotients, by a group acting on positions, are the
\emph{analytic} functors of Joyal's species, already outside the image of
$\Ext{-}$; powerset lies further out still, its idempotent identification coarser
than any group action. Either way the lesson is one line: the
coproduct-of-representables form is closed under connected limits and under
coproducts, and not under quotients.

The failure is not merely generic: there is a minimal two-container witness, and it
is exactly the swap that turns ordered pairs into unordered ones.

\begin{Proposition}[$\Ext{-}$ does not preserve coequalisers]\label{prop:coeq-fail}
$\Cont$ has coequalisers whenever the base category has equalisers and coequalisers,
but $\Ext{-}$ does not preserve them. For a witness, let
$\mathsf{Pair}=\cont 1{\{1,2\}}$ be the container of ordered pairs, with
$\Ext{\mathsf{Pair}}X=X^{2}$, and let
$\id,\ \sigma\colon\mathsf{Pair}\rightrightarrows\mathsf{Pair}$ be the identity and
the swap --- both the single shape map, with $\sigma$'s backward position map
transposing $1\leftrightarrow2$. Their coequaliser in $\Cont$ is $\cont 1\varnothing$:
one shape, and its positions are the \emph{equaliser} of $\id,\sigma\colon\{1,2\}
\rightrightarrows\{1,2\}$ in $\Set$ (positions run backward, so a coequaliser of
shapes is an equaliser of positions), which is empty since the swap has no fixed
point. Hence $\Ext{\cont 1\varnothing}=X^{\varnothing}$ is the constant functor $1$.
But the coequaliser of $\id,\sigma$ in $[\Set,\Set]$ is the \emph{unordered}-pair
functor $X\mapsto X^{2}/\!\!\sim$, $(x,y)\sim(y,x)$, which is not constant. The two
disagree, so $\Ext{-}$ carries neither to the other: it does not preserve this
coequaliser.
\prov{Cited: Abbott--Altenkirch--Ghani, \emph{Categories of containers} (FoSSaCS 2003)
\cite{AAGcat}, Prop.~4.3 ($\Cont$ has coequalisers when the base has equalisers and
coequalisers) and Ex.~4.4 (this witness); Ex.~4.5 treats filtered colimits the same way.}
\end{Proposition}

The witness is the mechanism of the paragraph above at its smallest: quotienting by
the swap introduces a position automorphism (the transposition), and the extension of
the quotient --- the unordered-pair functor --- is exactly one of the analytic
functors that lie outside the image of $\Ext{-}$.

\begin{Remark}[$\Cont$ is complete and cocomplete; $\Ext{-}$ is not cocontinuous]
$\Cont\cong\Fam(\Set\op)$ is complete and cocomplete, so \emph{$\Cont$ has all
limits and all colimits}. The subtlety is entirely about which the bridge $\Ext{-}$
transports: connected limits and (co)products, yes; coequalisers, no
(Proposition~\ref{prop:coeq-fail}). This is the first place the polynomial world and
the functor world part company, and it is worth holding as the boundary marker it is.
\prov{MacBeth; the coequaliser half is Proposition~\ref{prop:coeq-fail}.}
\end{Remark}

\section{Two witnesses that are genuine (co)monads}
\label{sec:two-witnesses}

The coequaliser of the last section carried an object out of the container world, and
the object it produced --- the unordered-pair functor --- was a bare functor, assembled
by hand as a colimit. A sceptic could still hope the boundary is an artefact of that
assembly: perhaps every functor one actually \emph{uses} --- every functor that carries
algebraic structure, that is a monad or a comonad --- stays inside. This section closes
that hope, and closes it symmetrically. There is a bona fide \emph{monad} on $\Set$ that
is not a container, and, mirroring it almost line for line, a bona fide \emph{comonad} on
$\Set$ that is not a container. Both fail the container test in \emph{exactly the same
way} --- they do not preserve the connected pullback $A\to 1\leftarrow B$ --- and both
fail it for the same reason: each is a \emph{quotient} of a polynomial structure by a
symmetry, and the quotient throws away the very provenance a container must remember.

Take the monad first, because you can hold it in your hand.

\subsection*{The monad side: the multiset monad is not a container}

The finite-multiset monad $\mathsf{Bag}$ sends a set $X$ to the set $\mathsf{Bag}\,X$ of
finite multisets of elements of $X$, with unit $\eta\,x=\{x\}$ and multiplication
$\mu=\uplus$ (multiset union). It is a genuine monad --- the free commutative monoid
monad --- and at first sight it looks exactly like a container. A multiset
$m=\{x_{1},\dots,x_{n}\}$ has $n$ \emph{leaves}, labelled by the $x_{i}$, and
$\mathsf{Bag}\,f$ relabels the leaves without disturbing their number; so $\mathsf{Bag}$
is leaf-supported, just as $\Ext{\cont SP}X=\coprod_{s}X^{P(s)}$ presents each element by
a shape and a labelling of positions. The multiplication is even better behaved than one
could ask: $\mu$ merges the inner multisets by disjoint union, so its effect on leaves is
a \emph{bijection} --- no leaf of $\mu\,mm$ appears from nowhere, every one has a unique
inner token of the same label. Every pointwise reason to be a container is present.

And yet $\mathsf{Bag}$ is not a container. The test of \S\ref{sec:test} is preservation of
connected limits, and the cheapest connected limit to probe is the pullback of the cospan
$A\to 1\leftarrow B$, whose limit is $A\times B$. Preservation asks that the comparison map
\[
  \mathsf{Bag}(A\times B)\ \longrightarrow\ \mathsf{Bag}\,A\ \times_{\mathsf{Bag}\,1}\ \mathsf{Bag}\,B
\]
be a bijection. Here is the point that makes this a real test rather than a triviality:
$\mathsf{Bag}\,1=\mathbb N$, the set of \emph{sizes}, so the right-hand side is not a plain
product but the pullback of \emph{equal-size} multisets --- a genuinely connected limit
mediated by $\mathbb N$. Now count, with $A=B=\{0,1\}$, at size $2$. On the left, the
size-$2$ multisets over the four-element set $A\times B$ number $\binom{4+2-1}{2}=10$. On
the right, a size-$2$ multiset over $A$ paired with a size-$2$ multiset over $B$: there are
three of each, so $3\times 3=9$. Ten does not equal nine, and the comparison map cannot be a
bijection.

\begin{Proposition}[The multiset monad is not a container]\label{prop:bag-not-container}
$\mathsf{Bag}$ does not preserve the connected pullback $A\to 1\leftarrow B$, so it is not
the extension of any container. Concretely, at $A=B=\{0,1\}$ the two distinct size-$2$
multisets
\[
  w_{1}=\{(0,0),(1,1)\},\qquad w_{2}=\{(0,1),(1,0)\}\quad\text{over }A\times B
\]
have the same image in $\mathsf{Bag}\,A\times_{\mathsf{Bag}\,1}\mathsf{Bag}\,B$, namely the
pair $(\{0,1\},\{0,1\})$: the multiset has forgotten \emph{which} $A$-element was paired
with \emph{which} $B$-element.
\prov{MacBeth, computed (\texttt{bag\_not\_container.py}: $\lvert\mathsf{Bag}(2\times
2)\rvert_{2}=10\neq 9$, and $\lvert\cdot\rvert_{3}=20\neq16$). The underlying fact --- the
multiset functor is analytic, not polynomial --- is classical.}
\end{Proposition}

The collision $w_{1}\neq w_{2}$ is the whole story in miniature. A container remembers, for
each position, exactly which value sits there; a multiset remembers only \emph{how many} of
each value, and $w_{1},w_{2}$ differ precisely in a bookkeeping that the multiset discards ---
the pairing of the two coordinates. That lost bookkeeping is a symmetry: a multiset of size
$n$ is an ordered $n$-tuple with its order quotiented away by the symmetric group $S_{n}$.
Quotient the ordered tuples --- which \emph{are} a container, the list functor
$X\mapsto\coprod_{n}X^{n}$ with shape $n$ and positions $\{1,\dots,n\}$ --- by that action
and you leave the polynomial world for the \emph{analytic} one, exactly as the swap did for
ordered pairs. The difference now is only
that $\mathsf{Bag}$ carries this all the way up to a monad: a legitimate algebraic structure
on $\Set$ can sit outside $\Cont$. Chapter~\ref{ch:comonoids} puts precisely this witness to
work (Theorem~\ref{thm:bag-refutes}): being leaf-supported and reverse-total is not enough to
make a monad a \emph{container} monad, and $\mathsf{Bag}$ is what separates the two.

\subsection*{The comonad side: the multigraph comonad}

Now the mirror. Directed multigraphs --- quivers --- are the presheaves on the two-object
category $\cdot\rightrightarrows\cdot$, and they are the coalgebras of a \emph{polynomial}
comonad on $\Set$, the one whose extension is $F_{\mathrm{dir}}(X)=X^{3}+X$ (an edge is a
source, a target, and a name; a vertex is a vertex). This is a container, and under the
equivalence $\Cont$-comonoids $\cong$ small categories that Chapter~\ref{ch:comonoids} will
prove, its category is the walking parallel pair $\cdot\rightrightarrows\cdot$ itself: two
objects, with three arrows out of the source (whence the $X^{3}$) and one out of the target
(whence the $X$).

Undirected multigraphs are what you get by forgetting the orientation of each edge --- by
quotienting the two parallel arrows of $\cdot\rightrightarrows\cdot$ by the swap. Aaron David
Fairbanks observed that this quotient can be taken \emph{at the level of comonads}: the
undirected-multigraph comonad $F$ is the coequaliser, in the category of comonads on $\Set$,
of the identity and the comonad automorphism $\varphi$ induced by that swap
\cite{Fairbanks25}. Its extension is
\[
  F(X)=\frac{X^{3}+X^{2}}{2}+X,
\]
the count of undirected edges (a vertex, a vertex, a name, with the two endpoints now
unordered) plus the vertices. And like $\mathsf{Bag}$, it fails the container test on the
kernel pair of $X\to 1$. Fairbanks's computation is one line and worth reading as the exact
dual of the multiset collision. Here $F(1)=2$ --- an edge slot and a vertex slot --- and
$F(!)$ sends the whole edge summand to the first point and the vertices to the second; so the
pullback of $F(!)\colon F(X)\to F(1)$ along itself, being the pairs of elements that agree in
$F(1)$, splits as the sum of the two summands squared \emph{separately}:
\[
  \Bigl(\tfrac{X^{3}+X^{2}}{2}\Bigr)^{2}+X^{2},\qquad\text{whereas}\qquad
  F(X^{2})=\frac{X^{6}+X^{4}}{2}+X^{2},
\]
and these disagree --- \emph{squaring every summand is not the same as substituting a squared
variable}. That slogan is the polynomial condition itself: a container's extension
substitutes into positions, so $F(X^{2})$ must be ``$F$ of the squared variable'', not ``$F$
squared summand-wise''. The halving --- the orientation thrown away by the quotient --- is
exactly what breaks the substitution.

\begin{Proposition}[A comonad that is not a container]\label{prop:multigraph-comonad}
The undirected-multigraph comonad $F(X)=\tfrac12(X^{3}+X^{2})+X$ on $\Set$ is a genuine
comonad --- undirected multigraphs are comonadic over $\Set$ --- yet it does not preserve
pullbacks (it fails on the kernel pair of $X\to 1$), so its underlying functor is not
polynomial and $F$ is not a container.
\prov{Cited: A.~D.~Fairbanks, MathOverflow answer to ``Examples of non-polynomial comonads on
$\Set$'' \cite{Fairbanks25} (question~457580), verified verbatim by MacBeth. The
pullback-non-preservation is Fairbanks's own displayed computation.}
\end{Proposition}

\begin{teachbox}[The boundary is two-sided]
$\mathsf{Bag}$ and $F$ are a matched pair. Each is a legitimate algebraic structure on
$\Set$ --- one a monad, the other a comonad --- and each is a \emph{quotient of a container
by a symmetry}: ordered tuples modulo $S_{n}$ for $\mathsf{Bag}$, the directed-quiver comonad
modulo edge-reversal for $F$. In both cases the quotient discards \emph{provenance} --- the
pairing of coordinates, the orientation of an edge --- and provenance is exactly what a
polynomial functor is obliged to remember. And in both cases the failure is detected by the
\emph{same} connected limit, the kernel pair of $X\to 1$: $\mathsf{Bag}$ gives $10\neq 9$,
$F$ gives squared-summands $\neq$ squared-variable. ``Polynomial, not merely analytic'' is
therefore not a one-sided caution about which monads are container monads; it is a two-sided
law that constrains, by one test, both the monads that can be container monads and the
comonads that can be small categories. The same boundary returns one categorical level up:
in Chapter~\ref{ch:comonoids} (Remark~\ref{rem:analytic-excluded}) the symmetric functors
$\mathrm{Sym}^{2}$ and $\mathsf{Bag}_{2}$ are excluded as monad \emph{liftings} because they
admit no natural counit --- ``no canonical element to point at'' being the lifting-level
shadow of ``no provenance to remember''.
\end{teachbox}

\section{The boundary of the theory}
\label{sec:boundary}

Step back and the chapter's findings line up into one picture. A container is a
family of sets; $\Cont\cong\Fam(\Set\op)$; and its extension is a coproduct of
representables. Being a coproduct of representables is a property a functor can be
\emph{tested} for --- preservation of connected limits --- and when it holds, the
shapes are $F(1)$ and the positions are named by generic elements, not by fibres.
And the same form that these coproducts take is what fails under quotients: powerset
from the value side, coequalisers from the colimit side, and --- \S\ref{sec:two-witnesses} ---
a genuine monad ($\mathsf{Bag}$) and a genuine comonad ($F$) from the algebraic side,
each a quotient of a container by a symmetry. All of these are one
failure. Containers are what you get by \emph{summing} representables; you leave them
the moment you \emph{quotient}.

That boundary is the polynomial boundary the book keeps returning to. Inside it live
the objects on which the whole equivalence chain runs --- directed containers,
comonoids, small categories, the monoidal structures of the later chapters --- and
each of those structures will silently use that its underlying functor is a
coproduct of representables, that its shapes are recoverable, that connected limits
pass through. Outside it lie the analytic and symmetric functors, the probability
and powerset monads, the quotients: a rich world, but one where the equivalence
chain fractures, and where compositional correctness has to be re-earned rather than
inherited. Knowing which functors are containers is knowing where you are standing.

Chapter~\ref{ch:comonoids} puts the boundary to work: it asks which containers carry
a \emph{comonoid} structure for the composition product, and finds --- the central
surprise of the book --- that the answer is small categories.

\chapter{Algebraic structure on \texorpdfstring{$\Cont$}{Cont}}
\label{ch:algebra}

This chapter is, by design, partly an \emph{outline with cited statements}. The
monoidal and closed structures on containers are most cleanly developed in the
polynomial-functor setting, where $\Cont$ embeds as the full subcategory $\Poly$
of polynomial endofunctors. We record the structures, cite their sources, and
flag carefully which MacBeth has personally verified. With a single exception ---
the closure criterion of §\ref{sec:closed}, which is the book's own contribution
--- we do \emph{not} reprove the cited structures here, and we invent no proofs of
them.

\begin{Remark}[Containers and polynomial functors]
A container $\cont SP$ presents the polynomial functor
$X\mapsto\sum_{s\in S}X^{P(s)}$; the category $\Poly$ of polynomial functors and
natural transformations is equivalent to $\Cont$ (this is
Theorem~\ref{thm:rep} read as ``$\Poly\simeq\Cont$''). We use the two names
interchangeably and adopt Spivak's catalogue of monoidal structures on $\Poly$
\cite{Spivak21}.
\end{Remark}

\section{Products and coproducts}

\begin{Proposition}[Coproducts in $\Cont$]\label{prop:coprod}
$\Cont$ has all coproducts: the coproduct of $\cont{S_i}{P_i}$ has shape set
$\coprod_i S_i$ and positions inherited shapewise,
$P\bigl((i,s)\bigr)=P_i(s)$. On extensions it is the pointwise coproduct of
functors, $\Ext{\coprod_i C_i}=\coprod_i\Ext{C_i}$.
\prov{Cited: AAG05; Spivak 2021}
\end{Proposition}

\begin{Proposition}[Products in $\Cont$]\label{prop:prod}
$\Cont$ has all products: the product of $\cont{S_1}{P_1}$ and $\cont{S_2}{P_2}$
has shape set $S_1\times S_2$ and positions
$P\bigl((s_1,s_2)\bigr)=P_1(s_1)\sqcup P_2(s_2)$, so that
$\Ext{C_1\times C_2}(X)=\Ext{C_1}(X)\times\Ext{C_2}(X)$.
\prov{Cited: Spivak 2021}%
\footnote{Stated as in \cite{Spivak21}; the coproduct (Prop.~\ref{prop:coprod})
MacBeth has checked directly on extensions, the product MacBeth has \emph{not}
independently re-derived in full and reports it as cited.}
\end{Proposition}

\section{The monoidal menagerie}

Spivak \cite{Spivak21} catalogues several monoidal structures on $\Poly$. We list
the ones relevant to this book; the reader should treat each as
\emph{cited, not independently checked} unless tagged otherwise.

\begin{Definition}[The four products, after Spivak 2021]\label{def:menagerie}
On $\Poly\simeq\Cont$ there are, among others:
\begin{itemize}[nosep]
  \item the \emph{coproduct} $+$ (Prop.~\ref{prop:coprod}), unit the initial
        polynomial $0$;
  \item the \emph{product} $\times$ (Prop.~\ref{prop:prod}), unit the terminal
        polynomial $1$;
  \item the \emph{Dirichlet (parallel) tensor} $\otimes$: on shapes
        $S_1\times S_2$, on positions $P_1(s_1)\times P_2(s_2)$, so that
        $\Ext{C_1\otimes C_2}(X)=\sum_{s_1,s_2}X^{P_1(s_1)\times P_2(s_2)}$; unit
        the identity polynomial $X\mapsto X$ (\,$1\triangleleft 1$\,);
  \item the \emph{composition (substitution) product} $\triangleleft$, with
        $\Ext{C_1\triangleleft C_2}=\Ext{C_1}\circ\Ext{C_2}$; unit the identity
        polynomial. This is the \emph{sequencing} product, and comonoids for it
        are the subject of Chapter~\ref{ch:comonoids}.
\end{itemize}
\prov{Cited: Spivak 2021}
\end{Definition}

\begin{Remark}[Honesty about $\otimes$ vs.\ $\times$]
The symbols are easy to confuse. The \emph{Dirichlet} tensor $\otimes$ multiplies
\emph{position sets} ($P_1\times P_2$); the categorical \emph{product} $\times$
takes a \emph{disjoint union} of position sets ($P_1\sqcup P_2$). They have the
same shape set $S_1\times S_2$ but different position data and different
extensions. Both are due to Spivak's catalogue \cite{Spivak21}; MacBeth has
verified the shape/position formulas combinatorially on small cases but defers
the unit and coherence laws to \cite{Spivak21}.
\end{Remark}

The free-monad and cofree-comonad constructions associated with $\triangleleft$
are developed in Chapter~\ref{ch:moncomon}, where they belong to the functors-over-$\Cont$
programme rather than to the structure of $\Cont$ itself.

\section{Closing the structures}
\label{sec:closed}

A monoidal category asks one question above all others: is its tensor
\emph{closed}? Can every map out of a tensor, $p\odot q\to r$, be traded for a
map into a single object, $p\to[q,r]$, that internalises ``the maps from $q$ to
$r$''? For the four structures on $\Cont$ the answers are not uniform --- and the
reason they are not uniform is the most instructive thing in this chapter. This
section keeps one promise: \emph{by its end you will be able to decide, for any of
these tensors, whether it is closed --- and you will see that closure on
containers is polynomiality on sets, read through the extension.}

There is a warning to hear first. If an internal hom exists it is a right Kan
extension, and a right Kan extension of container-valued data has no obligation to
be a container. So the honest question is not ``what is the internal hom'' but
``when is there one at all.'' The rest of the chapter cites; this section is where
it earns one theorem of its own, and the theorem answers exactly that.

Throughout, write $y^A$ for the \emph{representable} container --- a single shape
whose position set is $A$ --- so that $\Ext{y^A}X=X^A$; in particular $y:=y^1$ is
the identity container. Every container is a coproduct of representables,
$p\cong\sum_{i\in S_p}y^{p[i]}$, and we use two facts about $\Cont$ constantly:
co-Yoneda, $\Cont(y^R,q)\cong\Ext qR$ (a map $y^R\to q$ is a shape of $q$ and a
map of its positions into $R$), and the fact that $\Cont(-,q)$ turns the coproduct
of representables into a product, $\Cont\bigl(\sum_i y^{p[i]},q\bigr)\cong\prod_i\Ext q{(p[i])}$.

\subsection{The Dirichlet closure, as a hom of morphisms}

Begin with the Dirichlet tensor $\otimes$, whose closure is the cleanest, and
compute the transpose by hand. Recall that $p\otimes q$ has shape set
$S_p\times S_q$ and, over $(s,t)$, position set $p[s]\times q[t]$. A morphism
$r\otimes p\to q$ is therefore a forward shape map
$U\colon S_r\times S_p\to S_q$ together with, over each $(a,i)$, a backward map
$q[U(a,i)]\to r[a]\times p[i]$. Split each backward map into its two components
($\Hom$ into a product is a product of $\Hom$s), then apply the type-theoretic
axiom of choice twice --- once to pull the $S_q$-choice out over $S_p$, once to
pull the whole package out over $S_r$. What falls out is: for each $a\in S_r$, a
container morphism $f_a\colon p\to q$, together with a single map
$\sum_{i\in S_p}q[f_a\,i]\to r[a]$. But that is precisely a morphism
$r\to[p,q]$, where
\[
   [p,q]\;:=\;\Bigl(\Cont(p,q),\ \ f\mapsto\textstyle\sum_{i\in S_p}q[f_1\,i]\Bigr)
   \;=\;\sum_{f\colon p\to q}y^{\,\sum_{i\in S_p}q[f_1 i]} .
\]

Read this object slowly, because it is the whole point. A \emph{shape} of $[p,q]$
is a morphism $p\to q$. A \emph{position} over that shape $f$ is a $p$-shape $i$
paired with a $q$-position lying over the shape $f_1(i)$ that $f$ sends $i$ to. In
the reading that will drive this book's applications: a \emph{prompt} is a
morphism $p\to q$; a \emph{response} to it is a $p$-shape together with a
$q$-position over the shape the prompt delegated. The shapes of the internal hom
are the morphisms of the category.

\begin{Theorem}[The Dirichlet tensor is closed]\label{thm:dirclosed}
For each container $p$ the functor $(-\otimes p)$ is left adjoint to $[p,-]$:
there is a bijection
\[
   \Cont(r\otimes p,\ q)\;\cong\;\Cont\bigl(r,\ [p,q]\bigr),
\]
natural in $r$ and in $q$. Hence $(\Cont,\otimes,y)$ is a symmetric closed
monoidal category, with internal hom $[p,q]$ as displayed above.
\prov{Cited: Niu--Spivak 2023, Ex.~4.78, eq.~(4.79) \cite{NiuSpivak23};
Spivak 2022, eq.~(44) \cite{SpivakRef}; MacBeth, Lean-verified:
\texttt{DirichletClosed.lean}}%
\footnote{The formula is not new: it is Niu--Spivak's Exercise~4.78 and the
object Spivak calls ``the Dirichlet closure.'' The derivation above reproduces it
character for character. What is MacBeth's is a mechanisation:
\texttt{lean/Containers/Containers/DirichletClosed.lean} constructs the currying
bijection and both round-trip identities, each by \texttt{rfl}
($\eta$ for $\Sigma$ and $\times$ identifies the two composites on the nose), and
the declaration \texttt{dirichlet\_closure} depends on no axioms. To our knowledge
this is the first machine-checked internal hom for the Dirichlet tensor.}
\end{Theorem}

\subsection{The same hom, read as a product of composites}

Here is the first surprise. The object we just assembled by hand --- shapes are
morphisms, positions are delegated positions --- is also, \emph{on the nose}, a
product of composites of the simplest containers:
\[
   [p,q]\;\cong\;\prod_{i\in S_p}\cont{q}{(p[i]\cdot y)},
   \qquad\text{equivalently}\qquad
   \Ext{[p,q]}R\;=\;\prod_{i\in S_p}\Ext q{(R\times p[i])} .
\]
Here $p[i]\cdot y=\sum_{p[i]}y$ is the monomial container of the functor
$X\mapsto p[i]\cdot X=X\times p[i]$, and $\cont{q}{(-)}$ is the composition product
of §\ref{sec:comonoid-table}'s catalogue, whose extension is composition of
functors, so that $\Ext{\cont{q}{(p[i]\cdot y)}}R=\Ext q{(R\times p[i])}$.

That two such differently-built containers coincide is the type-theoretic axiom of
choice wearing a disguise. Unfold the right-hand extension:
$\prod_i\Ext q{(R\times p[i])}=\prod_i\sum_{t}(R\times p[i])^{q[t]}$; distribute
$\prod\sum\cong\sum\prod$ over $i$, and the sum-index that appears is exactly a
forward map $u\colon S_p\to S_q$ together with the backward data $\prod_i
p[i]^{q[u\,i]}$ --- that is, a morphism $p\to q$ --- while the remaining
$R$-exponent collapses to $\sum_i q[u\,i]$. The morphism form and the product form
are the same container, seen from opposite sides of one choice isomorphism.

\begin{Theorem}[The product form of the Dirichlet hom]\label{thm:pihom}
For all containers $p,q$ there is an isomorphism
$[p,q]\cong\prod_{i\in S_p}\cont{q}{(p[i]\cdot y)}$, natural in $q$.
\prov{MacBeth, Lean-verified: \texttt{DirichletHomPi.lean}}%
\footnote{Machine-checked as \texttt{Container.ihomPiIso} in
\texttt{lean/Containers/Containers/DirichletHomPi.lean}: an isomorphism of
containers whose shape component is the choice reindexing
$\prod\sum\cong\sum\prod$. Axiom-free. The anticipated transport along the choice
bijection never materialised --- the two containers are definitionally close enough
that the shape map is a plain rearrangement.}
\end{Theorem}

The product form is worth the detour because it is the one that \emph{generalises}.
Nothing in it is special to $\times$; it is built from a fixed target $q$, the
shapes of $p$, and one binary operation --- here $\times$ --- feeding each position
of $p$ into $q$. Change that operation and you change the tensor. That is the next
subsection.

\subsection{When is a convolutional tensor closed?}

The product $\times$, the Dirichlet tensor $\otimes$, and a third family of
Spivak's are not three unrelated constructions: each arises by the \emph{same
recipe} from a monoidal structure on the base category $\Set$. Given a monoidal
structure $(\star,I)$ on $\Set$, its \emph{Day convolution} is the tensor
$\odot_\star$ on $\Cont$ with
\[
   p\odot_\star q\;=\;\Bigl(S_p\times S_q,\ (s,t)\mapsto p[s]\star q[t]\Bigr),
   \qquad\text{unit }y^I .
\]
Taking $\star=(+,\varnothing)$ gives the categorical product $\times$ (positions
$p[s]\sqcup q[t]$); taking $\star=(\times,1)$ gives the Dirichlet tensor $\otimes$
(positions $p[s]\times q[t]$); taking Garner's $\star=(\vee_S,\varnothing)$ gives
Spivak's third family $\rhd_S$.

\begin{Theorem}[Uniform closure criterion]\label{thm:uniformclosed}
Let $(\star,I)$ be a monoidal structure on $\Set$. The Day tensor $\odot_\star$ is
\emph{left closed} --- i.e.\ $(-)\odot_\star p$ has a right adjoint for every
container $p$ --- \textbf{if and only if} the functor $(-)\star B\colon\Set\to\Set$
is polynomial (preserves connected limits) for every set $B$. When it holds, the
internal hom is
\[
   [p,q]_\star\;=\;\prod_{i\in S_p}\cont{q}{(p[i]\star y)},
   \qquad\text{i.e.}\qquad
   \Ext{[p,q]_\star}R\;=\;\prod_{i\in S_p}\Ext q{(R\star p[i])} .
\]
\prov{MacBeth (both directions); the three named instances are Cited: Spivak 2022,
eqs.~(38)--(40) \cite{SpivakRef}}
\end{Theorem}

\begin{proof}[Why, in two breaths]
\emph{Sufficiency.} If each $(-)\star B$ is polynomial then so is
$R\mapsto\Ext q{(R\star p[i])}$ (a composite of polynomial functors), and so is
their $S_p$-indexed product ($\Poly$ is closed under composition and small
products). Call the resulting container $[p,q]_\star$. The adjunction now reads off
co-Yoneda: $\Cont(r\odot_\star p,q)$, expanded over the representables of $r$ and
$p$, is $\prod_{J}\prod_{i}\Ext q{(r[J]\star p[i])}$, which is exactly
$\Cont\bigl(r,[p,q]_\star\bigr)$.

\emph{Necessity.} This is the half that looks like it should be hard and is a
single line. Suppose $\odot_\star$ is closed, and evaluate the right adjoint at the
one hom $[y^B,y]_\star$. Co-Yoneda and the adjunction give
\[
   \Ext{[y^B,y]_\star}R\;\cong\;\Cont\bigl(y^R,[y^B,y]_\star\bigr)
   \;\cong\;\Cont\bigl(y^R\odot_\star y^B,\ y\bigr)
   \;=\;\Cont\bigl(y^{R\star B},\ y\bigr)\;\cong\;R\star B ,
\]
using $y^R\odot_\star y^B=y^{R\star B}$. So $(-)\star B$ is the extension of a
container --- it is polynomial. Closedness is a statement about \emph{all} homs;
but to extract the obstruction you interrogate only one. The polynomial condition
is not bolted onto closure --- it \emph{is} closure, evaluated at a point.
\end{proof}

\begin{Remark}[What is new here, and what is not]\label{rem:closednovelty}
The three concrete closures are prior art, and this book claims none of them: the
cartesian is Niu--Spivak's Theorem~5.31 \cite{NiuSpivak23} (and, historically,
Altenkirch--Levy--Staton), the Dirichlet is Exercise~4.78 there and Spivak's
eq.~(44) \cite{SpivakRef}, and $\rhd_S$ is Spivak's eq.~(40). Spivak writes the
three internal homs \emph{side by side}; what is MacBeth's is the observation that
they are one formula, together with the biconditional that says exactly which
convolutional tensors are closed and the one-line necessity argument that makes it
a criterion rather than a catalogue. It is a remark on the structure, not a deep
theorem --- but it is the sentence that turns three coincidences into a law.
\end{Remark}

\subsection{Is the condition ever really a condition?}
\label{sec:vacuity}

Theorem~\ref{thm:uniformclosed} has a right-hand side --- ``$(-)\star B$ is
polynomial for every $B$'' --- and an honest reader should ask whether that side
is ever \emph{false}. Every monoidal structure on $\Set$ one can name off the top
of one's head ($+$, $\times$, Garner's $\vee_S$, and their finite iterates) is
polynomial in each variable, so its Day tensor is closed; and it is tempting, on
that evidence, to guess that the condition is automatic --- that \emph{every}
monoidal structure on $\Set$ is polynomial in each variable, and that
Theorem~\ref{thm:uniformclosed} carries a hypothesis it never needs. The guess is
wrong. There is a symmetric monoidal structure on $\Set$ whose Day convolution is
a bona fide tensor on $\Cont$ with no internal hom at all. On containers
\emph{convolutional strictly contains left-closed}, and the side-condition of
Theorem~\ref{thm:uniformclosed} earns its keep.

\begin{Definition}[The collapse tensor]\label{def:collapse}
Define $\star\colon\Set\times\Set\to\Set$, with unit $\varnothing$, by
\[
   A\star B \;=\;
   \begin{cases}
     B & A=\varnothing,\\
     A & B=\varnothing,\\
     1 & A\neq\varnothing \text{ and } B\neq\varnothing,
   \end{cases}
\]
where $1$ is a fixed one-point set. On maps $f\colon A\to A'$, $g\colon B\to B'$,
the arrow $f\star g$ is the surviving factor --- $g$ when $A'=\varnothing$, $f$
when $B'=\varnothing$ --- and the unique map to $1$ when $A',B'$ are both
nonempty. \prov{MacBeth}
\end{Definition}

Two nonempty sets, tensored, \emph{collapse} to a single point; hence the name.
The action on maps is forced and consistent, because a function out of a nonempty
set has nonempty image: $A\neq\varnothing$ forces $A'\neq\varnothing$, so the
``both nonempty'' clause of the target is reachable only from the ``both
nonempty'' clause of the source, where the unique map to $1$ is the only thing one
could write.

\begin{Proposition}[A convolutional tensor with no internal hom]\label{prop:collapse}
$(\star,\varnothing)$ of Definition~\ref{def:collapse} is a symmetric monoidal
structure on $\Set$, yet $(-)\star 2$ is not polynomial. Consequently its Day
convolution $\odot_\star$ is a symmetric monoidal structure on $\Cont$ that is
\emph{not} left closed: $(-)\odot_\star y^2$ has no right adjoint.
\prov{MacBeth}%
\footnote{The monoidality laws and the non-polynomiality were cross-checked
exhaustively on all sets of cardinality $\le 3$
(\texttt{scratch/collapse-tensor/collapse\_hostile2.py}), with a two-element
factor present to witness the non-injective collapse $2\to 1$. This is a
remark-level refinement of Theorem~\ref{thm:uniformclosed}'s landscape, not a Lean
formalisation.}
\end{Proposition}

\begin{proof}[Why, in two breaths]
\emph{Monoidal.} The unitors are identities. The associator
$(A\star B)\star C\to A\star(B\star C)$ and the braiding $A\star B\to B\star A$ are
the evident bijections: if two or more of the factors are nonempty both sides
equal $1$; if exactly one is nonempty both sides are a copy of it; if all are
empty both sides are $\varnothing$. Every arrow in every coherence diagram ---
pentagon, triangle, naturality of $\alpha$ and of $\beta$ --- is therefore one of
just two things. Either it is the \emph{unique} map into $1$, which happens
whenever the node it touches multiplies two or more nonempty factors; such
diagrams commute because $1$ is terminal. Or it is the functorial action on the
\emph{single} surviving factor, which happens when it does not; such diagrams are
the coherence of the trivial unit-only structure, i.e.\ plain functoriality of one
factor. There is no third case, so this is a complete finite proof.

\emph{Not polynomial.} By the strict left unit, $R_2(\varnothing)=\varnothing\star
2=2$, whereas $R_2(1)=1\star 2=1$ because $1,2\neq\varnothing$. But every
polynomial functor $F=\sum_i y^{a_i}$ satisfies
$|F\varnothing|=\#\{i:a_i=\varnothing\}\le\#\{i\}=|F1|$, and here
$|R_2\varnothing|=2>1=|R_2 1|$. So $R_2$ is not polynomial --- concretely, it
sends the mono $\varnothing\hookrightarrow 1$ to the non-mono $2\to 1$. By the
necessity half of Theorem~\ref{thm:uniformclosed}, $\odot_\star$ is not left
closed.
\end{proof}

\begin{teachbox}[Why the collapse tensor slips through]
The moral of the near-misses survives, and it is exactly what makes the collapse
tensor legal. A polynomial functor is one that \emph{remembers}, in its position
sets, which inputs each output used; $\vee_S$ is coherent \emph{because} its normal
form keeps a summand for every relevant subset, and that bookkeeping \emph{is} the
exponent that makes it polynomial. The support tensor
$A\sqcup B\sqcup\{\bullet\}$ tried to discard the provenance and was punished with
no natural associator: its lone separator $\bullet$ could not record \emph{which}
elements it had separated. The collapse tensor discards the provenance too, and
gets away with it, for one reason --- it collapses not to a \emph{marked} point but
to \emph{the} point, a terminal object, which has no provenance to be inconsistent
about. There is nothing for naturality to demand back. And that is, in the same
breath, why it is not polynomial: $1\star B$ is too small to see $B$. Provenance,
polynomiality and coherence are still three names for one thing --- but a tensor is
free to keep no provenance at all, provided it also builds nothing that needs any,
and it forfeits polynomiality by exactly that choice.
\end{teachbox}

The mechanism is worth naming, because it says where to look next. Closure fails
because the unit insertion $\eta_B\colon B\cong\varnothing\star B\to 1\star B$ is
not injective --- multiplying by the unit-image $1$ can \emph{shrink}. (Contrast
the support tensor, whose analogous defect would \emph{enlarge} $1\star B$ with a
phantom element. The two obstructions are independent, which is why a search
calibrated to one sails past the other; the collapse tensor was missed for
precisely this reason.) Injectivity of $\eta_B$ --- \emph{tautness} --- is thus
\emph{necessary} for closure, as is a pullback-preservation condition on the
naturality squares of $\eta$. These conditions locate the failure. But for a
\emph{symmetric} tensor we can do better than locate it: we can \emph{list} the
survivors, and the list is short.

\subsection{The complete list}
\label{sec:classification}

The closure condition of Theorem~\ref{thm:uniformclosed} sorts the monoidal
structures on $\Set$ by a single question: does the Day convolution close? The
collapse tensor showed that both answers occur. We have spent the last pages at the
wall between yes and no; the better question is what lives on the \emph{yes} side.
Which symmetric monoidal structures on $\Set$ have $(-)\star B$ polynomial for every
$B$?

Here is the answer, and it is the cleanest kind: the two families we already met are
the whole pile. Up to isomorphism, a symmetric $\star$ that survives the condition is
either the product $\times$ or one of Garner's $\vee_S$ --- so on containers the
closed symmetric convolutional tensors are exactly $\otimes$ and the $\rhd_S$ family.
(One caveat rides along, an arity bound; it is real, and we meet it honestly at the
end.) Neil's instinct that these closures are a law and not a coincidence is, in the
bounded case, a theorem. Let us prove it.

Start by writing down what polynomiality gives us to work with. Each $R_B=(-)\star B$
has a normal form
\[
   X\star B\;\cong\;\sum_{u\in 1\star B} X^{A_{B,u}}\qquad\text{(naturally in $X$)},
\]
so $R_B$ is pinned by two pieces of data: its \emph{index set} $1\star B=R_B(1)$, and
over each index $u$ its \emph{arity}, the exponent set $A_{B,u}$. Write
$d(B):=\sup_u|A_{B,u}|$ for the largest arity --- the \emph{degree} of $R_B$ --- and
$\kappa:=\sup_B d(B)$ for the degree of the whole family. Our two survivors are
\emph{affine} (degree $\le 1$): the product has $X\star B=B\times X$, index set $B$
and every arity a single point; $\vee_S$ has $X\star B=B+(1+S\times B)\times X$, a
constant part $B$ and a linear part $1+S\times B$. We will show nothing else of
bounded degree survives. Three moves do it.

\medskip\noindent\emph{Move 1: the unit is small.}

\begin{Lemma}[The unit is small]\label{lem:unitsmall}
Under the closure condition, $|I|\le 1$.
\end{Lemma}

\begin{proof}[Why, in two breaths]
$R_1=(-)\star 1$ is polynomial, say $R_1(X)=\sum_k X^{M_k}$. The right unitor makes
$(-)\star I$ the identity, so $1\star I=1$, and by symmetry $R_1(I)=I\star 1=1$. If
$|I|\ge 2$ then each term $I^{M_k}$ has at least one element, so a sum equal to $1$
forces a single term with $M_k=\varnothing$: $R_1\equiv 1$ is constant. Then
$1\star\varnothing=1$ as well, so $R_\varnothing$ has one index and
$X\star\varnothing=X^A$ for a fixed set $A$. Evaluate at $I$: the left unit gives
$\varnothing=I\star\varnothing=I^A$, which is impossible for nonempty $I$. So
$|I|\ge 2$ cannot happen. \end{proof}

The unit is $\varnothing$ or $1$ --- nothing exotic slips in underneath.\footnote{The
finite cardinality search agrees: over $\mathbb N$, the symmetric associative unital
polynomial operations are exactly $a\cdot b$ (unit $1$) and $a+b+s\,ab$ (unit $0$,
$s\ge 0$); units $\ge 2$ admit none
(\texttt{scratch/cardinality-classification.py}).}

\medskip\noindent\emph{Move 2: degrees multiply, and that kills high arity.}

This is the engine. Associativity says $(X\star C)\star B\cong X\star(C\star B)$, that
is $R_B\circ R_C\cong R_{C\star B}$: the family is closed under composition. Compose
the normal forms and read the exponents off,
\[
   R_B\bigl(R_C(X)\bigr)
   \;=\;\sum_{b\in 1\star B}\ \sum_{\varphi:A_{B,b}\to 1\star C}
        X^{\,\sum_{i\in A_{B,b}} A_{C,\varphi(i)}},
\]
so the arities of the composite are the sums $\sum_{i\in A_{B,b}}A_{C,\varphi(i)}$,
and the largest of them has size $d(B)\cdot d(C)$. Hence
\[
   d(C\star B)\;=\;d(C)\cdot d(B):
\]
\emph{degree is a monoid homomorphism} into the cardinals under multiplication. Now
watch the trap close. Suppose some arity had two or more elements, so $\kappa\ge 2$,
and suppose the family had \emph{bounded} degree, $\kappa<\infty$. Choose $B$
realising the maximum, $d(B)=\kappa$. Then $d(B\star B)=\kappa^2>\kappa$ --- a degree
strictly above the supposed maximum. That is a contradiction, and it leaves only one
possibility.

\begin{Lemma}[Bounded degree forces affine]\label{lem:affine}
If $\kappa=\sup_B d(B)$ is finite then $\kappa\le 1$: every $R_B$ is affine,
$X\star B\cong C_B+D_B\times X$ for sets $C_B$ (the empty-arity indices) and $D_B$
(the singleton-arity indices).
\end{Lemma}

The whole force of the classification is that squaring a degree-$\ge 2$ operation
overflows any finite ceiling. It is the same phenomenon as $\vee_S$ staying affine
under its own composition --- $1$ squared is $1$ --- turned into an obstruction for
everything else.

\medskip\noindent\emph{Move 3: an affine, symmetric $\star$ is $\times$ or $\vee_S$.}

Affineness is most of the distance; a single identity covers the rest, and it is the
heart of the matter. Split on the unit (Lemma~\ref{lem:unitsmall}).

If $I=1$: the right unit $X\star 1=X$ gives $C_1=\varnothing$, $D_1=1$, and in
particular $R_1=\id$. Then $1\star\varnothing=\varnothing\star 1=R_1(\varnothing)
=\varnothing$, and since $R_\varnothing$ is affine and takes the value $\varnothing$
at $1$, it is identically $\varnothing$: $X\star\varnothing=\varnothing$ for all $X$.
Hence $C_B=\varnothing\star B=\varnothing$, and $C_B+D_B=1\star B=B$ gives $D_B=B$. So
$X\star B=B\times X$: the tensor is the product.

If $I=\varnothing$: the left unit $\varnothing\star B=B$ gives $C_B=B$, so
$X\star B=B+D_B\times X$. Now bring in symmetry. It says $X\star B\cong B\star X$,
which spelled out is
\[
   B + D_B\times X\;\cong\;X + D_X\times B\qquad\text{naturally in both $X$ and $B$.}
\]
Stare at this identity, because it does all the work. Read both sides as functors of
$B$ with $X$ held fixed. The right-hand side is manifestly \emph{affine in $B$} ---
it is $X+D_X\times B$, and $D_X$ carries no $B$. So the left-hand side is affine in
$B$ too; its only not-obviously-affine part is $D_B\times X$, which forces $D_B$
itself to be an affine functor of $B$:
\[
   D_B\;\cong\;D_\varnothing + S\times B
\]
for some set $S$, the linear coefficient. Put $X=\varnothing$ in the identity:
$B\cong D_\varnothing\times B$ naturally, so $D_\varnothing=1$. Therefore
$D_B=1+S\times B$, and
\[
   X\star B\;=\;B+(1+S\times B)\times X\;=\;X+B+S\times X\times B\;=\;X\vee_S B.
\]

\begin{Lemma}[The symmetry identity]\label{lem:symid}
Let $\star$ be symmetric with unit $\varnothing$ and every $R_B$ affine, so
$X\star B=B+D_B\times X$. Then the symmetry isomorphism forces $D_B\cong 1+S\times B$
for a unique set $S$, and hence $\star\cong\vee_S$. \prov{MacBeth}
\end{Lemma}

Uniqueness of $S$ is free: it is the linear coefficient of the affine functor
$B\mapsto D_B$, and associativity independently pins the associator by making $D$ a
strong monoidal functor $(\Set,\star,\varnothing)\to(\Set,\times,1)$, which
$1+S\times B$ satisfies. The three moves now assemble into the payoff.

\begin{Theorem}[The complete list, bounded arity]\label{thm:classification}
Let $(\Set,\star,I)$ be a symmetric monoidal structure with $(-)\star B$ polynomial
for every set $B$ and with bounded arity ($\kappa<\infty$). Then $\star$ is
isomorphic to exactly one of
\begin{itemize}
\item the product $\times$ (unit $1$), whose Day convolution is the Dirichlet tensor
      $\otimes$; or
\item $\vee_S$ (unit $\varnothing$) for a unique set $S$, whose Day convolution is
      Spivak's $\rhd_S$ --- with $\vee_\varnothing=+$ giving the categorical product
      $\times$ on $\Cont$.
\end{itemize}
Equivalently: the closed symmetric convolutional tensors on $\Cont$ of bounded arity
are precisely $\otimes$ and the $\rhd_S$ family.
\prov{MacBeth (bounded arity)}%
\footnote{The reconstruction was cross-checked on finite sets for $\times$ and
$\vee_S$, $S\le 3$: the affine normal form, associativity
$R_B\circ R_C=R_{C\star B}$, the symmetry identity, and the strong-monoidal
compatibility $D_{B\star C}=D_B\times D_C$, $D_\varnothing=1$ all hold
(\texttt{scratch/verify\_reconstruction.py}).}
\end{Theorem}

This is the sentence the section was built to earn. The three closures we assembled
by hand --- cartesian, Dirichlet, and Spivak's third family --- are not a sampler
from a larger catalogue we have not yet drawn. In the bounded case there \emph{is} no
larger catalogue: they are all of it. The census of Theorem~\ref{thm:uniformclosed}
becomes effective on the closed locus, and the collapse tensor is revealed for what
it is --- not a member we forgot, but the minimal structure ejected for keeping no
provenance at all.

\begin{Remark}[The infinite-arity boundary]\label{rem:infarity}
The hypothesis $\kappa<\infty$ is not a convenience; dropping it is a genuine open
problem, and we state exactly what is missing rather than dress it as a formality.
The engine of Move~2 was the inequality $\kappa^2>\kappa$, and that inequality
\emph{fails} for infinite cardinals, where $\kappa^2=\kappa$. The failure is
structural, not a gap in the bookkeeping. Reframed, ``every $R_B$ is affine'' is
exactly ``every $R_B$ preserves connected \emph{colimits}'', whereas closure delivers
only that each $R_B$ preserves connected \emph{limits} --- and for polynomial
functors those two preservation properties are logically independent, so no
categorical shortcut converts one into the other. One can even write down a formal
fixed point of the arity recursion at an infinite seed: taking $R_2=y+y^\lambda$ with
$\lambda$ infinite, the associativity recursion reproduces
$R_{2\star 2}=y+2^\lambda\cdot y^\lambda$ with no cardinality obstruction and an
associator natural in the visible variable. Counting is therefore provably blind to
the infinite case. If an infinite-arity closed tensor is in fact impossible, the
obstruction must live in the element-level pentagon, jointly in all variables --- a
coherence computation this book does not settle. So the classification is complete
for bounded arity; the unbounded-arity case, and the non-symmetric case (a $\star$
that is left closed without being right closed, which the symmetry identity above
never sees), are left open. \prov{MacBeth; open problem}
\end{Remark}

\subsection{The other three, and a boundary}

The remaining structures round out the picture and each carries a small lesson.

\emph{The categorical product.} $(\Cont,\times,1)$ is cartesian closed, with
exponential $\Ext{q^p}R\cong\prod_{s\in S_p}\Ext q{(R+p[s])}$ --- which is
Theorem~\ref{thm:uniformclosed} at $\star=+$, since $\times$ is the Day convolution
of $(+,\varnothing)$. This is Niu--Spivak's Theorem~5.31 \cite{NiuSpivak23}, first
observed by Altenkirch, Levy and Staton. The cautionary half: $\Cont$ is
\emph{not} locally cartesian closed --- its slices are not all cartesian closed ---
so, despite the exponentials, it does not model extensional dependent type theory.
Closure and pointwiseness pull in opposite directions: $\Cont$ is closed for the
tensor whose extension is \emph{not} pointwise ($\otimes$) and fails to be closed
in the strong local sense for the one whose extension \emph{is} pointwise
($\times$).

\emph{The composition product.} $\triangleleft$ is not symmetric, so ``closed''
bifurcates. It is not closed on the left, but it does have a right coclosure ---
a functor $C_q$ with $\Cont(p,\cont rq)\cong\Cont(C_q\,p,\,r)$, computed by Josh
Meyers \cite[eqs.~(68)--(69)]{SpivakRef}, \cite[Prop.~6.57]{NiuSpivak23}. A reader
collating the two sources should be warned that the same isomorphism is called a
\emph{left} coclosure by Niu--Spivak and a \emph{right} coclosure by Spivak; the
variance convention differs, not the mathematics.

\emph{Outside the census.} The Day recipe manufactures a tensor --- closed or not
--- from every monoidal structure on $\Set$, but not every tensor on $\Cont$ worth
having is a Day convolution. Dorta, Jarvis and Niu's two \emph{uninterpreted}
tensors $\ltimes,\rtimes$ \cite{DJN} are non-convolutional --- their positions are not a
fixed binary operation $p[s]\star q[t]$ applied fibrewise --- and they behave
differently under closure:
$\ltimes$ is \emph{not} closed, while $\rtimes$ is closed on the left in a directed
sense. They mark the boundary of the classification of this section; their full
treatment belongs with the fibrational material and is deferred.
\emph{Attribution.} The tensors themselves are \cite{DJN}~\S6, which gives the
two formulas and then asks: ``We would like to know if there are interpretations
or applications for these monoidal products as there are for the parallel and
composition products.'' Calling them the \emph{Dialectica} tensors ---
$\ltimes$ de~Paiva's $\otimes_{\mathrm{Dial}}$ extended off the homogeneous
slice, $\rtimes$ its directed variant --- is my own identification and answers
that open question; it is not asserted anywhere in \cite{DJN}. (What
\cite{DJN} \emph{does} establish about Dialectica is Prop.~2.13,
$\mathrm{Hmg}(2)\simeq\mathbf{Dial}(\mathbf{Set})$, for the category, not for
$\ltimes,\rtimes$.)

\begin{Remark}[The sentence to carry away]
Every monoidal structure on $\Set$ induces one on containers by Day convolution,
and it is closed exactly when tensoring by a fixed set is a polynomial operation.
The internal hom is then always the same shape: hold $q$, run the shapes of $p$,
feed each position of $p$ through $\star$ into $q$, and take the product. Cartesian
closure, Dirichlet closure, and Spivak's third family are one formula read three
times --- and you learn whether closure holds without ever leaving $\Set$. That
last clause is not a formality: the collapse tensor of
Definition~\ref{def:collapse} is convolutional and \emph{not} closed, so the
polynomial condition is a genuine dividing line and not a fact every Day tensor
enjoys for free.
\end{Remark}

\section{Monoids and comonoids for the four structures}
\label{sec:comonoid-table}

We now have four monoidal structures on $\Cont$: the coproduct $+$, the product
$\times$, the Dirichlet tensor $\otimes$, and the composition product
$\triangleleft$. Each carries a notion of internal \emph{monoid}
$(\mu\colon c\odot c\to c,\ \eta\colon I\to c)$ and internal \emph{comonoid}
$(\delta\colon c\to c\odot c,\ \varepsilon\colon c\to I)$. That is eight
questions. The purpose of this section is to answer all eight, and --- more
usefully --- to say which of the answers are worth remembering.

Throughout, a container is $c=(S,P)$ with $P\colon S\to\Set$; its extension is
$\Ext c(X)=\sum_{s\in S}X^{P(s)}$; a morphism $c\to c'$ is a forward shape map
$\varphi\colon S\to S'$ together with backward position maps
$\varphi^\sh_s\colon P'(\varphi s)\to P(s)$ (the variance of
Definition~\ref{def:contmor}). The four units are the initial container $0=(\emptyset,!)$
for $+$, the terminal container $1=y^0=(1,\ast\mapsto\emptyset)$ for $\times$, and
the identity polynomial $y=y^1=(1,\ast\mapsto 1)$ for both $\otimes$ and
$\triangleleft$.

Here is the whole table; the rest of the section unpacks it.

\begin{center}\small
\renewcommand{\arraystretch}{1.35}
\begin{tabular}{@{}p{2.6cm}p{5.4cm}p{5.4cm}@{}}
\toprule
$\odot$ (unit) & \textbf{comonoids} & \textbf{monoids}\\
\midrule
$+$\quad (\,$0$\,)
  & only $0$ itself \hfill\textsf{\scriptsize[folklore]}
  & every container, uniquely \hfill\textsf{\scriptsize[folklore]}\\
$\times$\quad (\,$1$\,)
  & every container, uniquely \hfill\textsf{\scriptsize[folklore]}
  & monoid on shapes with an \emph{empty} unit-fibre \hfill\textsf{\scriptsize[MacBeth, proved]}\\
$\otimes$\quad (\,$y$\,)
  & families of monoids $\sum_s y^{M_s}$ \hfill\textsf{\scriptsize[MacBeth, proved]}
  & monoid on shapes $+$ oplax fibre-functor \hfill\textsf{\scriptsize[MacBeth, proved]}\\
$\triangleleft$\quad (\,$y$\,)
  & small categories $=$ directed containers \hfill\textsf{\scriptsize[cited]}
  & polynomial monads \hfill\textsf{\scriptsize[cited]}\\
\bottomrule
\end{tabular}
\end{center}

The two \emph{canonical} structures $+$ and $\times$ give nothing: their
(co)monoids collapse. The two \emph{Day-exotic} structures $\otimes$ and
$\triangleleft$ carry all the content, and there it comes in dual pairs. We treat
the collapse in one box and spend the rest of the section on the exotic pair.

\begin{teachbox}[The cartesian/cocartesian collapse]
Because $\times$ is the categorical \emph{product}, every container is a
$\times$-comonoid in exactly one way: the comultiplication is the diagonal
$\Delta\colon c\to c\times c$ and the counit is the unique map $c\to 1$. Dually,
because $+$ is the \emph{coproduct}, every container is a $+$-monoid in exactly one
way, via the codiagonal $\nabla\colon c+c\to c$ and the initial map $0\to c$. The
two opposite corners degenerate the other way: a $+$-\emph{comonoid} needs a
counit $c\to 0$, i.e.\ a map $S\to\emptyset$, which forces $c\cong 0$; so the only
$+$-comonoid is $0$ itself.

Only the $\times$-\emph{monoid} corner has any texture, and even it is thin. A
unit $\eta\colon 1\to c$ has a backward part $P(e)\to 1[\ast]=\emptyset$, which
exists \emph{only if the unit shape has no positions}. So a $\times$-monoid is an
ordinary monoid $(S,\cdot,e)$ on the shape set whose identity $e$ has an empty
fibre, together with a coherent backward routing $P(s\cdot t)\to P(s)\sqcup P(t)$
of each output position into one factor. A container all of whose shapes have
non-empty fibres admits \emph{no} $\times$-monoid at all; when every fibre is
empty one recovers plain monoids on $S$ (counts $1,4,33$ for $|S|=1,2,3$). It is a
genuine, if minor, container-specific notion --- worth naming once and not
dwelling on.
\prov{MacBeth, proved (empty-unit-fibre refinement)}
\end{teachbox}

\subsection{The composition product: categories and monads}

For $\triangleleft$ the answers are the two halves of the polynomial dictionary,
and both are the subject of later chapters, so we only name them here.

\begin{Theorem}[$\triangleleft$-comonoids and $\triangleleft$-monoids]
\label{thm:comp-comonoid-monoid}
A comonoid in $(\Cont,\triangleleft,y)$ is exactly a \emph{directed container},
equivalently a \emph{small category} with object set $S$
(Chapter~\ref{ch:comonoids}). A monoid in $(\Cont,\triangleleft,y)$ is exactly a
\emph{polynomial monad}, equivalently a monad on $\Set$ whose functor is
polynomial.
\prov{Cited: Ahman--Uustalu 2016 \cite{AU16}; Dorta--Jarvis--Niu 2023
Thm 4.3 \cite{DJN}; Gambino--Kock 2013 Thm 4.5 \cite{GK}}
\end{Theorem}

\noindent The comonoid side is developed in full in
Chapter~\ref{ch:comonoids}, including the morphism refinement
$\DCont\cong\Cof$; MacBeth's Lean formalisation of the comonad law
(Theorem~\ref{thm:comonad}) is the object-level half. The monad side is
Gambino--Kock \cite{GK}; the indexed refinement --- monoids in
$(I\text{-}\Cont,\triangleleft,y)$ are monads on $\Set^I$ --- is
De~Pascalis--Uustalu--Veltr\`i \cite{DUV}, and the free-monad instance is the
grafting construction developed in Chapter~\ref{ch:moncomon}. The point to carry
forward is the shape of the dictionary: \emph{$\triangleleft$-comonoid $=$
category, $\triangleleft$-monoid $=$ monad}. It is the pattern we are about to see
again, in a purer form, for $\otimes$.

\subsection{The Dirichlet tensor: a lax/oplax duality}

The Dirichlet tensor multiplies position sets:
$\Ext{c\otimes d}(X)=\sum_{s,t}X^{P(s)\times Q(t)}$, with unit $y$. Its
(co)monoids are the prettiest entry in the table, because the comonoid/monoid
duality here \emph{is} a lax/oplax duality --- visibly, in coordinates. We state
the two classifications and then read the duality off them.

\begin{Theorem}[$\otimes$-comonoids are families of monoids]
\label{thm:dir-comonoid}
A comonoid in $(\Cont,\otimes,y)$ is exactly a container $c=\sum_{s\in S}y^{M_s}$
in which each position set $M_s:=P(s)$ carries an \emph{arbitrary} monoid
structure. The comultiplication is forced to be the diagonal on shapes,
$\delta_1=\Delta\colon S\to S\times S$, and on positions it is the fibrewise
multiplication $M_s\times M_s\to M_s$; coassociativity is associativity of each
$M_s$ and the counit is its unit. On morphisms this is a fibrewise monoid
homomorphism, so
\[
  \mathrm{Comon}(\Cont,\otimes,y)\ \cong\ \Fam(\mathbf{Mon}\op),
\]
the cocommutative comonoids being $\Fam(\mathbf{CMon}\op)$.
\prov{MacBeth, proved; Lean-verified (forward)}%
\footnote{Proved by unwinding the comonoid axioms in container coordinates
(\texttt{proofs/2026-07-17-bare-dirichlet-comonoid.md}); the forward map
$\otimes\text{-comonoid}\to\Fam(\mathbf{Mon}\op)$ is machine-checked in
\texttt{lean/Containers/Containers/DirichletComonoid.lean}
(\texttt{toFamilyOfMonoids}, \texttt{\#print axioms}${}=[\mathtt{Quot.sound}]$).
This answers the $\Poly/\otimes$ slice of Niu--Spivak's open Question~5
of Chapter~9 \cite{NiuSpivak23} by an elementary computation --- not a deep
theorem, but a clean one.}
\end{Theorem}

\begin{Theorem}[$\otimes$-monoids are monoids with an oplax fibre-functor]
\label{thm:dir-monoid}
A monoid in $(\Cont,\otimes,y)$ is exactly a monoid $(S,\cdot,e)$ on the shape set
together with an \emph{oplax monoidal functor}
$P\colon(S,\cdot,e)\to(\Set,\times,1)$ on the fibres --- that is, an oplax
structure map $\phi_{s,t}\colon P(s\cdot t)\to P(s)\times P(t)$ and an oplax unit
$P(e)\to 1$, subject to the oplax coherence (unit and associativity) laws. The
multiplication $\mu\colon c\otimes c\to c$ has forward part the shape
multiplication $\cdot$ and backward part $\phi$; the unit $\eta\colon y\to c$
picks $e$ with the forced terminal map $P(e)\to 1$.
\prov{MacBeth, proved; Lean-verified (forward)}%
\footnote{Proved as the dual/sibling of Theorem~\ref{thm:dir-comonoid}
(\texttt{proofs/2026-07-19-dirichlet-monoid-classification.md}), with two
independent enumerations agreeing on the counts. The forward map is
machine-checked in \texttt{lean/Containers/Containers/DirichletMonoid.lean}
(\texttt{toShapeMonoidOplaxFibres}, \texttt{\#print axioms}${}=\varnothing$).
Niu--Spivak flag the $\otimes$-monoids as future work in Remark~3.78
\cite{NiuSpivak23}; this is an elementary answer to that item, orthogonal to the
$\triangleleft$-monoid classification of De~Pascalis--Uustalu--Veltr\`i
\cite{DUV}, and is not claimed as a deep result.}
\end{Theorem}

Now put the two theorems side by side. A comonoid comultiplies, so its shape part
is a map $S\to S\times S$ --- a comonoid in $(\Set,\times)$, and the only one is
the \emph{diagonal}. Forced. A monoid multiplies, so its shape part is a map
$S\times S\to S$ --- a monoid in $(\Set,\times)$, and that can be \emph{any}
monoid. Free. The single arrow reversal $\delta\colon c\to c\otimes c$ versus
$\mu\colon c\otimes c\to c$ flips the shape map between ``forced diagonal'' and
``arbitrary monoid,'' and simultaneously flips the fibre data between a
\emph{monoid on each fibre} (a lax assignment: $M_s\times M_s\to M_s$) and an
\emph{oplax functor} (structure maps $P(s\cdot t)\to P(s)\times P(t)$ pointing the
other way).

\begin{Remark}[The astonishment]
For the Dirichlet tensor, the comonoid/monoid duality \emph{is} the lax/oplax
duality. A $\otimes$-comonoid is a \emph{lax} gadget --- a monoid living inside
each fibre --- glued to the forced diagonal on shapes; a $\otimes$-monoid is the
exact \emph{oplax} mirror --- structure maps out of each fibre --- glued to a free
choice of monoid on shapes. Reversing the one comultiplication/multiplication
arrow does both flips at once. This is the cleanest instance in the book of a
recurring theme: dualising a structure need not merely swap inputs and outputs; it
can swap the \emph{kind} of structure --- lax for oplax, forced for free.
\end{Remark}

\begin{Remark}[One table, four statements about polynomial functors]
The extension functor $\Ext{-}\colon\Cont\to[\Set,\Set]$ is strong monoidal for
all four structures at once: it sends $+,\times,\otimes,\triangleleft$
respectively to the pointwise coproduct, the pointwise product, the Dirichlet
product, and composition of polynomial endofunctors. Consequently each row of the
table is simultaneously a statement about polynomial endofunctors: the
$\triangleleft$ row is the comonad/monad dictionary, the $\otimes$ row classifies
Dirichlet-product (co)monoids, and the $+,\times$ rows the pointwise ones. The
directed containers of the next chapter are the $\triangleleft$-comonoid cell,
read as comonads on $\Set$.
\prov{Cited: Niu--Spivak 2023 (extension is strong monoidal)}
\end{Remark}

\section{An aside: the lens subcategory}
\label{sec:lens-subcat}

There is one full subcategory of $\Cont$ worth pointing out before we leave the
algebra behind, because it is exactly where compositional systems people already
live. Restrict to the containers whose position set is \emph{constant} in the
shape --- the \emph{monomials} $S\cdot y^A=(S,\ s\mapsto A)$
\cite[Def.\ 2.15]{NiuSpivak23}.

\begin{teachbox}[Monomials are bimorphic lenses]
The monomials span a full subcategory of $\Cont$ equivalent to the category of
\emph{bimorphic lenses}. Unwinding the container morphism $S\cdot y^A\to T\cdot
y^B$: the forward shape map is a \emph{get} $f\colon S\to T$, and the backward
position map, constant in the shape, is a \emph{put} $f^\sh\colon S\times B\to A$.
These are exactly the two components of a bimorphic lens between $(S,A)$ and
$(T,B)$, and the container composition law reproduces lens composition.
\prov{Cited: Niu--Spivak 2023, Example 3.41 (monomials $\simeq$ bimorphic lenses,
after Hedges); get/put as in eqn.\ (3.42)}%
\footnote{Convention: we present the raw morphism data $f\colon S\to T$,
$f^\sh\colon S\times B\to A$. Niu--Spivak's eqn.\ (3.42) relabels the get as
$J\to I$ under the functional-programming reading ``state $=$ domain''; the
underlying data is the same.}
\end{teachbox}

\noindent Which of the categorical structures of this chapter does this
subcategory inherit? The answer is a clean dividing line --- and it is worth
stating carefully, because the intuition ``lenses are a poorer setting, they lack
most of the structure'' overshoots.

\begin{Proposition}[What the lens subcategory keeps and loses]
\label{prop:lens-closure}
The full subcategory of monomials in $\Cont$ is closed under
\begin{itemize}[nosep]
  \item \emph{binary products}: $S\cdot y^A\times T\cdot y^B=(S\times T)\cdot
        y^{A+B}$, again a monomial;
  \item the \emph{terminal object} $1=1\cdot y^0$ and the \emph{initial object}
        $0=\emptyset\cdot y^A$;
\end{itemize}
and is \emph{not} closed under
\begin{itemize}[nosep]
  \item \emph{binary coproducts}: $S\cdot y^A+T\cdot y^B$ is a monomial only when
        $|A|=|B|$ (the two summands carry different, unequated fibres);
  \item \emph{initial algebras} ($W$-types, and the free monad of
        Chapter~\ref{ch:moncomon}): the leaf-position set varies with the depth of
        the tree, so the fixpoint is not a monomial.
\end{itemize}
\prov{MacBeth, verified (hand $+$ exhaustive check on small monomials)}%
\footnote{The product formula is also Niu--Spivak's (p.\ 134); the coproduct and
initial-algebra failures are checked by hand and confirmed by an exhaustive
$81/81$ enumeration in
\texttt{scratch/2026-07-19-lens-monomial-closure.md}.}
\end{Proposition}

So the line does not fall where one might guess. Products, the terminal object and
the initial \emph{object} all stay inside the lens world; what escapes it is
precisely the \emph{colimit\,/\,inductive} side --- binary coproducts and initial
algebras. The lens subcategory is poorer in colimits and recursion, not in limits.
That is the clean teaching point: bidirectional transformations compose and take
products freely, but the moment you \emph{branch} (coproduct) or \emph{recurse}
(initial algebra) you leave the monomials and need the full generality of
$\Cont$.

\section{Two predicate liftings: \texorpdfstring{$\mathsf{All}$}{All} and \texorpdfstring{$\mathsf{Exists}$}{Exists}}
\label{sec:predicate-liftings}

Here is a construction Neil Ghani arrived at independently, and it turns out to
name two things this book has kept apart. Give a container $X=\cont{S_X}{P_X}$ and
a \emph{predicate} on its positions --- for now, just a rule that attaches a set to
each position, depending on the value sitting there. There are two obvious ways to combine those sets over the
positions of a shape: take them \emph{all} at once, or ask for \emph{some} one of
them. ``All'' is a product $\prod$; ``some'' is a coproduct $\coprod$. Two
liftings, born of a single container, differing only in $\prod$ versus $\coprod$.

The promise of this section is that these two liftings are not new objects at all.
One of them is the composition product $\triangleleft$ you already met in the
menagerie (\S\ref{sec:closed}) --- so the ``exists'' lifting \emph{is} sequencing.
The other is a construction you cannot always perform, and \emph{why} you cannot is
the boundary the whole rest of the book circles: it exists on all of $\Cont$ only
when the map it rides is position-preserving. The punchline is a clean law relating
the two --- a Fubini identity that makes ``for all, and then for all'' into ``for
all pairs.''

\subsection{The two liftings, defined}

Fix a container $X=\cont{S_X}{P_X}$. A \emph{predicate} we take to be a second
container $Y=\cont{S_Y}{P_Y}$: its shapes $S_Y$ are the \emph{values} the predicate
can report, and its positions $P_Y(t)$ are the \emph{witnesses} carried by the
value $t$. (This is exactly the container fibration $\Cont\to\Set$,
$\cont SP\mapsto S$: an object over a shape set is a position family, i.e.\ a
container. We return to it below.) Evaluate the predicate at a filled shape of $X$:
an element of $\Ext X{S_Y}$ is a pair $(s,g)$ with $s\in S_X$ and
$g\colon P_X(s)\to S_Y$ --- a shape of $X$ with a value $g(p)$ hung at each
position $p$.

\begin{Definition}[The two liftings of a container]\label{def:all-exists}
For $X$ and $Y$ as above and $(s,g)\in\Ext X{S_Y}$,
\[
  (\mathsf{All}\ X\ Y)(s,g)\;=\;\prod_{p\in P_X(s)}P_Y\bigl(g(p)\bigr),
  \qquad
  (\mathsf{Exists}\ X\ Y)(s,g)\;=\;\coprod_{p\in P_X(s)}P_Y\bigl(g(p)\bigr).
\]
Each assembles into a bifunctor on $\Cont$ with the \emph{same} shapes
$\Ext X{S_Y}=\coprod_s S_Y^{P_X(s)}$ and differing only on positions:
\[
  \mathsf A\,X\,Y=\bigl(\Ext X{S_Y},\ \mathsf{All}\ X\ Y\bigr),
  \qquad
  \mathsf E\,X\,Y=\bigl(\Ext X{S_Y},\ \mathsf{Exists}\ X\ Y\bigr).
\]
\prov{The liftings and the two operations $\mathsf A,\mathsf E$ are Neil Ghani's
(personal communication, 2026); the identifications and laws below are MacBeth's.}
\end{Definition}

\begin{Example}[A node with two children]\label{ex:all-exists}
Let $X$ have a single shape with two positions, $P_X=\{1,2\}$ (a binary node), and
let $Y=\mathsf{Maybe}=\cont{\{\mathrm n,\mathrm j\}}{P_Y}$ with $P_Y(\mathrm n)=\varnothing$,
$P_Y(\mathrm j)=\{*\}$ (Example~\ref{ex:menagerie}). A filled shape is a choice
$g\colon\{1,2\}\to\{\mathrm n,\mathrm j\}$ --- one of the four words
$\mathrm{nn},\mathrm{nj},\mathrm{jn},\mathrm{jj}$. Then
$\mathsf{All}$ takes the \emph{product} of the two witness sets and $\mathsf{Exists}$
their \emph{sum}:
\[
\begin{array}{c|c|c}
 g & \mathsf{All}=P_Y(g1)\times P_Y(g2) & \mathsf{Exists}=P_Y(g1)+P_Y(g2)\\\hline
 \mathrm{nn} & \varnothing\times\varnothing=\varnothing & \varnothing+\varnothing=\varnothing\\
 \mathrm{nj} & \varnothing\times 1=\varnothing & \varnothing+1=1\\
 \mathrm{jn} & 1\times\varnothing=\varnothing & 1+\varnothing=1\\
 \mathrm{jj} & 1\times 1=1 & 1+1=2
\end{array}
\]
$\mathsf{All}$ is nonempty only when \emph{every} position reports a witness
($\mathrm{jj}$); $\mathsf{Exists}$ is nonempty as soon as \emph{some} position does,
and at $\mathrm{jj}$, where both positions report, it records \emph{which} one ---
the two ways to pick that $\mathsf{All}$ collapses to a single point. Hold on
to that last cell: it is the whole difference between the two liftings.
\end{Example}

\subsection{The first surprise: \texorpdfstring{$\mathsf E$}{E} is the composition product}

Unfold $\mathsf E\,X\,Y$ against the definition of $\triangleleft$. A shape of
$X\triangleleft Y$ is an $X$-shape with a $Y$-shape planted at each $X$-position,
i.e.\ a pair $(s,g)$ with $g\colon P_X(s)\to S_Y$ --- the shapes of $\mathsf E\,X\,Y$
on the nose. A position of $X\triangleleft Y$ over $(s,g)$ is an $X$-position $p$
\emph{together with} a $Y$-position beneath the value there, i.e.\ an element of
$\coprod_{p\in P_X(s)}P_Y(g(p))$ --- the positions of $\mathsf{Exists}$ on the nose.

\begin{Proposition}[$\mathsf E=\triangleleft$]\label{prop:E-is-comp}
$\mathsf E\,X\,Y=X\triangleleft Y$, and thus $(\Cont,\mathsf E,1)=(\Cont,\triangleleft,y)$
--- the monoidal spine of this book (Definition~\ref{def:menagerie}). $\mathsf E$ is a
bifunctor on \emph{all} of $\Cont$.
\prov{MacBeth (identification); the $\triangleleft$-monoidal structure is
Lean-verified, \texttt{lean/Containers/Containers/Monoidal.lean} (pentagon and
triangle by \texttt{rfl}).}
\end{Proposition}

The reader hook Neil offers is worth repeating, because it makes $\coprod$
inevitable. Read $X$ as a conversation with several prompts (its positions) and $Y$
as a follow-up conversation. A \emph{reply} in the composite $X\triangleleft Y$ is a
reply to one prompt of $X$ \emph{plus} a reply to the follow-up it opened --- you
choose \emph{which} prompt you are answering, then answer downstream of it. Choosing
which, then answering, is exactly $\coprod_p P_Y(g(p))$. Sequencing is
existential.

\subsection{The second surprise: \texorpdfstring{$\mathsf A$}{A} lives only on
cartesian morphisms}

Now try the same for $\mathsf A$, and watch it fail in an instructive way. Both
liftings share their shape map: a container morphism $\alpha=(u,f)\colon X\to X'$
sends $(s,g)\mapsto(u\,s,\ g\circ f_s)$, using the backward position map
$f_s\colon P_{X'}(u\,s)\to P_X(s)$ to re-address the values. That is defined for
\emph{every} $\alpha$. The liftings part on \emph{positions}, and this is where the
variance bites.

\begin{itemize}[nosep]
  \item For $\mathsf E$ the induced position map is on \emph{sums},
        $\coprod_{p'}P_Y(g\,f p')\to\coprod_p P_Y(g\,p)$, sending $(p',b)\mapsto(f p',b)$.
        A coproduct pushes forward \emph{along} $f$ --- always available.
  \item For $\mathsf A$ it would be on \emph{products},
        $\prod_{p'}P_Y(g\,f p')\to\prod_p P_Y(g\,p)$. A product is
        \emph{contravariant} in its index set: to produce a coordinate at each
        $p\in P_X(s)$ from coordinates indexed by $p'\in P_{X'}(u\,s)$ you must
        re-address $p$ \emph{backward} through $f$ --- you need a section of $f$.
\end{itemize}

How many such sections are there? Exactly $\prod_{p\in P_X(s)}f_s^{-1}(p)$, the
choices of a preimage for each position. This one count settles everything.

\begin{Theorem}[$\mathsf A$ is a cartesian-only bifunctor]\label{thm:A-cartesian}
The natural pushforwards $\mathsf A\,X\,Y\to\mathsf A\,X'\,Y$ over
$\alpha=(u,f)$ are indexed by $\prod_{s}\prod_{p\in P_X(s)}f_s^{-1}(p)$. Hence:
\begin{itemize}[nosep]
  \item if some $f_s$ is \emph{not surjective}, that product is \emph{empty} ---
        \emph{no} natural map exists;
  \item if every $f_s$ is surjective but some is not injective, maps exist but
        \emph{none is canonical};
  \item if every $f_s$ is a \emph{bijection} ($\alpha$ \emph{cartesian}), the
        product is a \emph{singleton}: the pushforward is forced, $\theta_s=f_s^{-1}$.
\end{itemize}
So $\mathsf A$ extends canonically to a bifunctor $\Cont_{\mathrm{cart}}\times\Cont\to\Cont$,
and to no non-cartesian morphism. In its \emph{second} argument $\mathsf A$ is
functorial for all morphisms, which act inside each factor pointwise.
\prov{MacBeth, proved. The section-set count is a one-line Yoneda calculation
(a natural transformation between the representable functors $\Set^{P}(D,-)$,
with $D(p)=f^{-1}(p)$); details and a $3000/3000$ random check in
\texttt{proofs/2026-08-08-A-E-predicate-liftings.md}, \S2.}
\end{Theorem}

This is the object-level shape of a flag Neil raised: \emph{``I cannot see how to
define $\mathsf A$ on polynomial functors, and it is functorial in the first
argument only on cartesian morphisms.''} The theorem says exactly why, and that it
is not a defect of imagination: along a morphism that drops a position, there is no
natural pushforward to be had, because $\prod$ has no coordinate to carry at the
dropped position. Contrast $\coprod=\mathsf E$, which simply forgets that position's
contribution and moves the rest forward.

\begin{Remark}[This is $T_M$, one level down]\label{rem:A-is-TM}
$\mathsf{All}=\prod$ is the proof-relevant $\prod$-lifting $T_M$ of Ahman--Bauer
\cite{AhmanBauer24} (``proof-relevant'' is the type theorist's name for
\emph{records witnesses, not truth}; we use the fibrational wording and log the
synonym here, once). Take the first argument to be a monad's multiplication
$\alpha=\mu^M\colon M\triangleleft M\to M$. Theorem~\ref{thm:A-cartesian} then reads:
$\mu^M$ lifts to a \emph{canonical} $T_M$-multiplication \emph{iff} $\mu^M$ is
cartesian. So the tempting ``core theorem'' --- \emph{every $\Set$ monad lifts to a
container monad via $T_M$} --- is \textbf{false}. It fails outright for the reader
and state monads, whose $\mu^M$ \emph{drops} leaves: this is exactly
Theorem~\ref{thm:A-cartesian}'s non-surjective case, where no natural pushforward
exists at all. It \emph{survives} for a wide class --- the ``polynomial'' monads ---
that includes $\mathsf{List}$ (whose $\mu$ is cartesian) and, more delicately,
finite powerset $\mathsf{Pf}$, whose $\mu=\bigcup$ \emph{merges} leaves rather than
dropping them and so still supports a multiplication, though not the
cartesian-canonical one. So cartesianness is the border for the \emph{canonical}
lift, while merging buys a non-canonical one and dropping buys nothing: the precise
ladder --- with two fibrations, a $\mathbb Z/2$ grading, and the $\coprod$-lifting
that Reader \emph{does} carry --- is \S\ref{sec:moncomon-fibration} in
Chapter~\ref{ch:moncomon}. Neil's functoriality flag and our ``multiplication drops
leaves'' are one statement.
\prov{MacBeth, proved (\texttt{proofs/2026-08-07-proof-relevance-boundary.md});
$\prod$-lift is \cite{AhmanBauer24}. Two-fibration framing:
Hermida--Jacobs \cite{HermidaJacobs98}.}
\end{Remark}

\subsection{The law: \texorpdfstring{$\mathsf A$}{A} is a module of \texorpdfstring{$\triangleleft$}{comp}}

Cartesian-only or not, $\mathsf A$ obeys a clean composition law at the level of
objects, and it is the reason Neil wanted these two operations named side by side.
The law is a Fubini identity for the dependent product.

\begin{Theorem}[The action law]\label{thm:action-law}
For all containers $X,Y,C$,
\[
  \mathsf A\,X\,(\mathsf A\,Y\,C)\;\cong\;\mathsf A\,(X\triangleleft Y)\,C,
\]
naturally in $C$, with unit $\mathsf A\,y\,C=C$. Thus $\mathsf A$ is a
\emph{left module} of the $\triangleleft$-monoidal category $(\Cont,\triangleleft,y)=(\Cont,\mathsf E,1)$
acting on $\Cont$: the object-level associativity and unit equations hold
unconditionally; as a \emph{functorial} action it lives on $\Cont_{\mathrm{cart}}$
(Theorem~\ref{thm:A-cartesian}).
\prov{MacBeth, proved (\texttt{proofs/2026-08-08-A-E-predicate-liftings.md}, \S3;
shape/position multiset check $2000/2000$).}
\end{Theorem}

\begin{proof}[Why it holds]
On positions the identity is pure Fubini for $\prod$. A position of the left-hand
side over a filled shape is an element of
$\prod_{p\in P_X(s)}\bigl(\prod_{q\in P_Y(t_p)}P_C(\cdots)\bigr)$ --- for every
$X$-position, a choice over every $Y$-position beneath it --- and
$\prod_p\prod_q=\prod_{(p,q)}$ collapses this to a single choice over the composite
positions $\coprod_p P_Y(t_p)$ of $X\triangleleft Y$, which is precisely a position
of the right-hand side. On shapes the matching iso is the distributivity of $\prod$
over $\coprod$ (type-theoretic choice) followed by currying. Nothing about a base
monad enters: the law is a fact about $\prod$ and $\triangleleft$ alone.
\end{proof}

The intuition Neil gives is debate preparation. To hold
$\mathsf A\,X\,(\mathsf A\,Y\,C)$ is to prepare, for \emph{every} line of
questioning $p$ your opponent might open, and then under it for \emph{every}
follow-up $q$, a full answer. Fubini says that is the same as being ready for every
\emph{pair} $(p,q)$ at once --- a single ``for all'' over the composite conversation
$X\triangleleft Y$. Universal preparation composes.

\begin{Remark}[The contrast that names the two structures]
The parallel law for $\mathsf E$,
$\mathsf E\,X\,(\mathsf E\,Y\,C)=\mathsf E\,(X\triangleleft Y)\,C$, is nothing but
associativity of $\triangleleft$ (Proposition~\ref{prop:E-is-comp}): $\mathsf E$
acts by \emph{being} the monoidal product. $\mathsf A$ acts \emph{alongside} it, as
a module. The mixed law does \emph{not} hold: $\mathsf A\,X\,(\mathsf E\,Y\,C)$ and
$\mathsf E\,(\mathsf A\,X\,Y)\,C$ carry $\prod_p$ against $\coprod_p$ in the same
slot, and these agree iff every shape of $X$ carries \emph{exactly one} position
($X$ linear) --- strictly finer than non-branching, since an empty-position shape
already breaks it ($\coprod_\varnothing=0\neq1=\prod_\varnothing$). $\prod$ and
$\coprod$ genuinely refuse to commute, and that refusal is the same branching gap
that obstructs the arrow-category composition of Chapter~\ref{ch:moncomon}.
\prov{MacBeth, proved (\texttt{proofs/2026-08-08-A-E-predicate-liftings.md}, \S4).}
\end{Remark}

\subsection{Where this points}

Two liftings, one container. $\mathsf E=\coprod$ is sequencing, defined everywhere,
and its self-composition is just associativity of $\triangleleft$. $\mathsf A=\prod$
is universal quantification, defined only over cartesian morphisms, and its
composition with itself is the module law above. The asymmetry --- $\coprod$ pushes
forward, $\prod$ needs a section --- is the same variance that decides, in
Chapter~\ref{ch:moncomon}, which $\Set$-monads lift to $\Cont$ and which only lift
their existential shadow. We have met the boundary here in miniature, on a single
morphism; there we meet it as a fibration and read the whole ladder off it.

\begin{Remark}[Honesty]
The action law is stated as a canonical isomorphism of containers. Whether it is a
strict definitional equality in a given encoding --- as it is for the
$\triangleleft$-associator, whose coherence is \texttt{rfl} in Lean --- is a
formalisation question flagged but not settled here. We have checked the two module
\emph{equations} (associativity and unit) at the object level; the module
\emph{coherence} $2$-cells (the pentagon and triangle relating the associator to
$\triangleleft$'s own) are automatic for a strict encoding and remain to be
verified for a bicategorical one.
\end{Remark}

\chapter{Monoids and comonoids in \texorpdfstring{$\Cont$}{Cont}: directed containers, categories, and the free monoid}
\label{ch:comonoids}\label{ch:phase2dummy}

This is the heart of the book. We study the two internal algebraic structures for the
composition product $\triangleleft$. On the \emph{comonoid} side --- the bulk of the chapter
--- we show that comonoids for $\triangleleft$ are \emph{directed containers}, that directed
containers are \emph{small categories}, and --- the refinement that the rest of the grant turns
on --- that the right \emph{morphisms} are \emph{cofunctors}, not functors. On the
\emph{monoid} side, we recall that a $\triangleleft$-monoid is a polynomial monad
(Theorem~\ref{thm:comp-comonoid-monoid}) and, at the close of the chapter, \emph{state} the
one construction we will need in Chapter~\ref{ch:moncomon}: the \emph{free}
$\triangleleft$-monoid on a container. Its universal property is deferred --- and proved --- in
the next chapter; here we only pin down what it is.

\section{Directed containers}

\begin{Definition}[Directed container]\label{def:dcont}
A \emph{directed container} is a container $\cont SP$ together with
\begin{itemize}[nosep]
  \item a \emph{root} $\mathsf o(s)\in P(s)$ for each $s$;
  \item a \emph{sub-shape} operation $s\downarrow p\in S$ for $p\in P(s)$;
  \item a \emph{shift} $p\oplus q\in P(s)$ for $p\in P(s)$, $q\in P(s\downarrow p)$;
\end{itemize}
subject to the five laws (for all $s$, $p\in P(s)$, $q\in P(s\downarrow p)$,
$q'\in P((s\downarrow p)\downarrow q)$):
\begin{align*}
  &\text{(D1)} & s\downarrow\mathsf o(s)&=s;
  & &\text{(D2)} & \mathsf o(s)\oplus q&=q;
  & &\text{(D3)} & p\oplus\mathsf o(s\downarrow p)&=p;\\
  &\text{(D4)} & s\downarrow(p\oplus q)&=(s\downarrow p)\downarrow q;
  & &\text{(D5)} & (p\oplus q)\oplus q'&=p\oplus(q\oplus q').
\end{align*}
\end{Definition}

\noindent D2, D3, D5 make each $P(s)$ a monoid-like structure with unit
$\mathsf o(s)$; D1, D4 track how sub-shapes nest. (The well-typedness of D5 needs
D4 --- exactly the dependency that recurs as a dependent-cast obstacle in the Lean
formalisation.)

\section{Directed containers are comonads}

\begin{Definition}[The comonad of a directed container]\label{def:dcomonad}
On $D:=\Ext{\cont SP}$ define
\[
  \varepsilon(s,v)=v(\mathsf o(s)),\qquad
  \delta(s,v)=\bigl(s,\ \lambda p.\,(s\downarrow p,\ \lambda q.\, v(p\oplus q))\bigr).
\]
\end{Definition}

\begin{Theorem}[Comonad laws $\Longleftrightarrow$ D1--D5]\label{thm:comonad}
$(D,\varepsilon,\delta)$ is a comonad, and each comonad law is exactly one
directed-container law: left counit is D1 (shape) and D2 (value); right counit is
D3; coassociativity is D4 (shapes) and D5 (deepest values). Conversely a comonad
of this form forces D1--D5.
\prov{Cited: Ahman--Chapman--Uustalu 2014; MacBeth, Lean-verified:
Directed.lean}
\end{Theorem}

\begin{proof}[Proof sketch]
Direct calculation. For left counit, $\varepsilon(\delta(s,v))$ reduces to
$(s\downarrow\mathsf o(s),\,\lambda q.\,v(\mathsf o(s)\oplus q))$, which is $(s,v)$
by D1 then D2. Right counit is D3; coassociativity unfolds both composites to a
common triple nest, equal by D4 (inner shapes) and D5 (deepest values). The full
computation, and the converse, are checked with zero \texttt{sorry} in the
committed file \texttt{lean/Containers/Containers/Directed.lean}
(\texttt{left\_counit}, \texttt{right\_counit}, \texttt{coassoc}); the
``container is a comonad'' direction is Ahman--Chapman--Uustalu \cite{ACU14}.
\end{proof}

\section{Directed containers are small categories}

\begin{Theorem}[Object dictionary]\label{thm:objdict}
Directed-container structures on $\cont SP$ are in bijection with small categories
$\mathcal C$ with object set $S$ and, for each $s$,
$P(s)=\coprod_{s'}\mathcal C(s,s')$ the morphisms \emph{out of} $s$, under
\[
  \ob\mathcal C=S,\quad
  \mathcal C(s,s')=\{p\in P(s):s\downarrow p=s'\},\quad
  \id_s=\mathsf o(s),\quad
  q\circ p=p\oplus q .
\]
In particular $s\downarrow p=\cod(p)$, and every small category arises from a
unique directed container.
\prov{Cited: Ahman--Uustalu 2016}
\end{Theorem}

Combining Theorems~\ref{thm:comonad} and \ref{thm:objdict} with the polynomial
incarnation gives the object-level chain promised in the Preface: a directed
container is a comonoid in $(\Poly,\triangleleft)$, hence
\[
  \textbf{Containers}\ \supseteq\ \textbf{Directed Containers}
  \ \cong\ \textbf{Polynomial Comonoids}\ \cong\ \textbf{Small Categories}.
\]

\section{The morphism refinement: \texorpdfstring{$\DCont\cong\Cof$}{DCont = Cof}}

A dictionary on objects invites a dictionary on morphisms. The naive guess --- a
morphism of directed containers is a \emph{functor} --- is \emph{false on
variance}. A container morphism has a backward position map
$f^\sharp_s\colon P'(fs)\to P(s)$, whereas a functor pushes morphisms forward. The
maps native to directed containers are \emph{cofunctors}.

\begin{Definition}[Cofunctor / retrofunctor]\label{def:cofunctor}
A \emph{cofunctor} $\varphi\colon\mathcal C\co\mathcal D$ is an object map
$\varphi_0\colon\ob\mathcal C\to\ob\mathcal D$ together with a \emph{lift}: for each
object $c$ and each $h\colon\varphi_0 c\to d$ in $\mathcal D$, a morphism
$\varphi^\uparrow(c,h)\colon c\to c'$ in $\mathcal C$, subject to
\begin{align*}
  \text{(C0)}\ & \varphi_0(\cod\varphi^\uparrow(c,h))=\cod(h);\quad
  \text{(C1)}\ \varphi^\uparrow(c,\id_{\varphi_0 c})=\id_c;\\
  \text{(C2)}\ & \varphi^\uparrow(c,h'\circ h)=\varphi^\uparrow(c',h')\circ
        \varphi^\uparrow(c,h),\quad c'=\cod\varphi^\uparrow(c,h).
\end{align*}
Small categories and cofunctors form a category $\Cof$.
\prov{Cited: Aguiar 1997; Clarke 2020; Spivak 2021 (Def.\ 3.21)}
\end{Definition}

\begin{Remark}[$\Cof=\Cathash$]
$\Cof$ is exactly $\Cathash=\mathrm{Comod}(\Poly)$, the category of comonoids in
$(\Poly,\triangleleft)$ and their morphisms: a cofunctor \emph{is} a comonoid
homomorphism \cite{Spivak21,SS2405}. This is why the object dictionary
(Theorem~\ref{thm:objdict}) cannot extend to $\Cat$: functors are carried not by
comonoid morphisms but by bicomodules \cite{SS23}.
\end{Remark}

\begin{Definition}[Directed-container morphism]\label{def:dcm}
A \emph{morphism of directed containers} $D\to D'$ is a container morphism
$(f,f^\sharp)$ with $f\colon S\to S'$ and $f^\sharp_s\colon P'(fs)\to P(s)$,
satisfying, for all $s$, $h\in P'(fs)$, $k\in P'((fs)\downarrow' h)$,
\begin{align*}
  \text{(M0)}\ & f(s\downarrow f^\sharp_s(h))=(fs)\downarrow' h;\quad
  \text{(M1)}\ f^\sharp_s(\mathsf o'(fs))=\mathsf o(s);\\
  \text{(M2)}\ & f^\sharp_s(h\oplus' k)=
        f^\sharp_s(h)\ \oplus\ f^\sharp_{\,s\downarrow f^\sharp_s(h)}(k).
\end{align*}
\end{Definition}

\begin{Theorem}[The morphism dictionary]\label{thm:dcontcof}
The assignment $(f,f^\sharp)\mapsto(\varphi_0:=f,\ \varphi^\uparrow(s,h):=f^\sharp_s(h))$
is a bijection on every hom-set carrying (M0)--(M2) to (C0)--(C2) term for term,
and it respects identities and composition. It is moreover a bijection on objects
(Theorem~\ref{thm:objdict}). Hence there is an isomorphism of categories
\[
  \DCont\ \cong\ \Cof .
\]
\prov{Cited: Ahman--Uustalu 2016 \S3.2--3; Shapiro--Spivak 2023, Cor.\,5.12}%
\footnote{MacBeth has formalised this isomorphism in Lean (file
\texttt{Cofunctor.lean}, built and reported zero-\texttt{sorry} in
\texttt{papers/dcont-cof.tex}), but that source is \emph{not yet committed} to
this repository's \texttt{lean/} tree; so the statement is tagged \textsf{[Cited]}
here, not \textsf{Lean-verified}, pending the source landing in the repo. The
\emph{identification} of Ahman--Uustalu's ``relative split pre-opcleavages'' with
Spivak's cofunctors, and the bridge to delta lenses below, are MacBeth's
connective contribution; see \texttt{papers/dcont-cof.tex}.}
\end{Theorem}

\begin{proof}[Proof sketch]
Under the object dictionary a position $h\in P'(fs)$ is a morphism out of $fs$ and
$f^\sharp_s(h)$ a morphism out of $s$, so $(f,f^\sharp)$ and
$(\varphi_0,\varphi^\uparrow)$ are literally the same data. Reading $\oplus$ as
reverse composition, (M0)$\Leftrightarrow$(C0) (codomain), (M1)$\Leftrightarrow$(C1)
(unit), (M2)$\Leftrightarrow$(C2) (composition). Identities and composites match,
and objects are in bijection by Theorem~\ref{thm:objdict}; a bijective-on-objects,
fully faithful functor is an isomorphism of categories. The full term-by-term
match is in \texttt{papers/dcont-cof.tex} (Lemma~2.5--Theorem~2.7 there).
\end{proof}

\begin{Proposition}[Functors are the opposite variance]\label{prop:cov}
The \emph{covariant} position maps $\phi_s\colon P(s)\to P'(fs)$ subject to the
mirror conditions (N0)--(N2) of (M0)--(M2) are in bijection with \emph{functors}
$\mathcal C(D)\to\mathcal C(D')$. Thus a directed structure carries two notions of
map of opposite variance: covariant gives functors, contravariant gives
cofunctors, and only the latter is a morphism of the underlying container.
\prov{MacBeth}%
\footnote{Verified combinatorially over $36$ ordered pairs of small categories
drawn from $\{\mathbf1,\mathbf2,\mathbf3,\mathbb Z/2,\mathbb Z/3,M\}$;
\#DCont-morphisms $=$ \#cofunctors in all $36$ cases, \#op-variance-morphisms $=$
\#functors in all $36$, and \#functors $\neq$ \#cofunctors in $17$ of the $36$.
See \texttt{papers/dcont-cof.tex} \S5. The underlying fact that functors are
carried by bicomodules, not comonoid morphisms, is \cite{SS23}.}
\end{Proposition}

\begin{Corollary}[Cofunctors are the Put of a delta lens]\label{cor:lens}
A \emph{delta lens} $A\rightleftarrows B$ is a \emph{Get} functor together with a
\emph{Put} cofunctor sharing one object map, related by the section condition
(PutGet) $f\cdot\varphi(a,u)=u$. Under $\DCont\cong\Cof$ the Put cofunctor is
exactly a morphism of directed containers; so a directed-container morphism is the
update-propagating half of a bidirectional transformation.
\prov{Cited: Clarke 2022 (arXiv:2108.00390), Def.\ 4 / Thm.\ 9}%
\footnote{The section condition is (PutGet) $f\cdot\varphi(a,u)=u$, \emph{not}
``$f\circ\varphi=\id$''; MacBeth flags this because the loose form recurs in the
folklore.}
\end{Corollary}

\section{The dual side: monoids in \texorpdfstring{$\Cont$}{Cont} and the free monoid}
\label{sec:ch6-monoids}

The whole chapter so far has been about \emph{comonoids} for $\triangleleft$. Their
mirror-image, the $\triangleleft$-\emph{monoids}, we have already named: by
Theorem~\ref{thm:comp-comonoid-monoid} a monoid in $(\Cont,\triangleleft,y)$ is exactly a
\emph{polynomial monad} --- a monad on $\Set$ whose functor is a container extension. Where
comonoids gave us categories, monoids give us monads; where the comonoid comultiplication
$\delta$ \emph{split} a state into a state-of-states, the monoid multiplication $\mu\colon
c\triangleleft c\to c$ \emph{composes} two layers of operations into one.

We do not develop the monoid side here --- that is Chapter~\ref{ch:moncomon}. But one monoid
will be the engine of that chapter, and it is cleanest to \emph{state} it now, while the
$\triangleleft$-monoid laws are fresh, and \emph{use} it there. It is the \emph{free}
$\triangleleft$-monoid: the smallest monad built from a bare container.

\begin{Proposition}[The free $\triangleleft$-monoid on a container --- statement]
\label{prop:free-monoid-stmt}
Every container $C=\cont SP$ generates a $\triangleleft$-monoid $C^{*}=\cont{S^{*}}{P^{*}}$,
the \emph{free $\triangleleft$-monoid on $C$}, whose shapes are the well-founded $C$-trees with
a variable-leaf constructor and whose positions are the variable leaves:
\[
  S^{*}\ =\ \mu Y.\,\bigl(1+\textstyle\sum_{s:S}(P(s)\to Y)\bigr),
  \qquad
  P^{*}(\bullet)=1,\quad
  P^{*}(\mathsf{node}(s,c))=\textstyle\sum_{p:P(s)}P^{*}\!\bigl(c(p)\bigr).
\]
Its unit $\eta\colon y\to C^{*}$ picks the leaf $\bullet$, and its multiplication
$\mu\colon C^{*}\triangleleft C^{*}\to C^{*}$ is tree \emph{grafting} (plant a tree at each
variable leaf). The three $\triangleleft$-monoid laws --- $\bullet$ a two-sided unit for
grafting, and grafting associative --- hold.
\prov{Cited: Gambino--Kock 2009, Thm 4.5 (free polynomial monad) \cite{GK};
Abbott--Altenkirch--Ghani 2005 \cite{AAG05}. The three laws in container coordinates are
MacBeth's, Lean-verified (\texttt{Free.lean}).}%
\footnote{The construction $S^{*}$ is a $W$-type (an initial algebra in $\Set$), so it exists
by the initial-algebra theorem, and the position family $P^{*}$ is well-defined by recursion
\emph{over} that $W$-type --- an indexed $W$-type. The full construction, the grafting figure,
the three-law verification, and the \texttt{Free.lean} status are given at
Definitions~\ref{def:freemonad}--\ref{def:grafting} and Theorem~\ref{thm:free-monoid}; the
dependency of $P^{*}$ on $S^{*}$ is spelt out at Remark~\ref{rem:indexed-w}. We state the object
here and build it there.}
\end{Proposition}

\noindent What makes $C^{*}$ the \emph{free} monoid --- that every $\triangleleft$-monoid
receiving a map from $C$ receives a unique monoid map from $C^{*}$ --- is its universal
property, the \emph{free-monad Lemma} (Theorem~\ref{thm:free-lemma}). That is the first result
of the next chapter, where $C^{*}$ reappears as the \emph{syntax} generated by $C$ read as a
signature. We have stated the object; there we prove it is universal.

\chapter{Monads and comonads: the free and cofree constructions}
\label{ch:moncomon}

Read a container as a \emph{signature}. The shapes $S$ are operation symbols; the fibre
$P(s)$ is the \emph{arity} of $s$, the set of argument slots the operation $s$ has. The
container $1+X$ is the signature with one nullary symbol ($\mathsf{nil}$, arity
$\varnothing$) and one unary symbol ($\mathsf{cons}$, arity $1$); the container $1+X^{2}$
is one constant and one binary operation. Now two questions, of the kind a logician and a
control theorist would each ask about a signature:
\begin{itemize}[nosep]
  \item \emph{What is the language it generates?} What are the well-formed \emph{terms} you
        can build from these operations, with variables at the leaves?
  \item \emph{What behaviour can it exhibit?} If a machine can, at each step, perform one of
        these operations and branch according to its arity, what are the possibly-infinite
        \emph{execution trees} it can unfold?
\end{itemize}
The first answer is the \emph{free monad} on the signature; the second is the \emph{cofree
comonad}. They are the syntax and the semantics of the same data --- and the fact that
organises this whole chapter is that \textbf{both answers are again containers}, with shape
and position data you can write down, and that the second one is not merely a comonad but a
\emph{directed container}: a small category. The behaviour of a signature lands back on the
equivalence chain of Chapter~\ref{ch:comonoids}.

\medskip
\noindent\textbf{The promise.} By the end of this chapter you will be able to compute, for
any container $C$, the free monad $C^{*}$ and the cofree comonad $C^{\infty}$ as containers;
you will know why $C^{*}$ is the \emph{universal} monad built from $C$ (the free-monad
Lemma); and you will see that $C^{\infty}$ is the cofree \emph{directed} container on $C$,
whose category is the category of subtrees.

\medskip
\noindent\textbf{The landscape.} Chapters~\ref{ch:algebra} and \ref{ch:comonoids} studied
the monoids and comonoids that already \emph{sit inside} $\Cont$: a monoid for
$\triangleleft$ is a monad, a comonoid is a directed container. This chapter runs the other
way. It does not ask which containers happen to carry a (co)monad structure; it asks how to
\emph{manufacture} one --- the smallest monad receiving a map from a given container $C$, and
the largest comonad mapping to it. In categorical language these are two adjunctions, natural
in the monoid $M$ and the comonoid $D$,
\[
  \Mon(\Cont)\big(C^{*},\ M\big)\ \cong\ \Cont\big(C,\ U M\big),
  \qquad
  \Comon(\Cont)\big(D,\ C^{\infty}\big)\ \cong\ \Cont\big(U D,\ C\big),
\]
the free monad $C\mapsto C^{*}$ left adjoint to the forgetful $U\colon\Mon(\Cont)\to\Cont$,
the cofree comonad $C\mapsto C^{\infty}$ right adjoint to $U\colon\Comon(\Cont)\to\Cont$. The
free monad is built by an \emph{initial algebra} and the cofree comonad by a \emph{final
coalgebra}, so we begin by borrowing exactly that much machinery --- and no more.

\section{The tools we borrow: initial algebras and final coalgebras}
\label{sec:moncomon-tools}

Everything below is built from two dual constructions. This section states them, names the
existence theorems, and stops; the mathematics of the chapter is in the container instances,
not here.\footnote{Neil's steer for this book is that the general fixed-point machinery
belongs in a Preliminaries chapter, with this chapter kept to the container instances. The
book has no Preliminaries chapter \emph{yet}; until the global structure is fixed, this
section carries the borrowed tools in tight, cited form, and can be lifted out wholesale when
the Preliminaries chapter exists.}

Let $\mathcal E$ be a category and $F\colon\mathcal E\to\mathcal E$ an endofunctor. An
\emph{$F$-algebra} is an object $A$ with a map $FA\to A$; a morphism of algebras is a map
commuting with the structure. The \emph{initial} $F$-algebra, when it exists, is written
$\mu F$ with structure map an isomorphism $F(\mu F)\xrightarrow{\ \sim\ }\mu F$ (Lambek's
lemma); dually the \emph{final} $F$-coalgebra $\nu F$ carries an iso
$\nu F\xrightarrow{\ \sim\ }F(\nu F)$. We write $\mu X.\,FX$ and $\nu X.\,FX$ when we want to
name the bound variable.

\begin{Theorem}[Existence by the initial-algebra sequence]\label{thm:adamek}
If $\mathcal E$ has an initial object $0$ and colimits of $\omega$-chains, and $F$ preserves
those colimits, then $\mu F=\operatorname*{colim}_{n}F^{n}0$ exists. Dually, if $\mathcal E$
has a terminal object $1$ and limits of $\omega^{\mathrm{op}}$-chains and $F$ preserves them,
then $\nu F=\lim_{n}F^{n}1$ exists. More generally it is enough that $F$ be
\emph{accessible} --- that it preserve $\kappa$-filtered colimits for some regular cardinal
$\kappa$ --- in which case the sequence is run out to length $\kappa$.
\prov{Cited: Ad\'amek 1974 (initial-algebra sequence); Ad\'amek--Rosick\'y (accessible
functors). Classical.}
\end{Theorem}

The two shapes we need are the ones that turn a plain endofunctor into a monad and a comonad.

\begin{Definition}[Free monad and cofree comonad of an endofunctor]\label{def:tf-df}
For $F\colon\mathcal E\to\mathcal E$ the \emph{free monad} on $F$ is the endofunctor
\[
  T_{F}X\ =\ \mu Y.\,(X+FY),
\]
the initial algebra of $Y\mapsto X+FY$ with $X$ held fixed; its unit is the injection
$X\hookrightarrow T_{F}X$ and its multiplication is substitution at the variable leaves.
Dually the \emph{cofree comonad} on $F$ is
\[
  D_{F}X\ =\ \nu Y.\,(X\times FY),
\]
the final coalgebra of $Y\mapsto X\times FY$; its counit reads off the $X$-label at the root
and its comultiplication relabels each node by the subtree hanging from it.
\prov{Cited: Barr 1970 (free monad on an endofunctor); the cofree comonad is dual. Standard.}
\end{Definition}

That is the whole toolkit. The rest of the chapter is the single observation that when $F$ is
a container extension, both fixed points are computed \emph{inside} $\Cont$ --- and computed
concretely enough to program.

\section{$W$-types and their duals}
\label{sec:moncomon-wtypes}

The initial and final algebras we need are of the special functors that a container itself
determines, and these have a name.

\begin{Definition}[$W$-type and $M$-type]\label{def:wtype}
For a container $\cont SP$ the \emph{$W$-type} is the initial algebra
$W\,S\,P=\mu Y.\ \Ext{\cont SP}Y=\mu Y.\,\textstyle\sum_{s:S}(P(s)\to Y)$: the set of
\emph{well-founded trees} whose nodes are labelled by shapes $s$ and which branch over $P(s)$.
The \emph{$M$-type} is the final coalgebra $M\,S\,P=\nu Y.\,\sum_{s:S}(P(s)\to Y)$: the same
trees, now allowed to be infinitely deep.
\end{Definition}

\begin{Proposition}[$W$- and $M$-types exist in $\Set$]\label{prop:wm-exist}
For every container $\cont SP$ the types $W\,S\,P$ and $M\,S\,P$ exist in $\Set$.
\prov{Cited: Abbott--Altenkirch--Ghani 2005 (containers and $W$-types); Ad\'amek
(Theorem~\ref{thm:adamek}).}
\end{Proposition}

The two halves of this proposition are true for different reasons, and the asymmetry is
worth a moment because it is a payoff of Chapter~\ref{ch:cont}.

For the $W$-type, the point is that $\Ext{\cont SP}=\sum_{s}(P(s)\to-)$ is
\emph{accessible}: each representable $(P(s)\to-)$ preserves $\kappa$-filtered colimits once
$\kappa$ exceeds the cardinality of $P(s)$, and a coproduct of accessible functors is
accessible. So the initial-algebra sequence converges (Theorem~\ref{thm:adamek}) and $W\,S\,P$
is its colimit --- the finite-depth trees, assembled by depth.

For the $M$-type we could invoke the dual of accessibility, but $\Cont$ hands us something
sharper. The final-coalgebra sequence $1\leftarrow F1\leftarrow F^{2}1\leftarrow\cdots$ is a
diagram of \emph{connected} shape, and Chapter~\ref{ch:cont} proved that a container extension
\emph{preserves connected limits} (Theorem~\ref{thm:char}). So $\Ext{\cont SP}$ preserves the
limit that defines $\nu$, the sequence converges on the nose, and $M\,S\,P=\lim_{n}F^{n}1$ is
the set of possibly-infinite trees.

\begin{teachbox}[Why the cofree side is the \emph{easy} side here]
It is a small surprise worth flagging. Coinductive constructions are usually the delicate
ones --- infinite objects, no induction principle, bisimulation instead of equality. But for
containers the \emph{final} coalgebra falls out of a theorem we already own: connected-limit
preservation (Theorem~\ref{thm:char}). The final-coalgebra sequence is connected; container
extensions preserve connected limits; done. The polynomial boundary of Chapter~\ref{ch:cont}
is not a technical sidebar --- it is exactly what makes the behaviour construction go through.
\end{teachbox}

One more remark is needed before the constructions, because the free monad's positions are
themselves defined by a fixed point that lives \emph{over} the tree of shapes.

\begin{Remark}[Positions are an indexed $W$-type]\label{rem:indexed-w}
In the free-monad container below, the shapes $S^{*}$ form a $W$-type in $\Set$, and the
positions $P^{*}$ are then defined by recursion \emph{on that $W$-type} --- an initial algebra
in the slice $\Set^{S^{*}}$ (equivalently, an indexed, or dependent, $W$-type). Such
indexed $W$-types exist whenever the base $W$-type does; this is what makes $P^{*}$
well-defined as a function of the tree, and not merely a heuristic ``count the leaves''.
\prov{Cited: Gambino--Hyland 2004 (dependent polynomial functors, indexed $W$-types).}
\end{Remark}

\section{The free monad of a container}
\label{sec:moncomon-free}

We can now write the free monad down. It is the extension of the free
$\triangleleft$-monoid $C^{*}$ stated at Proposition~\ref{prop:free-monoid-stmt}; here we build
that object explicitly --- shapes, positions, and grafting --- and in the next section prove it
universal. Fix a container $C=\cont SP$.

\begin{Definition}[Free monad container]\label{def:freemonad}
$C^{*}=\cont{S^{*}}{P^{*}}$ where the shapes are the well-founded $C$-trees \emph{with a
variable-leaf constructor},
\[
  S^{*}\ =\ \mu Y.\,\bigl(1+\textstyle\sum_{s:S}(P(s)\to Y)\bigr)
  \ =\ W\,\bigl(1+S\bigr)\,\widehat P,
  \qquad
  \widehat P(\mathsf{inl}\,\ast)=\varnothing,\ \ \widehat P(\mathsf{inr}\,s)=P(s),
\]
generated by a nullary \emph{leaf} $\bullet$ and a node $\mathsf{node}(s,c)$ for each shape
$s$ and each family $c\colon P(s)\to S^{*}$ of subtrees; and the positions are the
\emph{variable leaves},
\[
  P^{*}(\bullet)=1,\qquad
  P^{*}(\mathsf{node}(s,c))=\textstyle\sum_{p:P(s)}P^{*}\!\bigl(c(p)\bigr).
\]
Equivalently $P^{*}(w)$ is the set of paths from the root of $w$ that end at a $\bullet$.
\end{Definition}

\begin{figure}[h]\centering
\begin{tikzpicture}[level distance=9mm,level 1/.style={sibling distance=20mm},
   level 2/.style={sibling distance=10mm},every node/.style={font=\small}]
\node[circle,draw,inner sep=1pt]{$s$}
  child{node[circle,fill=black!12,draw,inner sep=1.2pt]{$\bullet$} edge from parent node[left,font=\tiny]{$p_1$}}
  child{node[circle,draw,inner sep=1pt]{$s'$}
     child{node[circle,fill=black!12,draw,inner sep=1.2pt]{$\bullet$} edge from parent node[left,font=\tiny]{$q_1$}}
     child{node[circle,fill=black!12,draw,inner sep=1.2pt]{$\bullet$} edge from parent node[right,font=\tiny]{$q_2$}}
     edge from parent node[right,font=\tiny]{$p_2$}};
\end{tikzpicture}
\caption{A shape of $C^{*}$ for $C$ with a binary shape $s$ (arity $\{p_1,p_2\}$) and a
binary $s'$: a node $s$ whose first child is a variable leaf and whose second is a node $s'$
with two variable leaves. Its $P^{*}$ has three elements --- the three shaded leaves,
addressed by the paths $\langle p_1\rangle$, $\langle p_2,q_1\rangle$,
$\langle p_2,q_2\rangle$.}\label{fig:freetree}
\end{figure}

The monoid structure is tree \emph{grafting}. Writing trees as $t::=\mathsf{lf}\mid
\mathsf{nd}\,s\,\kappa$ (leaf, or node with $\kappa\colon P(s)\to S^{*}$) and
$\mathrm{leaves}$ for $P^{*}$ --- so that, in the signature-functor presentation
$1+\sum_{s}(P(s)\to-)$ of Definition~\ref{def:freemonad}, $\mathsf{lf}=\mathsf{inl}\,\ast$
and $\mathsf{nd}\,s\,\kappa=\mathsf{inr}(s,\kappa)$, and $P^{*}$ is the indexed
$W$-type of leaves (Remark~\ref{rem:indexed-w}):

\begin{Definition}[Grafting monoid]\label{def:grafting}
The unit $\eta\colon I\to C^{*}$ sends the unique shape of $I=\cont 11$ to $\mathsf{lf}$. The
multiplication $\mu\colon C^{*}\triangleleft C^{*}\to C^{*}$ grafts: given a tree $t$ and a
family $u$ assigning a tree $u_{\ell}$ to each leaf $\ell$ of $t$,
\[
  \mathrm{graft}(\mathsf{lf},u)=u_{()},\qquad
  \mathrm{graft}(\mathsf{nd}\,s\,\kappa,\,u)=\mathsf{nd}\,s\,\bigl(\lambda p.\
    \mathrm{graft}(\kappa_{p},\,u_{p})\bigr),
\]
plants $u_{\ell}$ in place of each variable leaf $\ell$. On positions $\mu^{\sharp}$ is the
canonical bijection
$\mathrm{leaves}(\mathrm{graft}(t,u))\cong\sum_{\ell:\mathrm{leaves}\,t}\mathrm{leaves}(u_{\ell})$
that cuts a leaf-path of the grafted tree at its unique $t$-prefix.
\end{Definition}

\begin{Theorem}[$C^{*}$ is a monoid for $\triangleleft$]\label{thm:free-monoid}
$(C^{*},\eta,\mu)$ satisfies the three $\triangleleft$-monoid laws: $\mathsf{lf}$ is a
two-sided unit for grafting and grafting is associative (with the position bijections
matching on the nose). Hence $\Ext{C^{*}}$ is a monad, and
$\Ext{C^{*}}A\cong\mu X.\,(A+\Ext C X)$ is the free monad on $\Ext C$: the trees with
$A$-labelled leaves.
\prov{Construction cited: Gambino--Kock 2009 (free polynomial monad); Abbott--Altenkirch--Ghani
2005. The three $\triangleleft$-monoid laws in container coordinates are MacBeth's; see the
footnote for their Lean status.}%
\footnote{The laws are checked with zero \texttt{sorry} in \texttt{Free.lean}
(\texttt{Container.freeMonoid}), using only \texttt{Quot.sound}; the associativity law reduces
to a \texttt{split}/\texttt{graft} compatibility (\texttt{graft\_assoc}, \texttt{split\_assoc}).
That source is not yet committed to this repository's \texttt{lean/} tree, so the tag reads
\textsf{[MacBeth]} rather than \textsf{Lean-verified}, per the Preface convention.}
\end{Theorem}

\begin{Example}[$\mathsf{Maybe}$ and binary trees]\label{ex:free}
For $C=1+X$ ($\mathsf{nil}$ with $P=\varnothing$, $\mathsf{cons}$ with $P=1$), a $C$-tree is a
finite list of $\mathsf{cons}$es ending in either $\mathsf{nil}$ or a variable leaf, so
$\Ext{C^{*}}A$ is the free monad on $\mathsf{Maybe}$: lists over $A$ with an optional
$\mathsf{nil}$ terminator. For $C=1+X^{2}$, $C^{*}$ is the container of finite binary trees
with variable leaves, and $\Ext{C^{*}}A$ is binary trees with $A$-labelled leaves --- the
term algebra of one constant and one binary operation. In each case grafting is exactly
substitution of terms for variables.
\end{Example}

\section{The free-monad Lemma}
\label{sec:moncomon-lemma}

Definition~\ref{def:freemonad} builds \emph{a} monad from $C$. The content of this section is
that it builds the \emph{universal} one: every monad receiving a map from $C$ receives a
unique map from $C^{*}$. This is the free-monad Lemma, and in container coordinates its proof
is a single structural induction whose base case is a monoid unit law and whose step is a
monoid associativity law --- the target monoid's own laws, applied, never proved.

Write $U\colon\Mon(\Cont)\to\Cont$ for the functor forgetting a $\triangleleft$-monoid down to
its underlying container. The \emph{insertion of generators} is the container morphism
\[
  \alpha=\alpha_{C}\colon C\longrightarrow U(C^{*}),\qquad
  \alpha_{1}(s)=\mathsf{nd}\,s\,(\lambda p.\,\mathsf{lf}),\quad
  \alpha^{\sharp}_{s}\langle p,()\rangle=p,
\]
sending a shape $s$ to the one-node tree of shape $s$ with all children variable leaves, and
identifying that node's leaves with $P(s)$.

\begin{Theorem}[The free-monad Lemma]\label{thm:free-lemma}
For every container $C$, $\alpha_{C}\colon C\to U(C^{*})$ is universal from $C$ to $U$: for
every $\triangleleft$-monoid $M$ and every container morphism $g\colon C\to U(M)$ there is a
\emph{unique} monoid morphism $\widehat g\colon C^{*}\to M$ with $U(\widehat g)\circ\alpha_{C}=g$.
Equivalently $C\mapsto C^{*}$ is left adjoint to $U$, with unit $\alpha$. Consequently, since
$\Ext{-}$ is strong monoidal and fully faithful, $\Ext{C^{*}}$ is the free monad on the
polynomial functor $\Ext C$.
\prov{Cited: Gambino--Kock 2009, Thm 4.5 (free monad on a polynomial functor; single-variable
case Gambino--Hyland 2004) and Spivak 2022 (arXiv:2202.00534) --- two independent prior-art
anchors for the adjunction. The container-coordinate proof of the universal property is
MacBeth's and is graded \emph{proved}.}%
\footnote{MacBeth's proof is in \texttt{proofs/2026-07-24-free-monad-universal-property.md}:
$\widehat g$ is defined by $W$-type recursion, and that it is a monoid morphism, that the
triangle holds, and that it is unique are each proved by induction on the tree. Partial Lean:
\texttt{FreeUniversal.lean} machine-checks the full triangle $\alpha;\widehat g=g$, the unit
law, and object-level uniqueness (\texttt{Quot.sound} only); the multiplication-homomorphism
law and backward uniqueness are not yet formalised. Not committed to \texttt{lean/}.}
\end{Theorem}

The proof is short to describe even though it is long to check, and the description is the
part worth carrying away. Define $\widehat g$ by recursion on the tree:
\[
  \widehat g_{1}(\mathsf{lf})=e_{M},\qquad
  \widehat g_{1}(\mathsf{nd}\,s\,\kappa)=\mu_{M}\bigl(g_{1}s,\ \lambda q.\,
     \widehat g_{1}(\kappa(g^{\sharp}_{s}q))\bigr),
\]
so a tree is evaluated by reading each node as an $M$-operation and \emph{multiplying up} from
the leaves. That $\widehat g$ is a monoid morphism is one induction: the base case
$\mathrm{graft}(\mathsf{lf},u)=u$ is $M$'s \emph{left-unit} law, and the inductive step,
comparing $\widehat g$ of a graft with the product of the $\widehat g$'s, is $M$'s
\emph{associativity} law. The triangle $\alpha;\widehat g=g$ uses $M$'s \emph{right-unit} law
--- fittingly, since $\alpha$ places a generator at the root with leaf children, exactly the
right-unit configuration. Uniqueness is forced because the position bijection of grafting
(Definition~\ref{def:grafting}) is invertible: any morphism agreeing with $g$ at the
generators is pinned, node by node, by the multiplication law.

\begin{teachbox}[Where the induction really lives]
Notice what is and is not proved. The laws of the \emph{free} monoid $C^{*}$ --- that grafting
is unital and associative --- are Theorem~\ref{thm:free-monoid}, checked once in tree
combinatorics. The free-monad \emph{Lemma} proves something different: that the induced map
into a \emph{given} monoid $M$ is a morphism. That proof never touches a law of $C^{*}$; it
discharges each obligation using a law of $M$. The same tree-induction serves twice --- once
to establish $C^{*}$'s own laws, once to establish morphisms out of $C^{*}$ --- with the roles
of ``prove'' and ``apply'' exchanged. This is the shape of every freeness proof, shown here in
coordinates concrete enough to have been run on a machine and on paper both.
\end{teachbox}

\section{The cofree comonad, and cofree $=$ cofree directed container}
\label{sec:moncomon-cofree}

Now the dual, and the chapter's payoff. Dualising Definition~\ref{def:freemonad} --- swap the
initial algebra for the final coalgebra, products for sums, leaves for nodes --- gives the
cofree comonad, and Proposition~\ref{prop:wm-exist} already guarantees it exists.

\begin{Definition}[Cofree comonad container]\label{def:cofreecomonad}
$C^{\infty}=\cont{S^{\infty}}{P^{\infty}}$ where
$S^{\infty}=\nu Y.\,\sum_{s:S}(P(s)\to Y)=M\,S\,P$ is the set of possibly-infinite
$C$-trees, and $P^{\infty}(w)$ is the set of \emph{nodes} of $w$, presented as finite paths
from the root:
\[
  P^{\infty}(\langle s,c\rangle)\ =\ 1+\textstyle\sum_{p:P(s)}P^{\infty}\!\bigl(c(p)\bigr)
\]
(the empty path $\bullet$, plus the paths continuing into each child $c(p)$). Every path is
finite even when $w$ is infinite.
\end{Definition}

The contrast with the free monad is exact and worth stating, because it is the source of a
persistent confusion. \emph{The free monad's positions are its leaves; the cofree comonad's
positions are its nodes} --- \emph{all} of them, the root included, not the leaves. A tree may
have no leaves at all (if it is infinite) yet always has nodes, and it is the nodes that index
$P^{\infty}$.

\begin{Proposition}[Extension and comonad structure]\label{prop:cofree-ext}
$\Ext{C^{\infty}}A\cong\nu X.\,(A\times\Ext C X)$ is the carrier of the cofree comonad on
$\Ext C$: the possibly-infinite $C$-trees with \emph{every node} labelled in $A$. The counit
$\varepsilon\colon C^{\infty}\to I$ reads the root label; the comultiplication
$\delta\colon C^{\infty}\to C^{\infty}\triangleleft C^{\infty}$ relabels each node $n$ by the
subtree $w/n$ hanging from it, and on positions concatenates paths, $\delta^{\sharp}(n,m)=n\cdot
m$.
\prov{Cited: Niu--Spivak 2023, \emph{Polynomial Functors}, Prop.\ 8.18 \& 8.33, Thm.\ 8.45
(cofree comonoid on a polynomial functor); carrier via Abbott--Altenkirch--Ghani 2005
($M$-types). Container-coordinate structure maps re-derived here.}
\end{Proposition}

Here is the punchline the chapter has been driving at. A cofree comonad is, a priori, just a
comonad --- a comonoid for $\triangleleft$. But Chapter~\ref{ch:comonoids} taught us that a
comonoid for $\triangleleft$ is never \emph{just} a comonoid: it is a \emph{directed
container}, hence a small \emph{category}. So the cofree comonad must secretly be a category.
Which one?

\begin{Theorem}[The cofree comonad is a directed container]\label{thm:cofree-dircont}
$C^{\infty}$ is a directed container (Definition~\ref{def:dcont}) with
\[
  \text{root}\ \ \mathsf o(w)=\bullet\ \text{(empty path)},\qquad
  \text{sub-shape}\ \ w\downarrow n=w/n\ \text{(subtree at $n$)},\qquad
  \text{shift}\ \ n\oplus m=n\cdot m\ \text{(concatenation)}.
\]
Laws \textup{D1--D5} hold; the only non-formal one is $w/(n\cdot m)=(w/n)/m$ behind
\textup{D4}. By Theorem~\ref{thm:comonad} this directed container \emph{is} the comonad of
Proposition~\ref{prop:cofree-ext}: the cofree comonad of a container is the cofree
\emph{directed} container on it. Its category (Theorem~\ref{thm:objdict}) is the
\emph{subtree category}: objects are trees, a morphism $w\to w'$ is a node $n$ of $w$ with
$w/n=w'$, identities are roots (empty paths), and composition is path concatenation.
\prov{Cited: Niu--Spivak 2023, Thm.\ 8.45 (cofree comonoid $=$ tree comonoid);
Ahman--Chapman--Uustalu 2014 (polynomial comonad $=$ category). The explicit
$\mathsf o/\!\downarrow\!/\oplus$ data in D1--D5 form, and their verification, are MacBeth's.}%
\footnote{A Lean formalisation of $C^{\infty}$ as a directed container is \emph{in progress}
but blocked: the cofree carrier is an $M$-type (coinductive), and the repository's Lean core
(v4.30.0, no Mathlib) has no coinduction. Adding Mathlib's \texttt{PFunctor.M} is the intended
route; until then no \textsf{Lean-verified} tag is claimed for the cofree side.}
\end{Theorem}

The identification is concrete enough to read as an algorithm. To compose two subtree-morphisms
--- a node $n$ of $w$ landing at $w/n$, then a node $m$ of $w/n$ landing at $(w/n)/m$ --- you
concatenate the addresses, $n\cdot m$, and land at $w/(n\cdot m)$; law D4 is the assertion that
this equals $(w/n)/m$, i.e.\ that addressing a subtree of a subtree is addressing along the
joined path. That is the whole content, and it is why the category is exactly the poset-like
category of ``$w'$ is a subtree of $w$, witnessed by a specific occurrence''.

\begin{Example}[Colists and infinite binary trees]\label{ex:cofree}
For $C=1+X$ the shapes $S^{\infty}=\nu Y.\,1+Y=\bar{\mathbb N}$ are the conaturals, and
$\Ext{C^{\infty}}A$ is the non-empty colists over $A$ --- the classical
$\mathsf{head}/\mathsf{tails}$ comonad, whose subtree category is $(\bar{\mathbb N},\le)$ read
as ``drop a finite prefix''. For $C=1+X^{2}$, $C^{\infty}$ is the possibly-infinite
$A$-labelled binary trees, $\delta$ relabelling each node by its subtree, and the subtree
category has a morphism $w\to w'$ for every occurrence of $w'$ as a subtree of $w$.
\end{Example}

\subsection{The cofree-comonad Lemma}
\label{sec:moncomon-cofree-lemma}

Definition~\ref{def:cofreecomonad} builds \emph{a} comonad from $C$; as with the free monad,
it builds the \emph{universal} one --- now on the co-side. Every comonad mapping \emph{to} $C$
factors uniquely through $C^{\infty}$. This is the exact dual of the free-monad Lemma, and its
proof is the exact dual induction: where the free proof discharged its obligations using the
laws of the target \emph{monoid}, this one discharges them using the laws D1--D5 of the source
\emph{directed container}.

Write $U\colon\Comon(\Cont)\to\Cont$ for the functor forgetting a comonoid --- equivalently a
directed container, equivalently a small category --- to its underlying container. The
\emph{read-the-root} morphism is
\[
  \pi=\pi_{C}\colon U(C^{\infty})\longrightarrow C,\qquad
  \pi_{1}(w)=\mathrm{root}(w),\qquad
  \pi^{\sharp}_{w}(p)=\langle p,\bullet\rangle,
\]
sending a tree to its root shape and, on positions, sending $p\in P(\mathrm{root}\,w)$ to the
depth-one path ``step to child $p$, stop at its root''. This $\pi$ is the counit of the intended
adjunction, and it is \emph{not} the comonoid counit $\varepsilon\colon C^{\infty}\to I$ of
Proposition~\ref{prop:cofree-ext} (which reads the root \emph{label} into $I$); $\pi$ reads the
root \emph{shape} out into $C$.

\begin{Theorem}[The cofree-comonad Lemma]\label{thm:cofree-lemma}
For every container $C$, $\pi_{C}\colon U(C^{\infty})\to C$ is couniversal from $U$ to $C$: for
every comonoid $D$ (a directed container) and every container morphism $g\colon U(D)\to C$ there
is a \emph{unique} comonoid morphism $\widehat g\colon D\to C^{\infty}$ with
$\pi_{C}\circ U(\widehat g)=g$. Equivalently $C\mapsto C^{\infty}$ is right adjoint to $U$, with
counit $\pi$. Consequently, since $\Ext{-}$ is strong monoidal and fully faithful,
$\Ext{C^{\infty}}$ is the cofree comonad on $\Ext C$.
\prov{Cited: Niu--Spivak 2023, Prop.\ 8.18 \& 8.33, Thm.\ 8.45, and Spivak 2022
(arXiv:2202.00534), Eq.\ (244)--(249) (the adjunction $U\dashv\mathfrak c$ and its limit tower)
--- two independent prior-art anchors. The container-coordinate proof of the couniversal
property is MacBeth's and is graded \emph{proved}.}%
\footnote{MacBeth's proof is in \texttt{proofs/2026-07-25-cofree-comonad-universal-property.md},
the coordinate dual of the free-monad proof. \emph{Lean status: blocked.} The forward map
$\widehat g_{1}$ is defined by corecursion into the $M$-type $S^{\infty}$, and the
repository's Lean core (v4.30.0, no Mathlib) has no coinduction; adding Mathlib's
\texttt{PFunctor.M} is the intended route (an infrastructure decision pending). The
\emph{backward/position} layer --- finite path recursion --- is portable now and is the natural
partial target dual to \texttt{FreeUniversal.lean}.}
\end{Theorem}

The proof is again short to describe. Fix a directed container $D=\cont TQ$ with root
$\mathsf o$, sub-shape ${\downarrow}$ and shift $\oplus$, and a morphism $g\colon U(D)\to C$
(so $g_{1}\colon T\to S$ and, backward, $g^{\sharp}_{\tau}\colon P(g_{1}\tau)\to Q(\tau)$). Write
$\tau_{i}:=\tau\downarrow g^{\sharp}_{\tau}(i)$ for the sub-shape of $\tau$ named by the $i$-th
position of its root image. Define $\widehat g$ in two layers:
\[
  \underbrace{\mathrm{root}(\widehat g_{1}\tau)=g_{1}\tau,\quad
    \mathrm{child}(\widehat g_{1}\tau,i)=\widehat g_{1}(\tau_{i})}_{\text{forward: corecursion on the tree}}
  \qquad
  \underbrace{\widehat g^{\sharp}_{\tau}(\bullet)=\mathsf o_{\tau},\quad
    \widehat g^{\sharp}_{\tau}\langle i,w\rangle
      =g^{\sharp}_{\tau}(i)\oplus\widehat g^{\sharp}_{\tau_{i}}(w)}_{\text{backward: recursion on the finite path}}.
\]
The forward layer \emph{unfolds} $D$'s state $\tau$ into the tree that reads $g_{1}$ at every
reachable sub-shape; the backward layer \emph{folds} a path back down into a single $D$-position
by iterated shift. That $\widehat g$ is a comonoid morphism is one induction on the path: the
comultiplication law's base case is $D$'s unit laws (D1 forward, D2 backward) and its inductive
step is $D$'s nesting and associativity laws (D4 forward, D5 backward); the counit law is the
base clause $\widehat g^{\sharp}_{\tau}(\bullet)=\mathsf o_{\tau}$. The triangle
$\pi\circ U(\widehat g)=g$ turns on D3 --- fittingly the dual of the free triangle's use of the
target monoid's \emph{right}-unit law, since $\pi^{\sharp}$ picks the depth-one path
$\langle i,\bullet\rangle$ and $g^{\sharp}_{\tau}(i)\oplus\mathsf o_{\tau_{i}}=g^{\sharp}_{\tau}(i)$
is exactly D3. Uniqueness has two halves: the forward map is forced by \emph{finality} of the
$M$-type $S^{\infty}$ (any comonoid morphism over $g$ is a coalgebra morphism into the final
coalgebra, hence $=\widehat g_{1}$), and the backward map is then forced, path by path, by the
comultiplication law.

\begin{teachbox}[The coinduction is confined to the shapes]
The free monad's two layers were both finite \emph{induction}: shapes are $W$-type trees
(finite), positions are leaves (finite). The cofree comonad breaks the symmetry in exactly one
place. Its shapes are $M$-type trees, possibly infinite, so the shape layer is genuine
\emph{coinduction} --- and it is used precisely twice, both times as finality of $S^{\infty}$:
once to \emph{define} $\widehat g_{1}$ (an anamorphism) and once to prove it \emph{unique}. But
its positions are \emph{nodes presented as finite paths}, so the entire backward/position layer
--- the transport-heavy half where the directed-container laws D2, D5 do their work --- stays
ordinary finite \emph{induction on the path}. This is the payoff of getting the position count
right (Definition~\ref{def:cofreecomonad}): had positions been leaves, they would not even be
closed under the path-prefix operation the backward induction runs on. Positions-are-nodes is
what lets the dual proof reduce cleanly to $D$'s five laws.
\end{teachbox}

Every step is the coordinate dual of the free-monad Lemma; the dictionary is worth tabulating,
because it makes the free/cofree duality of the whole chapter completely explicit.
\begin{center}
\renewcommand{\arraystretch}{1.2}
\begin{tabular}{@{}p{4.5cm}p{4.7cm}p{4.7cm}@{}}
\toprule
 & \textbf{free monad $C^{*}$ (target monoid $M$)} & \textbf{cofree comonad $C^{\infty}$ (source d.c.\ $D$)} \\
\midrule
carrier shapes      & $S^{*}=W$-type (finite trees)        & $S^{\infty}=M$-type (infinite trees) \\
carrier positions   & leaves (finite)                     & \emph{nodes} = finite rooted paths \\
induced map forward & $\widehat g_{1}$ by $W$-recursion    & $\widehat g_{1}$ by $M$-corecursion \\
induced map backward& $\widehat g^{\sharp}$ via $\mu^{\sharp}$ split & $\widehat g^{\sharp}$ via $\oplus$ (path rec.) \\
morphism law fwd    & $M$-left-unit (base), $M$-assoc (step) & \textbf{D1} (base), \textbf{D4} (step) \\
morphism law bwd    & $M$-left-unit\,/\,assoc backward     & \textbf{D2} (base), \textbf{D5} (step) \\
triangle            & $M$-right-unit                       & \textbf{D3} \\
uniqueness forward   & tree \emph{induction}               & tree \emph{coinduction} (finality) \\
uniqueness backward  & path ind.\ $+$ split invertible     & path ind.\ $+$ comult law \\
\bottomrule
\end{tabular}
\end{center}
The single asymmetry --- induction becomes coinduction only in the \emph{shape} row --- is
precisely the free\,/\,cofree $=$ $W$-type\,/\,$M$-type, leaves\,/\,nodes duality, and nothing
else in the two proofs differs.

\section{Monads on the base, comonads on \texorpdfstring{$\Cont$}{Cont}: the transfer}
\label{sec:moncomon-transfer}

The free monad and the cofree comonad manufacture a $\triangleleft$-(co)monad from a container.
This section runs a different machine, one that takes a monad $M$ living \emph{one level down},
on the base $\Set$, and transfers it to a comonad living \emph{up} on $\Cont$. It is short,
it is exact, and it turns on the one structural fact this book keeps returning to: \emph{positions
are contravariant}.

Let $M=(M,\eta,\mu)$ be any monad on $\Set$ --- $\mathsf{Maybe}$, a writer monad, a free monad,
anything. Postcompose it into the position fibres of a container.

\begin{Definition}[The transferred functor]\label{def:transfer-G}
For a monad $M$ on $\Set$ define $G\colon\Cont\to\Cont$ by
\[
  G\cont SP=\cont S{M\circ P},\qquad G(u,f^{\sharp})=(u,\{M(f^{\sharp}_{s})\}_{s}),
\]
i.e.\ replace each position set $P(s)$ by $M(P(s))$ and each backward map $f^{\sharp}_{s}$ by
$M(f^{\sharp}_{s})$. Equip $G$ with a \emph{counit} $\varepsilon\colon G\Rightarrow\mathrm{Id}$
and a \emph{comultiplication} $\delta\colon G\Rightarrow G^{2}$, both the identity on shapes and,
backward at each shape $s$,
\[
  \varepsilon^{\sharp}_{s}=\eta_{P(s)}\colon P(s)\to M P(s),
  \qquad
  \delta^{\sharp}_{s}=\mu_{P(s)}\colon MMP(s)\to MP(s).
\]
\end{Definition}

The variances line up exactly. A container morphism runs its position map \emph{backward}, so a
\emph{unit} $\eta_{A}\colon A\to MA$ of $M$ points the right way to be the \emph{backward} part of
a counit $\varepsilon$, and $\mu_{A}\colon MMA\to MA$ the right way to be the backward part of a
comultiplication. That is the whole idea, and it forces the theorem.

\begin{Proposition}[The monad$\to$comonad transfer]\label{prop:transfer}
$(G,\varepsilon,\delta)$ is a comonad on $\Cont$. Moreover each comonad law is \emph{exactly}
one monad law of $M$, read backward through the position-contravariance:
\[
  \text{counit-left}\Leftrightarrow M\text{ right-unit},\quad
  \text{counit-right}\Leftrightarrow M\text{ left-unit},\quad
  \text{coassociativity}\Leftrightarrow M\text{ associativity},
\]
and the correspondence is a biconditional: $G$ is a comonad iff $M$ is a monad. Dually, a
\emph{comonad} $W$ on $\Set$ transfers to a \emph{monad} $H\cont SP=\cont S{W\circ P}$ on $\Cont$,
with $\eta^{H}$ backward $=\varepsilon^{W}$ and $\mu^{H}$ backward $=\delta^{W}$.
\prov{MacBeth, proved (\texttt{proofs/2026-07-25-monad-comonad-transfer.md}); Lean-verified,
\texttt{MonadComonadTransfer.lean} (\texttt{Quot.sound} only). Construction $=$ fibred-category
folklore, specialised to $\Cont\to\Set$; distinguished from Ahman--Uustalu below.}
\end{Proposition}

\begin{proof}[Proof sketch]
Everything is identity on shapes, so every equation localises to a family of backward maps in
$\Set$, one per shape $s$; write $A=P(s)$. Functoriality of $G$ is functoriality of $M$
($M(f\circ g)=Mf\circ Mg$); naturality of $\varepsilon,\delta$ is naturality of $\eta,\mu$.
For the three laws, recall that container morphisms compose backward maps in the opposite order,
so a composite's backward map at $s$ is the composite of the pieces read right-to-left. Then:
counit-left is $\mu_{A}\circ\eta_{MA}=\id$, which is $M$'s right-unit law; counit-right is
$\mu_{A}\circ M\eta_{A}=\id$, $M$'s left-unit law; coassociativity is
$\mu_{A}\circ M\mu_{A}=\mu_{A}\circ\mu_{MA}$, $M$'s associativity. The converse uses the
single-shape container $(1,A)$, on which the three comonad laws \emph{are} $M$'s three laws at
the set $A$; letting $A$ range recovers all of $M$'s laws. The dual $H$ is the same computation
with $(\eta,\mu)\rightsquigarrow(\varepsilon^{W},\delta^{W})$. The Lean file discharges each of
the three comonad laws by \texttt{ext'} to the corresponding monad-law field, with no transport.
\end{proof}

\begin{teachbox}[Why it is a duality, in one line]
Send $\Cont$ to its base by $(S,P)\mapsto S$. This is a (bi)fibration whose fibre over $S$ is the
category of position families with \emph{vertical} (backward) maps --- that is,
$(\Set^{\mathrm{op}})^{S}$. ``Positions are contravariant'' \emph{is} the $(-)^{\mathrm{op}}$. A
monad $M$ on $\Set$ is a \emph{comonad} $M^{\mathrm{op}}$ on $\Set^{\mathrm{op}}$; postcomposing
it fibrewise gives a fibred comonad, hence a comonad on the total category $\Cont$. The single
construction $M\mapsto(M^{\mathrm{op}})_{*}$ produces $G$ from a monad and $H$ from a comonad. The
coordinate proof above is this mechanism written out; the fibrational viewpoint on $\Cont$ is von
Glehn's.
\end{teachbox}

There is also a clean \emph{universal} description of $G$, which is where Neil first located it:
$G$ is a left co-closure of the composition product.

\begin{Remark}[$G$ is the $\triangleleft$-left-coclosure with $M$ in the numerator]
\label{rem:transfer-coclosure}
The composition product is left co-closed (Niu--Spivak, \emph{Polynomial Functors}, Prop.~6.57,
after Meyers): there is $\{q/p\}=\sum_{i\in p(1)}y^{\,q(p[i])}$ with
$\Poly(p,\ r\triangleleft q)\cong\Poly(\{q/p\},\ r)$. Taking numerator the monad $M$ and
denominator $p=(S,P)$ gives $\{M/(S,P)\}=\sum_{s}y^{\,M(P s)}=G(S,P)$. For \emph{any} endofunctor
$M$ one has, by Yoneda,
\[
  \Poly\bigl(G p,\ r\bigr)\ \cong\ [\Set,\Set]\bigl(\Ext p,\ r\circ M\bigr)
  \qquad\text{naturally in }r,
\]
so $G(S,P)=\{M/(S,P)\}=\mathrm{Lan}_{(S,P)}M$ is a left Kan extension (Trimble's observation,
Niu--Spivak Ex.~6.63). The comonad data is the coclosure applied to $M$'s structure maps:
$\varepsilon=\{\eta/p\}$ and $\delta=\{\mu/p\}$. On $\Poly$ the transferred comonad
$\hat G$ \emph{applies $M$ to the fibres of the direction bundle}:
$\Ext{G(S,P)}(A)=\sum_{s}(MP s\to A)$, with counit re-exponentiating each fibre along $\eta$ and
comultiplication along $\mu$.
\prov{Cited: Niu--Spivak 2023, Prop.~6.57, Ex.~6.63 \cite{NiuSpivak23}. Identification and
universal property MacBeth's, proved.}
\end{Remark}

\begin{Remark}[Where the transfer sits in the literature]\label{rem:transfer-novelty}
The nearest neighbour is Ahman--Bauer's \emph{Comodule representations of second-order functionals}
\cite{AhmanBauer24}. They work in this very category $\Cont$, with the same covariant extension
$\Ext{\cont A P}X=\sum_{a}(Pa\to X)$ and --- crucially --- the same \emph{contravariant}
cointerpretation $\prod_{a}(Pa\times X)$, which they attribute to Ahman--Uustalu's update monads.
That cointerpretation \emph{is} the fibrewise $(-)^{\mathrm{op}}$ of the teachbox above: the framing
and the machinery are prior art, and we claim neither. Two of their results sit close enough to be
mistaken for the transfer, and are not it.
\begin{itemize}[nosep]
  \item Their monad$\,\leftrightarrow\,$comonad correspondence (Prop.~4.1--4.2) is the \emph{trivial}
        opposite-category duality --- a monad on $\mathcal C$ is a comonad on
        $\mathcal C^{\mathrm{op}}$, ``just a matter of taking opposites.'' It relabels one gadget; it
        does not \emph{apply} a base monad to a container at all.
  \item Their construction of monads \emph{on} $\Cont$ (Thm.~6.3) sends $\cont A P$ to
        $\cont{MA}{P^{\star}}$: the base monad $M$ acts on the \emph{shapes}, the positions are
        merely extended, and --- because, in their words, ``the shape part is independent of the
        position part'' --- the result is a \emph{monad}.
\end{itemize}
The transfer is the mirror image of the second. It applies $M$ to the \emph{positions}, leaves the
shapes fixed, and --- because positions are contravariant --- is \emph{forced} to be a comonad,
namely $G\cont SP=\cont S{M\circ P}=\mathrm{Lan}_{\cont SP}M$ (Remark~\ref{rem:transfer-coclosure}).
Neither the object $\cont S{M\circ P}$ nor the direction ``a base monad on the positions gives a
comonad'' appears in Ahman--Bauer. That shapes-versus-positions, monad-versus-comonad dichotomy ---
sharpened in the next teachbox --- together with the coclosure identity, is what is new here; the
coordinate proof turning each comonad law into the correspondingly-named monad law is its substance.

Two further pointers, neither a scoop. The Topos-Institute account of free PLTL algebras
\cite{ToposPLTL} builds a linking map $\lambda\colon M\mathsf P\Rightarrow\mathsf P I^{\mathrm{op}}$
from a modal monad $M$ acting on \emph{predicates} (through a hyperdoctrine $\mathsf P$) to an
\emph{already-given} comonad $I$ on the base; its slogan --- that the linking obstructs once the
interface branches, degree $\ge 2$ --- rhymes with the transfer's variance story, but the
underlying mathematics is disjoint, since there $M$ never touches positions. Older still, and
running the opposite way, Ahman--Uustalu's \emph{update monads} \cite{AU16} \emph{extract} a base
monad \emph{from} a directed container, while Purdy--Damato's distributive laws compose two
\emph{given} monadic containers horizontally; Hinze's transport of a (co)monad across an adjunction
is a related mechanism one level of abstraction up, with no containers in sight.
\prov{Lead cite: Ahman--Bauer 2024 \cite{AhmanBauer24} (deep-read; Prop.~4.1--4.2 and Thm.~6.3
distinguished). Also: Topos-Institute PLTL blog \cite{ToposPLTL} (deep-read; parallel slogan,
disjoint mathematics); Ahman--Uustalu 2016 \cite{AU16} (update monads, opposite direction);
Purdy--Damato 2025 and Hinze WG2.8 distinguished in prose (not repository sources). Contribution
--- positions$\to$comonad, the coclosure identity, and the coordinate proof --- MacBeth's, proved
and Lean-verified.}
\end{Remark}

\begin{teachbox}[Two ways to feed a $\Set$-monad into a container]
A monad $M$ on $\Set$ can enter a container $\cont S P$ along \emph{either} of its two axes, and the
axis you choose fixes, with no further freedom, whether you get a monad or a comonad.
\begin{itemize}[nosep]
  \item \textbf{Along the shapes:} $\cont S P\mapsto\cont{M S}{P^{\star}}$ applies $M$ to the shape
        set and extends the positions to match (Ahman--Bauer, Thm.~6.3). Shapes are
        \emph{covariant}, so this is a \emph{monad} on $\Cont$.
  \item \textbf{Along the positions:} $\cont S P\mapsto\cont S{M\circ P}$ applies $M$ to each
        position fibre and leaves the shapes alone (the transfer). Positions are
        \emph{contravariant}, so this is a \emph{comonad} on $\Cont$.
\end{itemize}
There is no third axis, and no choice in the outcome: the shape/position axis and the
monad/comonad axis are welded together by the fibrewise $(-)^{\mathrm{op}}$ that makes a container a
container. Apply $M$ where the variance runs forward and monads stay monads; apply it where the
variance runs backward and monads become comonads.
\end{teachbox}

\begin{Example}[$M=\mathsf{Maybe}$]\label{ex:transfer-maybe}
For $M=\mathsf{Maybe}=(1+-)$ the comonad $G$ adjoins a fresh \emph{basepoint} to every position
fibre: $G\cont SP=\cont S{1+P}$. The counit forgets the basepoint (backward $\eta$ includes
$P(s)\hookrightarrow 1+P(s)$) and the comultiplication merges the two basepoints of $1+1+P(s)$
(backward $\mu$). On extensions, $\Ext{G(S,P)}(A)=\sum_{s}A^{1+P s}$: an operation of arity $P(s)$
becomes one of arity $P(s)$ \emph{plus one distinguished slot} --- a default, or exception,
argument. This is the smallest non-trivial instance, and it is exactly the kind of ``add an
ambient slot to every operation'' move an orchestration layer performs on a signature.
\end{Example}

\subsection{The two feeds entwine}
\label{sec:moncomon-entwine}

The teachbox left two liftings of one monad $M$ standing side by side: the shape-monad
$T\cont SP=\cont{MS}{P^{\star}}$ of Ahman--Bauer, and the position-comonad $G\cont SP=\cont S{M\circ P}$
of the transfer. They are built from the same $M$; they act on the same category $\Cont$; one is a
monad and one is a comonad. A reader who has met a distributive law will want to know whether they
\emph{interact} --- whether the monad passes through the comonad coherently. They do, and the way they
do is the cleanest possible: it is $M$'s own failure to be a strong monoidal functor for products,
read on positions. This subsection states the interaction, identifies the law, and --- the part worth
slowing down for --- shows that the interaction runs in \emph{exactly one} direction, obstructed by
branching.

Fix, for this subsection, the $\prod$-cointerpretation of Ahman--Bauer (their \S6.1, \S6.5): writing an
element $m\in MS$ through its support of leaves $b\in\mathrm{lv}(m)$ with labels $x_{b}\in S$,
\[
  P^{\star}(m)\ =\ \prod_{b\in\mathrm{lv}(m)}P(x_{b}).
\]
In Neil's language $P^{\star}$ is a \emph{predicate lifting}: it lifts the position family
$P\colon S\to\Set$ along $M$ to a family $P^{\star}\colon MS\to\Set$, and the product formula is the
\emph{universal} (that is, $\prod$-based) such lifting. The monad unit $\eta^{T}$ is backward the Mendler
projection $i$ out of the singleton product $P^{\star}(\eta_{S}s)=\prod_{\{*\}}Ps=Ps$; the multiplication
$\mu^{T}$ is backward the Mendler restriction $j$ that reindexes a product along the leaf-map of
$\mu^{M}$. The flagship instances of this lifting are $M=\mathsf{Pf}$ (finite powerset),
$M=\mathsf{Maybe}$, $M=\mathsf{Writer}$, and $M=\mathrm{Id}$; the reader monad $A^{K}$ is \emph{not} one,
which is why the lifting must be weak.

\paragraph{The two composites.} Both $TG$ and $GT$ are the identity on shapes over $MS$, so it suffices
to read their positions at a fixed shape $m\in MS$. Localise at $m$, write $Z_{b}:=P(x_{b})$ for the
position set at each leaf, and compute:
\[
  \renewcommand{\arraystretch}{1.3}
  \begin{array}{c|c}
    \text{functor} & \text{positions at } m\\\hline
    G\,T & M\bigl(P^{\star}m\bigr)=M\Bigl(\textstyle\prod_{b}Z_{b}\Bigr)\\
    T\,G & (M\circ P)^{\star}(m)=\textstyle\prod_{b}M(Z_{b})
  \end{array}
\]
($GT$: run $T$, then fatten its positions with $M$. $TG$: fatten first, then take the product-lifting
of the $M$-fattened positions.) Between $M\bigl(\prod_{b}Z_{b}\bigr)$ and $\prod_{b}M(Z_{b})$ there is
exactly one canonical map, the \emph{oplax product-comparison}
\[
  \str_{(Z_{b})}\ \colon\ M\Bigl(\prod_{b}Z_{b}\Bigr)\longrightarrow\prod_{b}M(Z_{b}),
  \qquad w\longmapsto\bigl(M(\pi_{b})\,w\bigr)_{b},
\]
built from $M$ and the product projections alone. It exists for \emph{every} functor $M$ by the
universal property of the target product; it is the structure map exhibiting $M$ as an oplax monoidal
functor for $(\Set,\times)$. No algebra structure on $M$ is used to write it down.

\paragraph{The direction is forced.} $\str$ runs $GT$-positions to $TG$-positions. Container morphisms
run their position map \emph{backward}; so $\str$ is the backward part of a morphism the other way,
\[
  \lambda_{X}\ \colon\ T\,G\,X\longrightarrow G\,T\,X,\qquad
  \text{forward }\id_{MS},\quad\text{backward }\str,
\]
that is, a $2$-cell $\lambda\colon TG\Rightarrow GT$. This is the \emph{standard} orientation of a mixed
distributive law of the comonad $G$ over the monad $T$ (Beck; Power--Watanabe; the entwining-structure
language is Brzezi\'nski--Majid). The naive opposite orientation $GT\Rightarrow TG$ would need the
reverse map $\prod_{b}M(Z_{b})\to M(\prod_{b}Z_{b})$, which is not canonical --- and, we will see,
genuinely does not assemble into a law.

\begin{Theorem}[The two feeds entwine]\label{thm:entwine}
Let $M$ be a $\Set$-monad with the $\prod$-cointerpretation predicate lifting $(-)^{\star}$
(e.g.\ $\mathsf{Pf}$, $\mathsf{Maybe}$, $\mathsf{Writer}$, $\mathrm{Id}$), let $T,G$ be its two container
liftings, and let $\lambda\colon TG\Rightarrow GT$ be the $2$-cell backward $\str$. Then $(T,G,\lambda)$
is an entwining structure on $\Cont$: $\lambda$ is natural and satisfies the four entwining axioms.
Consequently $G$ lifts to a comonad on the category of $T$-algebras and $T$ lifts to a monad on the
category of $G$-coalgebras.
\prov{MacBeth, proved (\texttt{proofs/2026-07-27-monad-comonad-entwining.md}); axioms E1, E3, E4 by
naturality, E2 machine-verified for the $\prod$-lifting examples incl.\ branching $\mathsf{Pf}$. $T=$
Ahman--Bauer Thm 6.3 \cite{AhmanBauer24}; $G=$ the transfer, Prop.~\ref{prop:transfer}. Framework:
Beck \cite{Beck69}, Power--Watanabe \cite{PowerWatanabe02}, Brzezi\'nski--Majid \cite{BrzezinskiMajid98},
standard.}
\end{Theorem}

\begin{proof}[Proof sketch]
Everything is identity on shapes, so each axiom localises to a position-level equation in $\Set$, with
backward maps composing contravariantly; localise at $Z_{b}$ and write $\str=(M\pi_{b})_{b}$. The two
unit/counit axioms and the comultiplication axiom are each one naturality square:
\begin{itemize}[nosep]
  \item \textbf{Counit-$G$} ($\varepsilon_{G}$ against $\lambda$) is $M(\pi_{b})\circ\eta_{\prod Z}
        =\eta_{Z_{b}}\circ\pi_{b}$, the \emph{naturality of $\eta$} at $\pi_{b}$.
  \item \textbf{Unit-$T$} collapses to a single-leaf product, where $\str=M(\id)=\id$ and the Mendler
        $i$ is the singleton projection $=\id$; both sides are $\id_{MZ}$.
  \item \textbf{Comult-$G$} ($\delta_{G}$ against $\lambda$) is $M(\pi_{b})\circ\mu_{\prod Z}
        =\mu_{Z_{b}}\circ MM(\pi_{b})$, the \emph{naturality of $\mu$} at $\pi_{b}$.
  \item \textbf{Mult-$T$} is the only axiom touching the Mendler $j$. Both paths are built from $\str$
        and the reindexing of a product along the leaf-map of $\mu^{M}$; since $\str$ is componentwise
        $M$ applied to a projection, it commutes with that reindexing, and the two backward maps
        coincide. This is the \emph{naturality of $\str$ with respect to product-reindexing}, i.e.\ the
        $j$-naturality diagram of Ahman--Bauer Def.~6.2.
\end{itemize}
Naturality of $\lambda$ in $X$ is naturality of $\str$ and of the projections. The general symbolic
index-chase for Mult-$T$ over an arbitrary weak Mendler $j$ is mechanical and not spelled out; it is
machine-verified for the $\prod$-lifting examples, including the branching commutative case
$M=\mathsf{Pf}$.
\end{proof}

\begin{teachbox}[What $\lambda$ is, in one line]
The entwining says one thing: \emph{$M$ oplax-preserves products, coherently with its unit and
multiplication, transported onto positions.} Fibrationally (the teachbox after
Proposition~\ref{prop:transfer}, made precise in \S\ref{sec:moncomon-fibration}), $G=(M\op)_{*}$ is
vertical and $T$ covers the base monad $M$; $\lambda$ is the \emph{(oplax, generally non-invertible)
Beck--Chevalley mate} comparing \emph{apply $M$ on the base, then fatten positions} against
\emph{fatten positions, then apply $M$ on the base}. The comparison exists one way for free --- that is
$\str$ --- and the four axioms are exactly its coherence with the two monad structures $M$ carries on
the two feeds. It is invertible --- strict Beck--Chevalley --- precisely when $M$ is a \emph{pure writer}
$A\times(-)$, which is strictly stronger than non-branching (\S\ref{sec:moncomon-fibration}); the
interesting content is its \emph{failure}.
\end{teachbox}

\paragraph{One direction only: the branching obstruction.} Here is the part a reader should carry away.
The opposite orientation $GT\Rightarrow TG$ needs the \emph{lax} comparison
$\prod_{b}M(Z_{b})\to M(\prod_{b}Z_{b})$, which not every monad has; a commutative monad supplies the
cartesian one (for $\mathsf{Pf}$, $(S_{b})_{b}\mapsto\prod_{b}S_{b}$, the set of choice-tuples). Even
when it exists, it \emph{fails to be a distributive law} the moment $M$ branches.

\begin{Example}[$\mathsf{Pf}$ breaks the reverse orientation]\label{ex:entwine-branching}
Take $M=\mathsf{Pf}$ and the container $X=\bigl(\{a,b\},\ a\mapsto\{0,1\},\ b\mapsto\{0\}\bigr)$. The lax
$GT\Rightarrow TG$ map satisfies unit-$T$, counit-$G$, and comult-$G$, but fails mult-$T$. At the double
shape $\{\{a,b\},\{a\}\}\in\mathsf{Pf}\,\mathsf{Pf}\,S$ the two paths return, on one side, the
\emph{correlated} two-element set $\{((1,0),(1)),\,((0,0),(0))\}$ --- a union of products --- and on the
other the full four-element product of unions. The multiplication $\mu=\bigcup$ merges the leaf $a$
shared by $\{a,b\}$ and $\{a\}$, and the cartesian lax map cannot see that the two choices at that
shared leaf must agree: \emph{union of products $\neq$ product of unions}.
\end{Example}

The obstruction is \textbf{branching, not non-commutativity}. $\mathsf{Pf}$ is commutative and still
fails; what breaks the law is $\mu$ identifying leaves, which needs a shape with at least two positions
so that $\bigcup$ can overlap. When $M$ has arity at most one ($\mathsf{Maybe}$, $\mathsf{Writer}$,
$\mathrm{Id}$) no two leaves are ever shared, so $\mu$ can identify none, mult-$T$ holds trivially, and
\emph{both} orientations are entwinings at once. The dichotomy is clean.
\[
  \renewcommand{\arraystretch}{1.3}
  \begin{array}{c|c|c}
    M & TG\Rightarrow GT\ (\str,\text{ oplax}) & GT\Rightarrow TG\ (\text{lax})\\\hline
    \text{arity}\le 1\ (\mathsf{Maybe},\mathsf{Writer},\mathrm{Id}) & \text{entwining }\checkmark
      & \text{entwining }\checkmark\ (\text{no overlap})\\
    \text{genuine branching}\ (\mathsf{Pf},\mathsf{List},\dots) & \text{entwining }\checkmark\ (\textbf{universal})
      & \text{fails mult-}T\ \text{\ding{55}}
  \end{array}
\]
So the interaction of the two feeds is real, and it lives in the standard orientation, where it holds for
\emph{all} $M$; the orientation one might guess first is the one that branching forbids. This is the
pedagogically interesting fact: the variance that forced $G$ to be a comonad (positions are
contravariant) also forces the \emph{single} direction in which the monad feed can pass through it.

\paragraph{From a failed law to an arrow calculus.} A distributive law that fails looks, at first,
like a dead end. It is not. The reverse comparison $\kappa\colon GT\Rightarrow TG$ has a job to do
that does not require it to be a global law, and reading it through that job is where the two feeds
pay off. Name it: write
\[
  \kappa\ \colon\ G\,T\Longrightarrow T\,G,\qquad
  \text{forward }\id_{MS},\quad\text{backward the lax comparison }\textstyle\prod_{b}M(Z_{b})\to M(\prod_{b}Z_{b}),
\]
the $2$-cell of the right-hand column of the dichotomy table. The morphisms it composes are the
\emph{effect--coeffect arrows}. An arrow $p\rightsquigarrow q$ is a container map that reads its input
through the coeffect comonad $G$ and writes its output through the effect monad $T$ --- exactly the
shape of a computation that consults a context and returns with a side effect.

\begin{Definition}[The arrow profunctor $\Arr_{M}$]\label{def:arrowprof}
For a monad $M$ with the two feeds $T=T_{M}$, $G=G_{M}$, define
\[
  \Arr_{M}(p,q)\ :=\ \Cont\bigl(G_{M}\,p,\ T_{M}\,q\bigr).
\]
Because $G_{M}$ and $T_{M}$ are \emph{functors} on $\Cont$ (the transfer,
Proposition~\ref{prop:transfer}; Ahman--Bauer, Theorem~6.3), this is a
\emph{profunctor} $\Arr_{M}\colon\Cont\op\times\Cont\to\Set$ --- contravariant in $p$ by
precomposition with $G_{M}$, covariant in $q$ by postcomposition with $T_{M}$. An element of
$\Arr_{M}(p,q)$ is an effect--coeffect arrow $p\rightsquigarrow q$. The \emph{intended identity}
$p\rightsquigarrow p$ is $\eta^{T}_{p}\circ\varepsilon_{p}\colon G_{M}p\to p\to T_{M}p$, and the
\emph{intended composite} of $f\colon G_{M}p\to T_{M}q$ with $g\colon G_{M}q\to T_{M}r$ runs the
coeffect out, threads the first effect past the second coeffect through $\kappa$, and multiplies the
effects:
\[
  G_{M}p\xrightarrow{\ \delta_{p}\ }G_{M}G_{M}p
        \xrightarrow{\ G_{M}f\ }G_{M}T_{M}q
        \xrightarrow{\ \kappa_{q}\ }T_{M}G_{M}q
        \xrightarrow{\ T_{M}g\ }T_{M}T_{M}r
        \xrightarrow{\ \mu^{T}_{r}\ }T_{M}r.
\]
\prov{MacBeth. $\Arr_{M}$ as a profunctor is immediate from functoriality of the two feeds; the
biKleisli identity and composite are the standard ones for a comonad--monad pair welded by a
compatible $\kappa$.}
\end{Definition}

The profunctor is the free part. It exists, with these exact hom-sets, for \emph{every} monad
$M$: the underlying \emph{data} of ``an effect--coeffect arrow from $p$ to $q$'' is always available,
because it is nothing but a container map out of $G_{M}p$ into $T_{M}q$, and $G_{M}$, $T_{M}$ are
always there. What is \emph{not} free is that the composite above be associative and unital ---
that $\Arr_{M}$ be the hom-profunctor of an actual \emph{category}, with $\Cont$ mapping into it
identity-on-objects as its pure part (a Freyd category over $(\Cont,\times)$). That last step needs
$\kappa$ to be coherent, and $\kappa$ is the very reverse comparison the dichotomy table said branching
forbids. So the boundary is already drawn: the category structure on $\Arr_{M}$ exists precisely when
$M$ does not branch. We state the classification where it belongs, alongside the other two modes of
composition, in Theorem~\ref{thm:arrows} of \S\ref{sec:threemodes}; here we record only the shape of
the answer.

\begin{teachbox}[The profunctor is free; the category is what branching costs]
Every monad $M$ gives an effect--coeffect \emph{profunctor} $\Arr_{M}=\Cont(G_{M}-,\,T_{M}=)$ on
$\Cont$ --- no conditions, because both feeds are functors. Turning that profunctor into a
\emph{category} --- giving arrows an associative composition --- requires the compositor
$\kappa$, and $\kappa$ is a coherent law \emph{iff $M$ is non-branching} (every shape has arity
$\le 1$; Theorem~\ref{thm:arrows}). Branching does not destroy the arrows; it destroys their
\emph{composition}. And the non-branching monads of this ($\prod$-liftable) class are no
curiosity: they are exactly the \emph{writer-with-exceptions} monads $MX\cong E+A\times X$
(a monoid $A$ of logs acting on a set $E$ of exceptions), the effect monads of
logging-that-may-fail. The arrow calculus of a computation
that writes a log and may abort composes; the arrow calculus of a computation that branches does
not. That is the price of branching, stated as a boundary rather than a defect.
\end{teachbox}

\begin{Remark}[Scope, and what is open]\label{rem:entwine-scope}
The theorem is stated, and proved, for the $\prod$-cointerpretation predicate lifting, because $\str$
uses the product structure of $P^{\star}$. Whether a canonical $\lambda$ exists for a general (non-$\prod$)
weak Mendler algebra is open: there is no evident map $M(P^{\star})\to(M P)^{\star}$ without the product.
Within the $\prod$-class the content is complete; the general Mendler case, and a from-scratch uniqueness
proof for $\lambda$, are not claimed. What is new here is not the entwining formalism --- that is Beck,
Power--Watanabe, and Brzezi\'nski--Majid --- but the specific law $\lambda=\str$ welding one monad's two
container feeds, its forced orientation, and the branching obstruction that Example~\ref{ex:entwine-branching}
exhibits.

A nearby line of work should be distinguished from ours, lest a reader mistake it for prior art. Spivak's
\emph{Categories by Kan extension} \cite{SpivakKanExt} also builds structure from distributive laws of
monads over comonads --- but there the carrier is $(\Set,\triangleleft)$, the density comonad
manufactures a \emph{new base category} as the law's output (Lawvere theories, $\Delta\op$), and the
orientation is monad-over-comonad producing an object of $\Cathash$. Our $\lambda$ instead entwines two
liftings of a \emph{fixed} base monad $M$ on the \emph{fixed} category $\Cont$, with output a comonad on
$T$-algebras rather than a fresh category. Different carrier, orientation, and output; the overlap is
vocabulary.\footnote{The distinction rests on \cite{SpivakKanExt}'s statement of results; a full deep-read
is a standing to-do, so the comparison is drawn at the level of carrier and orientation rather than proof.}
\prov{MacBeth, proved for the $\prod$-class; general Mendler case open (registry: dead-end).
Contribution over prior art: the law, its orientation, and the obstruction.}
\end{Remark}

\subsection{The fibrational picture: why the two feeds had to be what they are}
\label{sec:moncomon-fibration}

Everything in this section has turned on one phrase --- \emph{positions are contravariant} --- invoked
first to force $G$ to be a comonad, then to fix the single direction of $\lambda$. It is time to name the
structure that phrase describes. There is exactly one, a fibration, and once it is in view the two feeds,
the entwining, and the branching obstruction stop being three separate computations and become three
features of one picture. None of the fibrational machinery here is ours; it is standard, and we cite it.
What is ours is only the observation of \emph{which} standard feature each phenomenon of this section is.

\paragraph{The container fibration.} Forget the positions of a container and keep its shapes:
\[
  p\colon\Cont\longrightarrow\Set,\qquad \cont SP\mapsto S,\qquad (u,f^{\sh})\mapsto u.
\]
The fibre over a set $S$ --- the containers with shape set exactly $S$ and the morphisms lying over
$\id_{S}$ --- is the category of position families $P\colon S\to\Set$ with \emph{vertical} (backward) maps
$\{Q(s)\to P(s)\}_{s}$; that is, it is $(\Set\op)^{S}$. This is where ``positions are contravariant''
becomes literal: the fibre is a power of $\Set\op$, and the fibrewise $(-)\op$ is the entire source of every
variance flip in this chapter. A morphism $(u,f^{\sh})\colon\cont SP\to\cont TQ$ is \emph{$p$-cartesian}
exactly when each backward map $f^{\sh}_{s}\colon Q(u\,s)\to P(s)$ is a bijection --- when it merely
\emph{re-indexes} positions without discarding or duplicating them, so that $P\cong Q\circ u$. The cartesian
lift of $\cont TQ$ along a shape map $u\colon S\to T$ is the reindexing
\[
  u^{*}\cont TQ=\cont S{Q\circ u}\qquad\text{(precompose the position family)},
\]
and because $\Set$ is locally cartesian closed, $u^{*}$ has adjoints on both sides,
$\Sigma_{u}\dashv u^{*}\dashv\Pi_{u}$: $p$ is a split bifibration. Every construction below is read against this
one gadget.
\prov{Cited: von Glehn \cite{vonGlehn18} ($\Cont\cong\int_{S}(\Set/S)\op$, so
$\Cont(q)=\Sigma_{q}(q\op)$); Streicher \cite{Streicher} for the fibred-category language and the
fibrewise-opposite. We claim none of the fibration itself.}

\paragraph{$G$ is a vertical fibred comonad.} Read the transfer (Proposition~\ref{prop:transfer}) through
$p$. The comonad $G$ is \emph{vertical} --- $p\circ G=p$, it never moves a shape --- and fibrewise it is
postcomposition by $M\op$, the pushforward $(M\op)_{*}$ along $M\op\colon\Set\op\to\Set\op$. Since a monad
$M$ on $\Set$ \emph{is} a comonad $M\op$ on $\Set\op$, pushing it into every fibre gives a \emph{vertical
fibred comonad} in Jacobs's sense (Ex.~1.7.9 \cite{Jacobs99}, the vertical case: a comonad in the
$2$-category $\mathbf{Fib}(\Set)$ of fibrations over a \emph{fixed} base, so with base component $\id$), and
a fibred comonad is automatically a comonad on the total category $\Cont$. That is the one-line proof of Proposition~\ref{prop:transfer} --- and it explains the odd
biconditional there. Each comonad law of $G$ is fibrewise, hence is the same-named law of the fibre comonad
$M\op$; localised at the representing fibre $(1,A)$ it is the same-named \emph{monad} law of $M$ read through
$(-)\op$. Counit is unit, comultiplication is multiplication, coassociativity is associativity: the
correspondence is the fibrewise op and nothing more. Moreover $G$ preserves cartesian morphisms for
\emph{every} $M$, with no hypothesis at all --- $u^{*}G=G\,u^{*}$, both sending $Q$ to $M\circ Q\circ u$ ---
so $G$ is a \emph{cartesian} vertical fibred comonad, and its counit and comultiplication are stable under
reindexing. Three of the four entwining axioms of Theorem~\ref{thm:entwine} (the two involving
$\varepsilon_{G},\delta_{G}$, and the unit of $T$) are \emph{exactly} this reindexing-stability; they cost
nothing once $G$ is seen to be fibred.

\paragraph{$T$ lies over the base monad $M$.} The shape feed sits differently, and the difference is the
whole story. $T$ does move shapes: $p\circ T=M\circ p$ strictly, and its unit and multiplication cover
$\eta^{M},\mu^{M}$. So $T$ does not live inside one fibre --- it \emph{lies over} the base monad $M$, as a
\emph{monad morphism}, or a \emph{lifting of $M$ along $p$}: a monad on the total category $\Cont$ whose
image under $p$ is the honest monad $M$ (Street's formal theory of monads \cite{Street72}; the general
apparatus for such liftings along a fibration --- logical-predicate liftings, $\top\top$-liftings --- is
Hermida's \cite{Hermida93} and Katsumata's \cite{Katsumata13}). This lifting is \emph{unconditional}; it
exists for every $M$. What is \emph{not} unconditional is the stronger word ``fibred'', and it is worth
saying exactly why, because the word carries the crown boundary of this chapter inside it. There are two
grades of ``fibred'' monad, and $T$'s home is not $G$'s:
\begin{itemize}[nosep]
  \item A \emph{vertical} fibred monad --- a monad in $\mathbf{Fib}(\Set)$, with base component $\id$
        (Jacobs, Ex.~1.7.9 \cite{Jacobs99}). This is $G$'s home, dualised; it is never $T$'s, since $T$
        moves shapes.
  \item A fibred monad over a \emph{nontrivial} base monad --- a monad in the larger $2$-category
        $\mathbf{Fib}$ whose base component is $M\neq\id$ (Hermida, Def.~5.4.1 \cite{Hermida93}). This is
        the right shape for $T$. But Hermida's definition demands that the \emph{total} functor be a
        \emph{fibred functor}, i.e.\ preserve cartesian morphisms --- and $T_{M}$ does that exactly when
        $M$ is a cartesian monad (the crown boundary, made precise in the stratification box below).
\end{itemize}
So the honest statement is: \emph{$T_{M}$ is a lifting of $M$ for every $M$, and a fibred monad over $M$
in Hermida's sense (Def.~5.4.1) precisely when $M$ is cartesian.} The word ``fibred'' is not decoration;
it bakes in cartesian-preservation, which is the very boundary this chapter keeps circling. What Neil calls
the \emph{predicate lifting} --- the product formula $P^{\star}(m)=\prod_{b}P(x_{b})$ --- is this lifting
in coordinates; we claim none of the apparatus, only the identification of which grade $T_{M}$ occupies.
And here is the asymmetry to carry away: $G$ is
cartesian for \emph{every} $M$, while $T$ is cartesian \emph{only} when $M$ is. All the branching fragility
of this chapter lives on the shape feed $T$; the position feed $G$ is serene.

\paragraph{$\lambda$, honestly: a mixed distributive law, not a Beck--Chevalley square.} Now the entwining
law. It is tempting to christen $\lambda$ a \emph{Beck--Chevalley mate} and be done --- the picture rhymes
with one --- but the word would be a genuine category error, and it is worth spelling out why. Beck--Chevalley,
in the fibrational sense (Jacobs, Def.~1.9.4 \cite{Jacobs99}), is a statement about the \emph{adjoints to
reindexing} $\Sigma_{u}\dashv u^{*}\dashv\Pi_{u}$: for every pullback square in the base, the canonical mate
comparing quantification with \emph{substitution} is an isomorphism. That is not what $\lambda$ is. The
$2$-cell $\lambda\colon TG\Rightarrow GT$ is a \emph{mixed distributive law} --- an \emph{entwining} of the
comonad $G$ with the lifting $T$ (Beck \cite{Beck69}; Power--Watanabe \cite{PowerWatanabe02}) --- which is
exactly the content of Theorem~\ref{thm:entwine}. It compares $M$ with a fibrewise product, not quantification
with reindexing; conflating the two constructions is the error the first draft made.

What \emph{is} true, and is the honest residue of the Beck--Chevalley intuition, is a statement about when
$\lambda$ is \emph{invertible}. Fibrewise the predicate lifting $P^{\star}(m)=\prod_{b}P(x_{b})$ is a
$\Pi$-pushforward --- a product $\Pi$ along the map $\mathrm{lv}(m)\to 1$ collapsing the leaves of $m$ --- and
$\lambda$'s fibrewise component is the canonical oplax comparison asking whether applying $M$ downstairs
commutes with taking that product upstairs. This comparison always exists --- it is built from $\str$ --- but
it is \emph{oplax}, and \emph{in general not invertible}. Now one must be careful, because it is tempting ---
and wrong --- to read off a single condition that makes it iso. The comparison is invertible, a
\emph{Beck--Chevalley-style} coherence (though not Def.~1.9.4 itself), exactly when $M$ preserves \emph{all}
those products, \emph{including the empty one}. The empty product is the nullary shape: $\str$ there is $M(1)\to 1$, and for this to be a
bijection $M$ must fix the singleton, $M(1)\cong 1$. So strict BC forces $M$ to be a \emph{pure writer}
$A\times(-)$ (arity constantly $1$, no nullary shape) --- and this is \emph{strictly stronger} than
non-branching. The exception monad $\mathsf{Maybe}=(1+-)$ is non-branching, yet its $\lambda$ is
\emph{not} invertible: at the nullary shape $\str$ is $(1+1)\to 1$, which is not injective. So the honest
statement is not ``$\lambda$ is Beck--Chevalley'' and not the collapse ``strict BC $=$ non-branching''; it
is that strict BC sits one level \emph{inside} non-branching, and the interesting content is the
\emph{failure} at each level --- $\mathsf{Maybe}$ breaks strict BC without branching; $\mathsf{Pf}$
(Example~\ref{ex:entwine-branching}) breaks even the reverse orientation, by branching.

\begin{teachbox}[The fibration stratifies the monad zoo; it does not collapse it]
The first draft of this book's payoff was a slogan --- ``containers preserve cartesian morphisms $=$ $M$
non-branching $=$ strict Beck--Chevalley'' --- and it is \emph{false}: those are three \emph{different}
levels, with named monads separating each pair. What the fibration actually buys is sharper and more
honest: a strict, four-step \emph{stratification} of the ($\prod$-Mendler) monad zoo,
\[
  \underbrace{A\times(-)}_{\text{pure writer}}
  \ \subsetneq\ \underbrace{E+A\times(-)}_{\text{writer}+\text{exception}}
  \ \subsetneq\ \{\,\text{cartesian monads}\,\}
  \ \subsetneq\ \{\,\text{polynomial}/\prod\text{-Mendler monads}\,\},
\]
each inclusion \emph{strict}, with a named splitter:
\begin{itemize}[nosep]
  \item \textbf{Pure writer} $A\times(-)$: $\lambda$ invertible $=$ \emph{strict} Beck--Chevalley. Splitter
        into the next level: $\mathsf{Maybe}=1+(-)$ is non-branching but has a nullary shape, so its
        $\lambda$ is not strict.
  \item \textbf{Writer-with-exception} $E+A\times(-)$ $=$ the \emph{non-branching} monads (arity $\le 1$):
        here, and here exactly, the reverse compositor $\kappa$ exists and the effect--coeffect arrows form
        a \emph{category} (Theorem~\ref{thm:arrows}). \emph{This is the one true biconditional} ---
        non-branching $\Leftrightarrow$ reverse $\kappa$ --- the arrow core of \S\ref{sec:threemodes}.
        Splitter into the next level: $\mathsf{List}$ (the free monoid) is a \emph{cartesian} monad, but it
        \emph{branches} (arity $\mathbb N$), so it is non-branching's outside.
  \item \textbf{Cartesian monads}: exactly the $M$ for which \emph{$T_{M}$ preserves cartesian morphisms}
        (no leaf-merging in $\mu^{M}$). $\mathsf{List}$ lives here but not below; $\mathsf{Pf}$ (whose
        $\mu=\bigcup$ merges leaves) lives \emph{outside} even this.
  \item \textbf{Polynomial / $\prod$-Mendler monads}: the ambient class on which the whole apparatus is
        defined.
\end{itemize}
So $G_{M}$ preserves cartesian morphisms for \emph{every} $M$ (the vertical-fibred-comonad fact, Lean-certified);
Neil's fibred intuition ``containers preserve cartesian morphisms'' is exactly right \emph{on the $G_{M}$
side}. The error was to carry it over to the $T_{M}$ side and equate it with non-branching and with strict
BC. The fibration's real gift is that it makes these four levels --- and the three monads that separate them
--- \emph{visible as one ladder}, rather than collapsing them into a single false line.

\smallskip
\textbf{One rung further out --- two fibrations, subobject and codomain.} The outermost class is still not
the edge, and to see past it we must name a choice that has been implicit all along: \emph{along which
fibration over $\Set$ are we lifting $M$?} There are two, and they are genuinely different fibrations.
\begin{itemize}[nosep]
  \item The \textbf{subobject fibration} $\mathrm{Sub}(\Set)\to\Set$: over a set $X$ sit the
        \emph{subobjects} $P\hookrightarrow X$, and a leaf carries a \emph{truth value} --- whether or not
        it lies in $P$.
  \item The \textbf{codomain / families fibration} $\Cont\to\Set$ --- the container fibration $p$ of this
        very section: over a shape set sit the position families, and a leaf carries an \emph{actual set of
        positions} --- a datum, not merely a truth value.
\end{itemize}
This is the distinction Hermida and Jacobs draw in their fibrational account of induction and coinduction
\cite{HermidaJacobs98}, and --- Neil's own words --- it is the axis this rung is really about. It is
tempting to summarise the boundary as ``the subobject lift survives, the codomain lift dies.'' That is the
right first cut but too coarse, and the honest statement, below, is a $2\times2$.

\emph{First, the headline is a failure.} The plausible ``core theorem'' one would want ---
\textbf{\emph{every $\Set$ monad lifts to a container monad via $T_{M}$}} --- is \textbf{false}. It holds
for cartesian $M$ and fails for the reader monad $A^{K}$ and for state, because their $\mu^{M}$ \emph{drops}
leaves: Reader keeps only the diagonal $mm[e][e]$, discarding every off-diagonal entry, so the Mendler
restriction $j$, which would have to source a position at \emph{every} inner leaf, has nothing to reindex a
dropped leaf from, and no $T_{M}$-monad forms.

Now the one sentence that made this land, and it lives on a single row --- fix the ``for all leaves''
quantifier and \emph{flip the fibration}:
\begin{quote}\itshape
A ``for all leaves'' \emph{proposition} only gets \emph{easier} to satisfy when leaves are dropped, so the
subobject lift $\Box P(m)=\forall b.\,P(x_{b})$ survives; but a \emph{data}-valued assignment to
\emph{every} leaf cannot be reindexed through a drop --- there is no datum to carry at the vanished leaf ---
so the codomain lift $\prod_{b}P(x_{b})=T_{M}$ dies.
\end{quote}
Make it fully concrete on Reader with $E=\{0,1\}$ and the diagonal multiplication. A double element
$mm\in(X^{E})^{E}$ has four inner entries; $\mu^{M}$ keeps the diagonal $(0,0),(1,1)$ and \emph{drops} the
off-diagonal $(0,1),(1,0)$. The proposition ``all four entries lie in $P$'' certainly implies ``the two
surviving entries lie in $P$'': $\Box$ lifts. But a rule that assigns a \emph{position} to each of the four
entries has, after the drop, no coordinate left to hand to the vanished $(0,1)$: $\prod$ has no
multiplication. Same drop, opposite fates --- because a truth value is only \emph{weakened} by forgetting a
conjunct, while a datum, once its leaf is gone, has nowhere to live.

\emph{But the codomain fibration does not die uniformly}, and this is the correction the $\forall$-row
alone hides. Line up all four canonical ways to lift a leaf-family along $\mu^{M}$: the two \emph{codomain}
(data-valued) lifts $\prod_{b}P(x_{b})$ and $\coprod_{b}P(x_{b})$, and the two \emph{subobject}
(truth-valued) lifts $\Box P=\forall b.\,P(x_{b})$ and $\Diamond P=\exists b.\,P(x_{b})$. Write $R$ for the
leaf-matching relation --- inner tokens against surviving leaves, matched on label. Each lift's
multiplication exists exactly when $R$ is total in a direction fixed by a clean $\mathbb Z/2$ parity:
\[
  \text{direction}\;=\;(\text{is-a-limit})\ \oplus\ (\text{records data, not truth}),
\]
\[
\renewcommand{\arraystretch}{1.3}
\begin{array}{r|cc}
 & \text{records data (codomain)} & \text{records truth (subobject)}\\\hline
\forall\ \text{(limit)}   & \prod\ (=T_{M},\ \text{Ahman--Bauer}):\ \text{forward} & \Box:\ \text{reverse}\\
\exists\ \text{(colimit)} & \coprod\ (=T^{\Sigma}):\ \text{reverse}                & \Diamond:\ \text{forward}
\end{array}
\]
where forward-total means \emph{dies on a drop} and reverse-total means \emph{survives a drop}. (A
brute-force check over all predicates confirms the grading with zero mismatches.) So \textbf{neither
fibration survives uniformly}: the survivor in the subobject world is the $\forall$-lift $\Box$, and the
survivor in the codomain world is the $\exists$-lift $\coprod$ --- data at \emph{one} chosen leaf, not at
all of them. Reader and State are reverse-total (each surviving leaf \emph{is} an inner token --- the
diagonal $(e,e)$ for Reader, the threaded $(s,h\,s)$ for State) and forward-partial (they drop), so both
reverse-column lifts $\Box,\coprod$ survive and both forward-column lifts $\prod,\Diamond$ die.
\emph{Merging} (as in $\mathsf{Pf}$, whose $\mu=\bigcup$ invents no new labels) is total both ways, so all
four survive; \emph{dropping} kills exactly the two forward ones.

In particular, losing the $\prod$-monad does \emph{not} cost Reader every data-recording lift. The
\emph{$\Sigma$-container} lift
\[
  T^{\Sigma}_{M}\cont SP=\cont{MS}{P^{\Sigma}},\qquad P^{\Sigma}(m)=\coprod_{b\in\mathrm{lv}(m)}P(x_{b}),
\]
with unit the codiagonal fold and multiplication a reindexing along the diagonal/threading section, is a
genuine monad on $\Cont$ for Reader (every $E$) and State (every state set). It is the $\coprod$ partner of
Ahman--Bauer's $\prod$, and it is the concrete answer to Neil's question ``does Reader have \emph{any}
lifting that records positions, not just truth values?'': yes --- the $\Sigma$ one, not the $\prod$ one.
(``Records positions rather than truth'' is exactly what the type theorist calls
\emph{proof-relevant}; we use Neil's fibrational wording throughout and log the synonym here, once.)

\smallskip
This boundary carries an independent machine-checked witness. In Orestis's Agda development of liftings to
$\Cont$ \cite{OrestisAgda}, the $\forall$-over-leaves codomain lift (his \texttt{\textup{$\Box$}Lift}, a type-valued
$\prod$) is \emph{proved} to yield a container monad only when the leaf-reindexing map $\mathrm{pr}$ is
split surjective --- i.e.\ only when $\mu$ drops nothing --- with \texttt{Reader} named, on the nose, as the
counterexample (\texttt{CoLift.agda}, ll.\ 163--184, comment l.\ 175); his state and nondeterminism
examples route instead through the surviving $\exists$/$\coprod$ lift. One honest caveat: his framework is
\emph{entirely} type-valued, so his $\Box$ and $\Diamond$ (his notation --- \emph{not} the truth-valued
$\Box,\Diamond$ of the table above) are our $\prod$ and $\coprod$, both sitting in the \emph{codomain}
column. He therefore witnesses the $\prod$-dies\slash$\coprod$-survives split \emph{within} the codomain
fibration; he has no truth-valued box, so he does \emph{not} touch the subobject leg. Even so,
it is the same boundary, reached independently and checked by a proof assistant.
\prov{MacBeth, proved. The $\prod$/$\Box$ boundary is
\texttt{proofs/2026-08-07-proof-relevance-boundary.md}, with the companion census
\texttt{proofs/2026-08-06-state-reader-ladder-census.md} (Lean-verified); the $\mathbb Z/2$ grading is the
brute-force \texttt{fourfold.py} (all predicates, zero mismatches). The $\Sigma$-monad of Reader (every $E$)
and State (every state set) is \emph{proved} (\texttt{proofs/2026-08-07-sigma-monad-proved.md}) and
\emph{Lean-verified for Reader}: a reduction lemma turns each monad law into one label-preserving index
identity (the backward maps are all coproduct-pushforwards $\Sigma_{\alpha}$, and $\coprod$ is faithful,
tested at $P\equiv 1$), discharged by the diagonal section (Reader) or the threaded section (State). The
subobject-vs-codomain fibration distinction is Hermida--Jacobs \cite{HermidaJacobs98}; the $\prod$-lift is
Ahman--Bauer \cite{AhmanBauer24}; the $\Box$/predicate-lifting apparatus is Hermida \cite{Hermida93} and
Katsumata \cite{Katsumata13}; the independent Agda witness is \cite{OrestisAgda}. \emph{Resolved} in
\S\ref{sec:sigma-is-comp}: the general implication ``reverse-total $\Rightarrow$ $\Sigma$-monad'' is
\emph{false} --- the true criterion is that $M$ be a $\triangleleft$-monoid (Corollary~\ref{cor:sigma-monad}),
with $\mathsf{Bag}$ separating the two (Theorem~\ref{thm:bag-refutes}); reverse-totality is only the
pointwise shadow of that polynomial fact. \emph{Still open:} the exhaustiveness of the parity
dichotomy, i.e.\ whether every data-valued lifting of Reader/State is $\prod$, $\coprod$, or a mixture.}
\end{teachbox}

Be honest about the reach. The fibration \emph{organizes} and \emph{characterizes}; it does not
\emph{eliminate} the one hard axiom. As noted, three of the four entwining axioms fall out for free from
$G$ being a fibred comonad. But the fourth, the multiplication axiom E2 of $T$, touches $T$'s fibre
multiplication $j$ --- the Mendler restriction of Ahman--Bauer's $\star$ --- and that datum is \emph{extra}
structure living in the $\star$-algebra, not forced by $p$. Cartesianness of $T$ does not by itself hand
over E2; the general symbolic chase stays where Remark~\ref{rem:entwine-scope} left it. The fibration is a
lens that makes the \emph{shape} of the obstruction legible; it is not a machine that discharges it.
\prov{The four-level stratification, the refutation of the collapse slogan, and the splitters
$\mathsf{Maybe}$ (strict BC $\subsetneq$ non-branching) and $\mathsf{List}$ (non-branching $\subsetneq$
cartesian) are MacBeth, proved (\texttt{proofs/2026-08-05-cartesian-preservation-nonbranching.md}). The
non-branching $\Leftrightarrow$ reverse-$\kappa$ biconditional is Theorem~\ref{thm:arrows}; $G_{M}$
cartesian for all $M$ is Lean-certified (\texttt{FibredTransfer.lean}). Fibrational vocabulary cited:
von Glehn \cite{vonGlehn18}, Jacobs \cite{Jacobs99}, Street \cite{Street72}.}

\begin{Remark}[A word of caution: comonadic because contravariant]\label{rem:contravariant-caution}
It is tempting to read the transfer as saying something deep about \emph{comonadic} computation --- that
$M$'s effects were secretly \emph{coeffects} all along. The fibrational picture counsels restraint, and
Plotkin and Turi's standing suspicion of too-easy monad/comonad dualities is well taken here. The comonad
structure of $G$ is nothing more than the \emph{contravariant image} of $M$'s monad structure: unit becomes
counit, multiplication becomes comultiplication, associativity becomes coassociativity, each by the single
fibrewise $(-)\op$ and by nothing else. $G$ ``looks comonadic'' exactly \emph{because} positions are
contravariant --- a structural fact about \emph{where} in the fibration one applies $M$, not a claim that
$M$ was a comonad in disguise or that the transfer captures an operational duality of effects and contexts.
The real content is elsewhere: that \emph{both} liftings exist, from the one $M$, and that they
\emph{entwine}. To present the comonad side as an operational-semantics discovery would be to oversell
contravariance. We take the transfer for the clean structural fact it is.
\end{Remark}

\subsection{Which monads lift: the \texorpdfstring{$\Sigma$}{Sigma}-lifting is \texorpdfstring{$M\triangleleft(-)$}{M comp -}}
\label{sec:sigma-is-comp}

The grading table of the last box did two things and left two undone. It \emph{located} the
surviving codomain lift --- the $\coprod$-lift $T^{\Sigma}_{M}$, data at one chosen leaf --- and it
\emph{exhibited} that lift as a genuine monad on $\Cont$ for Reader and State. What it
did not say is \emph{which} base monads $M$ carry the $\Sigma$-lift as a monad, nor \emph{why} the
``canonical section'' that threads the multiplication should cohere. Both were flagged open. Both are
answered by a single observation, and it is not a coincidence of sections but an identity of
endofunctors: \emph{the $\Sigma$-lifting is left composition by $M$}. Once you see it, the entire
08-07 verification for Reader and State collapses into an instance of a theorem you already proved in
Chapter~\ref{ch:comonoids}.

\begin{Proposition}[$T^{\Sigma}_{M}=M\triangleleft(-)$]\label{prop:sigma-is-comp}
Let $M=\cont{S_{M}}{P_{M}}$ be a container, carrying candidate morphisms $\eta^{M}\colon\y\to M$ and
$\mu^{M}\colon M\triangleleft M\to M$ (not yet assumed to satisfy any law). For every container
$C=\cont S P$ there is an equality of containers
\[
  T^{\Sigma}_{M}\,C\;=\;M\triangleleft C,
\]
natural in $C$; and under it the unit and multiplication of the $\Sigma$-lifting are the whiskerings
$\eta^{\Sigma}=\eta^{M}\triangleleft(-)$ and $\mu^{\Sigma}=\mu^{M}\triangleleft(-)$ of $M$'s own
container-unit and container-multiplication.
\prov{MacBeth, proved (\texttt{proofs/2026-08-08-sigma-monad-is-triangle-monoid.md}, Prop.~2.1) and
Lean-verified: \texttt{lean/Containers/Containers/SigmaLift.lean}, where
\texttt{sigmaLift\_eq\_seq} closes $M.\mathrm{sigmaLift}\,C=M\triangleleft C$ by \texttt{rfl}
(a \emph{definitional} equality; sorry-free, axiom-free).}
\end{Proposition}

The proof is a match of shapes and positions, and it is worth doing in the open because it explains
the definition of $T^{\Sigma}$ rather than merely confirming it. Write an element $m\in M X$ as a
shape $s\in S_{M}$ together with a labelling $x\colon P_{M}(s)\to X$, so its leaves are exactly
$\mathrm{lv}(m)=P_{M}(s)$ and the label at $b$ is $x_{b}=x(b)$. Now unfold the composite:
\[
  \Ext{M\triangleleft C}X=\Ext{M}\bigl(\Ext{C}X\bigr)
  =\coprod_{s\in S_{M}}\Bigl(\coprod_{s'\in S}X^{P(s')}\Bigr)^{P_{M}(s)}.
\]
Its \emph{shapes} are $\coprod_{s}S^{P_{M}(s)}=M(S)$ --- an $M$-structure whose leaves are labelled by
shapes of $C$, which is exactly the shape set of $T^{\Sigma}_{M}\cont S P$. Its \emph{positions} over
such a shape $m$ are $\coprod_{b\in P_{M}(s)}P(x_{b})=\coprod_{b\in\mathrm{lv}(m)}P(x_{b})$ --- pick an
$M$-leaf, then a $C$-position beneath the shape sitting there --- which is precisely $P^{\Sigma}(m)$.
The unit's backward map (the codiagonal fold) is $\eta^{M}$ read on positions; the multiplication's
backward map --- the ``canonical section'' $\sigma$ of the 08-07 proof --- is nothing but the backward
position map of $M$'s own multiplication $\mu^{M}$. There was never a section to construct: it was
$\mu^{M}$ all along.

This turns the mysterious coherence into a triviality and, better, into a \emph{characterisation}.

\begin{Corollary}[$\Sigma$-monad $\Leftrightarrow$ container monad]\label{cor:sigma-monad}
For a container $M$ with candidate $\eta^{M},\mu^{M}$ as in Proposition~\ref{prop:sigma-is-comp}, the
following are equivalent.
\begin{enumerate}[nosep,label=\textup{(\arabic*)}]
  \item $\bigl(T^{\Sigma}_{M},\eta^{\Sigma},\mu^{\Sigma}\bigr)$ is a monad on $\Cont$.
  \item $\bigl(M,\eta^{M},\mu^{M}\bigr)$ is a $\triangleleft$-monoid in $(\Cont,\triangleleft,\y)$ ---
        a \emph{container monad} (Theorem~\ref{thm:comp-comonoid-monoid}).
  \item $\bigl(\Ext M,\Ext{\eta^{M}},\Ext{\mu^{M}}\bigr)$ is a monad on $\Set$ whose functor
        \emph{and} whose unit and multiplication are polynomial.
\end{enumerate}
Under any of them the three coherence conditions of the 08-07 proof hold \emph{automatically}: they
are the backward halves of the $\triangleleft$-monoid's unit and associativity laws, read on positions.
\prov{MacBeth, proved (\texttt{2026-08-08-sigma-monad-is-triangle-monoid.md}, Thm.~3.1); the three
monad laws are Lean-verified in \texttt{SigmaLift.lean} as the whiskered $\triangleleft$-monoid laws
(each closes by the monoid law plus definitional coherence).}
\end{Corollary}

\begin{proof}[Why it holds]
Every monoid $A$ in a monoidal category $(\mathcal V,\otimes,I)$ makes $A\otimes(-)$ a monad, with
structure maps $m\otimes(-)$ and $e\otimes(-)$; conversely, evaluating the monad laws of $A\otimes(-)$
at the unit object $I$ returns the monoid laws of $A$. Take $(\mathcal V,\otimes,I)=(\Cont,\triangleleft,\y)$
and $A=M$: by Proposition~\ref{prop:sigma-is-comp} that monad \emph{is} $T^{\Sigma}_{M}$, so
(1)\,$\Leftrightarrow$\,(2). And (2)\,$\Leftrightarrow$\,(3) because the extension
$\Cont\simeq\Poly\hookrightarrow[\Set,\Set]$ is fully faithful and strong monoidal
($\triangleleft\mapsto\circ$, $\y\mapsto\mathrm{Id}$; Niu--Spivak \cite{NiuSpivak23}), and a fully
faithful strong monoidal functor reflects and creates monoids.
\end{proof}

So the surviving half of the proof-relevance boundary does not live in some exotic corner: it lands
\emph{exactly} on the composition-monoid spine of Chapter~\ref{ch:comonoids}. Directed containers are
the $\triangleleft$-\emph{comonoids} (Theorem~\ref{thm:objdict}); container monads are the
$\triangleleft$-\emph{monoids}, the dual partner; and the $\Sigma$-lift is a monad on $\Cont$ precisely
for the latter. The running witnesses are two of the smallest such monoids. Reader $y^{E}=\cont 1 E$
is the container monad whose $\triangleleft$-monoid structure is the \emph{unique comonoid on the set
$E$} --- the diagonal $E\to E\times E$ --- so its section $\sigma$ is $e\mapsto(e,e)$: the diagonal of
the 08-07 proof was $\mu^{M}$'s backward map. State is the store container monad, whose $\sigma$ is the
threading map. The entire 08-07 computation was, in retrospect, a hand-verification that Reader and
State are $\triangleleft$-monoids --- which Chapter~\ref{ch:comonoids} had already established.

\paragraph{The counterexample that names the real content.} The open flag asked whether
\emph{reverse-total} --- the pointwise existence of a label-preserving section of every $\mu^{M}$ ---
is enough on its own to make $T^{\Sigma}_{M}$ a monad. Corollary~\ref{cor:sigma-monad} says the honest
requirement is stronger: $M$ must be a container monad, a \emph{polynomial} fact, not a pointwise one.
A single monad realises the gap.

\begin{Theorem}[$\mathsf{Bag}$ refutes ``reverse-total $\Rightarrow$ $\Sigma$-monad'']\label{thm:bag-refutes}
The finite-multiset monad $\mathsf{Bag}$ (with $\eta\,x=\{x\}$ and $\mu=\uplus$) is reverse-total in
the strongest possible way --- $\mu$ is a leaf-\emph{bijection}, so a section exists with $\sigma=\id$
--- yet $T^{\Sigma}_{\mathsf{Bag}}$ is \emph{not a monad on $\Cont$}; it is not even a functor on
$\Cont$, because $\mathsf{Bag}$ is not a container.
\prov{MacBeth, proved (\texttt{2026-08-08-sigma-monad-is-triangle-monoid.md}, \S4; cardinality witness
in \texttt{scratch/sigma-monad-coherence/bag\_not\_container.py}).}
\end{Theorem}

\begin{proof}[Why it holds]
$\mathsf{Bag}$ fails to preserve the connected pullback $A\to 1\leftarrow B$, which every container
(equivalently, every polynomial functor) preserves. Concretely, at $A=B=\{0,1\}$ the two size-two
multisets $\{(0,0),(1,1)\}$ and $\{(0,1),(1,0)\}$ over $A\times B$ have the \emph{same} image under
$(\pi_{A},\pi_{B})$, namely $(\{0,1\},\{0,1\})$: the multiset has forgotten the pairing. (Counting by
total size confirms it: $\lvert\mathsf{Bag}(A\times B)\rvert_{2}=10\neq 9=\lvert\mathsf{Bag}\,A\times_{\mathsf{Bag}\,1}\mathsf{Bag}\,B\rvert_{2}$.)
So $\mathsf{Bag}\notin\Cont$. But $T^{\Sigma}_{\mathsf{Bag}}$ being a functor on $\Cont$ would supply,
for each relabelling $f$, a natural backward leaf-map --- exactly a container structure on
$\mathsf{Bag}$. None exists, because a relabelling that identifies two labels forces a
choice of \emph{which} symmetric token survives, and no such choice is natural.
\end{proof}

The discriminator is the oldest one in the book: \emph{polynomial}, not merely \emph{analytic}.
$\mathsf{Bag}$ is analytic --- built from symmetric powers, carrying nontrivial $S_{n}$ symmetries ---
and it is those symmetries that block a natural leaf-assignment. Reverse-total is the pointwise,
analytic shadow (``some backward map exists at each $\mu^{M}$''); the container-monad condition is the
coherent, polynomial fact (``the backward map is natural in the labels and unital-associative'').
Lining the two up:
\[
  \text{reverse-total} : \triangleleft\text{-monoid}
  \;::\;
  \text{analytic} : \text{polynomial}
  \;::\;
  \text{forgets provenance} : \text{tracks provenance}.
\]
$\mathsf{Bag}$ forgets which input each output came from; that is the precise sense in which it loses
provenance, and here that loss costs it the functoriality of $T^{\Sigma}$ itself.

\begin{teachbox}[Both legs of the codomain fibration, with clean criteria]
The codomain fibration $p\colon\Cont\to\Set$ carries \emph{two} data-valued leaf-liftings of a base
monad $M$, and each now has a crisp lift-criterion --- one on each side of the $\mathbb Z/2$ grading:
\[
\renewcommand{\arraystretch}{1.4}
\begin{array}{c|c|c}
 \text{leaf-lifting} & \text{is a monad on }\Cont\iff & \text{totality direction}\\\hline
 \prod\ (=T_{M},\ \text{Ahman--Bauer \cite{AhmanBauer24}}) & M\ \textbf{cartesian} & \text{forward}\\
 \coprod\ (=T^{\Sigma}_{M}=M\triangleleft-) & M\ \text{a }\textbf{$\triangleleft$-monoid} & \text{reverse}
\end{array}
\]
The two criteria are independent axes, not nested: cartesianness is about $\mu^{M}$ \emph{merging}
leaves, the $\triangleleft$-monoid condition about the polynomial coherence of $\mu^{M}$'s backward
map. $\mathsf{List}$ satisfies both; Reader and State are $\triangleleft$-monoids but non-cartesian, so
they carry the $\coprod$-lift and not the $\prod$-lift --- exactly the split the grading table
predicted. This is the payoff Neil's fibrational framing was pointing at: \emph{both} legs of the
non-thin (codomain) side of the boundary, which the classical predicate-lifting literature ---
confined to thin, faithful fibrations (Hermida--Jacobs \cite{HermidaJacobs98}) --- never reaches, now
have monad-lift criteria as clean as the cartesian one.

\smallskip
\emph{One caveat, at the boundary with a base monad's own join.} The identity
$T^{\Sigma}_{M}=M\triangleleft-$ and Corollary~\ref{cor:sigma-monad} are strict, object-level, and
base-monad-free: pure $\triangleleft$-algebra. This is a \emph{different} layer from Orestis's oplax
comparison $\Lambda P\circ\mathrm{join}\subseteq\Lambda(\Lambda P)$ \cite{OrestisAgda}, which asks how
the $\forall$-lift interacts with a \emph{base} monad's join and is an $\subseteq$, not an equality.
The pure composition law is strict here; it is once the base multiplication re-enters --- the
$T_{M}$-boundary of the stratification box --- that it degrades to Orestis's oplax $\subseteq$.
\end{teachbox}

\subsection{Which liftings, all of them: monad liftings are comonads over the base}
\label{sec:liftings-are-categories}

Corollary~\ref{cor:sigma-monad} said \emph{which} base monads carry the $\coprod$-lift, and the
grading table said the $\prod$-lift of Reader dies. But both are statements about \emph{two named
liftings}. A more basic question has gone unasked. Fix Reader as the base monad and let the
aggregator range freely: what is the \emph{complete} list of its data-valued monad liftings? Is every
lifting a $\coprod$, a $\prod$, or --- the natural guess --- some leafwise blend of the two? The
answer is none of those tidy shapes, and it is the most satisfying one this book could offer. The
liftings are \emph{small categories} --- the very objects Chapter~\ref{ch:comonoids} taught us
directed containers were (Theorem~\ref{thm:objdict}), reappearing one categorical level up, and for
exactly the same reason.

Recall the Reader container $R=\y^{E}=\cont 1 E$: one shape, $E$ positions, $R(S)=S^{E}$. Along the
shape fibration $p\colon\Cont\to\Set$ a fibred endofunctor over $R$ cannot see the shapes at all ---
they are fixed to be $S^{E}$ --- so all its content is a single recipe that turns the
positions-over-the-leaves into new positions. That recipe is an \emph{aggregator}.

\begin{Proposition}[Reduction to an aggregator]\label{prop:reader-reduction}
Fibred endofunctors of $\Cont$ lying over $R=\y^{E}$ correspond bijectively (up to natural iso) to
functors $L\colon\Set^{E}\to\Set$, via
\[
  T\cont S P=\bigl(S^{E},\ m\mapsto L(P\circ m)\bigr),\qquad m\colon E\to S \text{ a Reader-shape.}
\]
A monad structure on such a $T$ lying over $(R,\eta,\mu)$ is exactly a pair of natural
transformations
\[
   \varepsilon_{A}\colon L(\langle A\rangle)\to A,
   \qquad
   \delta_{D}\colon L\bigl(b\mapsto D(b,b)\bigr)\longrightarrow L\bigl(b\mapsto L(c\mapsto D(b,c))\bigr),
\]
($\langle A\rangle$ the constant family at $A$; $D\colon E\times E\to\Set$) subject to three laws.
\prov{MacBeth, proved (\texttt{proofs/2026-08-09-reader-liftings-are-categories.md}, Prop.~2.1). The
functor part uses that the reindexing arrow $(m,\id)$ is cartesian, so a fibred $T$ preserves it; the
monad part reads off $\varepsilon$ from $\eta^{T}$ over $\eta$ and $\delta$ from $\mu^{T}$ over
$\mu$.}
\end{Proposition}

Now look hard at the shape of that data. The unit $\varepsilon\colon L\circ\Delta\Rightarrow\mathrm{Id}$
is a \emph{counit}; the multiplication $\delta$ is a \emph{comultiplication}; and the three monad laws
become a counit law and coassociativity. The lifting is a \emph{monad} on $\Cont$, yet the aggregator
that defines it carries a \emph{comonad}. This is the one astonishment of the section, and it is not
an accident. The multiplication $\mu^{T}\colon TT\to T$ is a \emph{container morphism}, so its only
interesting half is a \emph{backward} position map --- and a backward multiplication is a
comultiplication. It is precisely the one-operation-two-faces phenomenon of the transfer,
\S\ref{sec:moncomon-transfer}: the fibrewise op turns a monad multiplication upstairs into a comonad
comultiplication downstairs. Monad liftings are comonads --- and comonads, we already know, are
categories.

Before the classification, one clean lemma explains \emph{why} $\coprod$ was the survivor and
$\prod$ was not, in a single line that recovers the cartesian boundary of the whole chapter.

\begin{Lemma}[$\prod$ needs a monoid, $\coprod$ needs nothing]\label{lem:monoid-comonoid}
Let $L$ be an aggregator built from the leaves of $E$.
\begin{enumerate}[nosep,label=\textup{(\arabic*)}]
  \item $\coprod$ --- coproduct over any index --- carries a canonical comonad structure, hence is a
        monad lifting of $R$, for \emph{every} index. It is the free comonoid on the leaves.
  \item $\prod$ over an index $I$ carries a comonad structure, hence lifts, \emph{iff} $I$ carries a
        \emph{monoid}.
\end{enumerate}
\prov{MacBeth, proved (\texttt{proofs/2026-08-09-lifting-dichotomy-exhaustiveness.md}, Lemma~U, both
sides; the $\prod$ side is the standard fact that $Q\mapsto Q^{I}$ is a comonad on $\Set$ iff $I$ is a
monoid, pulled back along the monad morphism $\mathrm{ev}\colon R\Rightarrow\mathrm{Id}$).}
\end{Lemma}

Reader's leaf-set $E$ carries the \emph{diagonal} comonoid $E\to E\times E$ --- the unique comonoid on
any set --- but generally no monoid. So $\coprod$ lifts and full $\prod$ over $E$ does not: a
cross-leaf product would demand a multiplication $E\times E\to E$, and Reader supplies none. This is
$T_{M}$-lifts-iff-$M$-cartesian (\S\ref{sec:moncomon-fibration}) read one categorical level down: a
cartesian $\mu^{M}$ is a bijective, monoid-like law on positions; a non-cartesian one drops leaves,
leaves no monoid, and $\prod$ has nowhere to multiply. An \emph{auxiliary} monoid $M$ supplies its own
multiplication, so the single-leaf power $\prod_{M}=(-)^{M}$ does lift --- and that single-leaf power,
we are about to see, is a one-object category.

\begin{Theorem}[Reader's liftings are small categories]\label{thm:reader-classification}
Assume the aggregator $L\colon\Set^{E}\to\Set$ is polynomial. Then the data-valued monad liftings of
Reader $R=\y^{E}$ are in bijection with $E$-indexed families of small categories
$(\mathcal C_{v})_{v\in E}$, under
\[
  L(B)=\coprod_{v\in E}\ \coprod_{i\in\ob\mathcal C_{v}}
        B_{v}^{\ \coprod_{j}\mathcal C_{v}(i,j)},
\]
with counit $\varepsilon$ the identity morphisms and comultiplication $\delta$ the composition.
\prov{MacBeth, proved (\texttt{2026-08-09-reader-liftings-are-categories.md}, \S3, Steps~A--E),
computationally corroborated by two independent engines (\texttt{enum.py} vs \texttt{catcount.py},
agreeing on \emph{$\#$liftings}\,$=$\,\emph{$\#$categories} on every tested out-degree profile). The
final step is the equivalence \emph{polynomial comonads $\cong$ small categories} ---
Theorem~\ref{thm:objdict}, from Ahman--Uustalu \cite{AU16} --- read here one categorical level up.}
\end{Theorem}

\begin{proof}[Why it holds]
By Proposition~\ref{prop:reader-reduction} a lifting is an aggregator $L$ with a counit $\varepsilon$
and coassociative comultiplication $\delta$. Writing $L$ in polynomial normal form
$L(B)=\coprod_{s}\prod_{v}B_{v}^{P_{L}(s,v)}$, the counit forces every shape to have a
\emph{distinguished} position (a marked identity), and the left- and right-unit laws force each shape
to read \emph{one} leaf only --- an off-diagonal, cross-leaf arity is annihilated by the empty
diagonal (this is where $\prod$, whose only shape reads two leaves, is excluded outright: it admits no
$\delta$). What remains, per leaf $v$, is a polynomial comonad on $\Set$: objects the shapes at $v$,
morphisms out of a shape its positions, identities the marked positions, composition the backward map
of $\delta$. By the equivalence of Theorem~\ref{thm:objdict} that is a small category. The
correspondence is mutually inverse.
\end{proof}

The clean liftings decode as the degenerate categories, and the exotic one becomes visible.

\begin{center}
\renewcommand{\arraystretch}{1.35}
\begin{tabular}{@{}l l c@{}}
 aggregator $L$ & category data & lifts?\\\hline
 $\prod_{v}B_{v}$ \ (Ahman--Bauer $T_{R}$) & one object reading \emph{all} leaves & \textbf{No}\\
 $\coprod_{v\in U}B_{v}$ \ ($\Sigma_{U}$) & \emph{discrete} category on $U$ & Yes\\
 $\coprod_{v}W_{v}\cdot B_{v}$ \ (weighted) & discrete category, objects $\textstyle\coprod W_{v}$ & Yes\\
 $B_{0}\times B_{0}$, swap $\delta$ & one-object group $\mathbb Z/2$ (non-$\Sigma$, non-$\prod$) & Yes\\
 $B_{0}\times B_{1}$ \ (cross-leaf shape) & a cross-leaf object & \textbf{No}
\end{tabular}
\end{center}

\noindent So the survivors are categories; $\Sigma_{U}$ are exactly the \emph{discrete} ones; and the
imagined ``leafwise mix of $\prod$ and $\coprod$'' was a mirage. The honest parameter is a small
category per leaf: a shape may read one leaf with multiplicity $\ge 2$ --- these are the hom-sets of a
\emph{non}-discrete category, the products that survive because they live \emph{within} a single leaf
--- but never across leaves, which is exactly the $\prod$-obstruction of
Lemma~\ref{lem:monoid-comonoid}.

\begin{Remark}[Polynomial is the boundary --- again]\label{rem:analytic-excluded}
The theorem assumes $L$ polynomial. The counit alone rules out the natural non-polynomial rivals: a
natural $\varepsilon\colon L(\langle-\rangle)\Rightarrow\mathrm{Id}$ must pick a \emph{canonical}
element, and a symmetric summand --- an unordered pair, the symmetric square, the size-two multiset
$\mathsf{Bag}_{2}$ --- has no swap-invariant canonical element, so it admits no $\varepsilon$ and is
not a lifting. It is the same discriminator as Theorem~\ref{thm:bag-refutes}, on the same functor:
\emph{polynomial}, not merely \emph{analytic}. (A full removal of the polynomiality hypothesis would
need the lemma ``every accessible $\Set$-comonad with a counit is polynomial'', which is not proved
here; the natural rivals are excluded, the general statement is flagged open.)
\prov{MacBeth, proved (\texttt{2026-08-09-reader-liftings-are-categories.md}, Prop.~5.1), verified in
\texttt{analytic.py}: $\mathrm{Sym}^{2}$ and $\mathsf{Bag}_{2}$ admit no natural $\varepsilon$, the
ordered square does.}
\end{Remark}

\begin{teachbox}[The whole boundary lands on $\Cat$]
Stand back. Chapter~\ref{ch:comonoids} proved that the $\triangleleft$-comonoids --- the directed
containers --- are the small categories (Theorem~\ref{thm:objdict}), the promised
$\textbf{Containers}\supseteq\textbf{Directed Containers}\cong\textbf{Small Categories}$. This section
proves that the proof-relevant monad \emph{liftings} of Reader are \emph{also} small categories, and
by the identical mechanism: a monad lifting is secretly a comonad (its multiplication is a backward
map), and comonads are categories. The two facts are the same equivalence, Theorem~\ref{thm:objdict},
read at two different heights. The directed-container spine of the book is not one theorem but a gravitational
well: the objects, their morphisms (the cofunctors of $\DCont\cong\Cof$), and now their liftings all
fall into $\Cat$. Neil's two feeds and the $\mathbb Z/2$ grading were the map; the destination is the category
of small categories.
\end{teachbox}

Reader was the base case $S_{M}=1$: one shape, one aggregator, one small category per leaf. What
happens at the opposite extreme --- the \emph{richest} finite store there is?

\subsection{The State pole: the store multiplication is invisible}
\label{sec:state-liftings}

Take $M=\text{State}$, $MX=(S\times X)^{S}$. Its shape set is $S^{S}$, the full transformation monoid
of the state set --- every possible way to overwrite $S$, composed under function composition. If any
base monad were going to force its liftings into some elaborate graded shape, it is this one: the
reduction of Proposition~\ref{prop:reader-reduction} now hands back not a single aggregator but a
whole \emph{family} $(A_{\sigma}\colon\Set^{S}\to\Set)_{\sigma\in S^{S}}$, one per store-shape, and
the multiplication of State \emph{threads} them together along the monoid $S^{S}$, coupling the source
states through the rule $\sigma_{\mu}(s)=t_{s}(T(s))$. Every instinct says the answer is an
``$S^{S}$-graded category'': a category that transforms as you push it around the store.

Here is the punchline, and it is the one this section was built to earn.

\begin{Theorem}[State's liftings are small categories]\label{thm:state-classification}
The polynomial data-valued monad liftings of $\text{State}=\cont{S^{S}}{S}$ are in bijection with
small categories $\mathcal C$, under $\mathcal C\mapsto\mathbb S\times\mathcal C$, where $\mathbb S$ is
the action category of $S^{S}\curvearrowright S$. Explicitly the aggregator is
\[
   A_{t}(Q)\ \cong\ \coprod_{s\in S}\ \coprod_{c\in\ob\mathcal C}\
        Q_{s}^{\,\coprod_{c'\in\ob\mathcal C}\mathcal C(c,c')}
   \qquad\text{(independent of the grade $t\in S^{S}$),}
\]
with $\varepsilon$ the $\mathcal C$-identities and $\delta$ the source-routed $\mathcal C$-composition
carrying \emph{trivial} transport. The whole $S^{S}$-grading is invisible: State's liftings are the
\emph{same} $\Cat$ as a single-object Reader's.
\prov{MacBeth, proved at the object level; the morphism level by the mirror argument and exhaustive
machine check at $|S|=2$ (\texttt{2026-08-11-state-liftings-holonomy-triviality.md}). Soundness
$\Cat\hookrightarrow$ liftings, $\mathcal C\mapsto\mathbb S\times\mathcal C$, is Lean-verified
(\texttt{StateProductLifting.lean}).}
\end{Theorem}

\begin{proof}[Why it holds --- the store multiplication is a mirage]
Two things could make State richer than Reader, and both collapse.

\emph{First, the grade could matter.} A priori the aggregator $A_{t}$ depends on the store-shape
$t\in S^{S}$. It does not. The left-unit law hands each grade a comparison map
$\mathrm{sh}_{t}\colon A_{t}\to A_{\id}$, and a factorisation whose threading returns to $\id$ hands
back a map $\mathrm{pr}_{t}$ the other way; the \emph{outermost-object component of associativity}
forces these two to be mutually inverse. So $A_{t}\cong A_{\id}$ for every grade --- this is
\emph{grade-independence} --- and $A_{\id}$, by the Reader theorem with $E=S$ (at grade $\id$ the
threading is trivial, $s\mapsto\id(s)=s$, so the store does not couple the states), is an $S$-indexed
family of small categories $(\widetilde{\mathcal C}_{s})_{s\in S}$.

\emph{Second, the coupling could carry holonomy.} The only freedom left is how $\delta$ transports the
fibre $\widetilde{\mathcal C}_{s}$ along a store-arrow $s\to g(s)$. Read the \emph{deepest} object of
the associativity law. Writing $\tau^{g}(s,c)$ for the transport of an object
$c\in\widetilde{\mathcal C}_{s}$ along $g$, associativity says
\[
   \tau^{g'}(s,c)\ =\ \tau^{t_{s}}\!\bigl(T(s),\,\tau^{T}(s,c)\bigr),
   \qquad g'(s')=t_{s'}(T(s')).
\]
Look at the asymmetry. The right-hand side sees only the middle map \emph{at the source state} $s$;
the left-hand side depends on the whole function $g'$, stitched from \emph{all} the middles. The only
way the two can agree for every choice of the off-$s$ middles is if the transport cannot see them
either --- that is, if $\tau^{g}(s,c)$ depends on $g$ \emph{only through the endpoint} $g(s)$. This is
\emph{endpoint-locality}, and it is strictly stronger than functoriality.

Endpoint-locality is exactly the statement that the transport is a functor out of the
\emph{codiscrete} category on $S$ --- the category with a unique arrow $s\to m$ for every ordered pair.
A functor out of a codiscrete category is a trivial coherent system of isomorphisms: choose a
basepoint $0\in S$, identify every fibre with $\widetilde{\mathcal C}_{0}=:\mathcal C$ along the unique
arrows, and every transport becomes the identity. The $S$ source-fibres fuse into one category
$\mathcal C$, and the store transport is trivial. What is left is
$A_{t}(Q)=\coprod_{s}\coprod_{c}Q_{s}^{\,\coprod_{c'}\mathcal C(c,c')}$ with threaded
$\mathcal C$-composition ---
precisely $\mathbb S\times\mathcal C$.
\end{proof}

\begin{Remark}[Why the graded category was a mirage]
The tempting rivals all satisfy functoriality but fail endpoint-locality, and endpoint-locality is
the sharper blade. A copresheaf transport $F\colon\mathbb S\to\Set$ --- say the representable
$\mathbb S(0,-)$ --- is a perfectly good functor, obeying $\tau^{b\circ a}=\tau^{b}\circ\tau^{a}$; yet
$F(g)$ reads the whole function $g$, not just $g(s)$, so it fails the deepest-object identity and
breaks associativity. Per-state-different fibres $\widetilde{\mathcal C}_{0}\not\cong
\widetilde{\mathcal C}_{1}$ are impossible too: $S^{S}$ acts transitively on $S$, so the codiscrete
category is connected and all fibres are forced isomorphic. The finer graded object one hoped for
simply does not exist; the honest answer is the \emph{coarser} $\Cat$.
\prov{MacBeth, proved/computed (\texttt{2026-08-10-state-liftings-holonomy-free.md} §5;
\texttt{2026-08-11-state-liftings-holonomy-triviality.md} §5): representables satisfy the unit laws
but fail associativity, and every non-trivial automorphism-twist of $\mathbb S\times\mathcal C$ fails
the laws.}
\end{Remark}

\begin{teachbox}[One theorem, two poles: the invariant is $\pi_{0}$ of the threading action]
Reader and State are now the same theorem, read at two settings of a single dial. Every polynomial
monad lifting of a container monad is governed by the \emph{position-threading action} --- the way the
multiplication carries a position to its endpoint --- and specifically by its \emph{action category}
$\mathcal A$ and the number of connected components $\pi_{0}(\mathcal A)$.
\begin{center}
\renewcommand{\arraystretch}{1.35}
\begin{tabular}{@{}l c c l@{}}
 base monad & action category $\mathcal A$ & $\pi_{0}(\mathcal A)$ & liftings \\\hline
 Reader $\y^{E}$ & \emph{discrete} on $E$ & $|E|$ & $E$-indexed families of small categories\\
 State $(S\times-)^{S}$ & \emph{codiscrete} on $S$ (collapsed) & $1$ & a single small category $\mathcal C$\\
\end{tabular}
\end{center}
Reader spreads its liftings across $|E|$ independent components --- one category per leaf, no transport
between them. State collapses to a single component --- one category, and the store multiplication that
looked so rich contributes \emph{nothing} to the classification. The dial is $\pi_{0}$ of the
transport, and the two monads sit at its extremes: maximally disconnected, and maximally connected.
So the whole classification would seem to be a \emph{count} --- read off $\pi_{0}$ of the threading
action and you are done. It is irresistible to conjecture exactly that. Resist it: the next subsection
shows the count is an illusion, and the illusion is the most instructive thing in the chapter.
\end{teachbox}

So $\pi_{0}$ organises the two poles. Does it \emph{classify}? Here the story turns --- and it is the
turn the whole chapter has been walking toward.

\subsection{The general law: holonomy, and the two poles as its degeneracies}
\label{sec:update-liftings-holonomy}

Fair is foul, and foul is fair. The two clean answers --- Reader's family of categories, State's single
category --- looked like the whole truth, and the $\pi_{0}$ dial looked like the law that governed
them. Both were degeneracies. There is a wider class of container monads on which the classification is
\emph{not} a count, on which the transport carries an invariant $\pi_{0}$ cannot see, and on which the
grading we twice declared a mirage is alive and irreducible. Reader and State are holonomy-free not
because monad liftings are holonomy-free, but because each of those two monads kills the holonomy in
its own particular way. Name the wider class, and both poles fall out as the two ways of being trivial.

\paragraph{The class that reveals it.} Between ``one shape'' (Reader) and ``the full transformation
monoid'' (State) sits the family that interpolates them: the \emph{update monads} of Ahman and Uustalu
\cite{AhmanUustalu13}. Fix a set $S$ of states and a \emph{monoid} $P$ of updates acting on $S$ on the
right, $\downarrow\colon S\times P\to S$, with $s{\downarrow}o=s$ and
$s{\downarrow}(p\oplus q)=(s{\downarrow}p){\downarrow}q$. The update monad is the polynomial functor
\[
   \mathrm{Upd}_{(S,P,\downarrow)}(X)\ =\ \sum_{f\in P^{S}} X^{S},
   \qquad
   \mu(F)(s)=\bigl(f(s)\oplus g_{s}(s{\downarrow}f(s)),\ x_{s}(s{\downarrow}f(s))\bigr),
\]
its shapes the update-fields $f\in P^{S}$, its positions the states $S$. The cardinality of the
position set never changes --- but the multiplication \emph{threads} a state $s$ to its
\emph{endpoint} $s{\downarrow}f(s)$, and that is where the positions genuinely move. Reader is the case
$P=1$ (nothing to apply, $s{\downarrow}o=s$); State is the case $P=$ the overwrite monoid
$\{o\}\cup\{w_{m}:m\in S\}$ with $s{\downarrow}w_{m}=m$. The threading is no longer a fixed background
fact; it is the datum $\downarrow$, and it varies from one update monad to the next.

\begin{Definition}[The position-threading action and its category]\label{def:action-category}
The \emph{position-threading action} of $\mathrm{Upd}_{(S,P,\downarrow)}$ is the monoid action
$\downarrow\colon P\curvearrowright S$ itself. Its \emph{action category} $\mathbb A(\downarrow)$ (the
translation category of the action) has
\[
  \text{objects } S,\qquad
  \text{an arrow } s\xrightarrow{\ p\ } s{\downarrow}p \text{ for each } p\in P,\qquad
  \text{composition } (s\xrightarrow{p}\,\xrightarrow{q}) = (s\xrightarrow{p\oplus q}),\quad
  \id_{s}=(s\xrightarrow{o}s).
\]
Its connected components are the orbits, $\pi_{0}(\mathbb A(\downarrow))=S/{\downarrow}$; its
\emph{isotropy monoid} at $s$ is the stabiliser $\mathrm{Stab}(s)=\{p\in P:s{\downarrow}p=s\}$.
For Reader, $\mathbb A(\downarrow)$ is the \emph{discrete} category on $S$; for State it is the
overwrite category, which we shall see collapses.
\end{Definition}

Now the reduction runs exactly as it did for the two poles --- grade-independence trims the family of
aggregators to one, the Reader argument makes each source-fibre a small category $\mathcal C_{s}$ --- and
the only freedom that remains is the transport: how the comultiplication carries an object
$c\in\mathcal C_{s}$ along an update $p$ to a new object in $\mathcal C_{s{\downarrow}p}$. Write it
$\rho_{s,p}\colon\mathcal C_{s}\to\mathcal C_{s{\downarrow}p}$. The deepest object of the associativity
law --- the same instance that forced grade-independence and endpoint-locality before --- now reads, for
all $s\in S$ and $p,q\in P$,
\[
   \rho_{s,\,p\oplus q}\ =\ \rho_{s{\downarrow}p,\,q}\circ\rho_{s,p},
   \qquad \rho_{s,o}=\id_{\mathcal C_{s}}.
\]
Read the display again: it is \emph{precisely} the statement that $\rho$ is functorial on
$\mathbb A(\downarrow)$. The transport is not extra structure bolted onto the classification --- it
\emph{is} a functor out of the action category, and associativity is its functoriality.

\begin{Theorem}[Update-monad liftings are functors on the action category]\label{thm:update-classification}
The degree-$1$ proof-relevant polynomial monad liftings of $\mathrm{Upd}_{(S,P,\downarrow)}$ are in
bijection with functors
\[
   \mathbb A(\downarrow)\longrightarrow\Cat,\qquad s\mapsto\mathcal C_{s},\quad
   (s\xrightarrow{p}s{\downarrow}p)\mapsto\bigl(\rho_{s,p}\colon\mathcal C_{s}\to\mathcal C_{s{\downarrow}p}\bigr).
\]
Equivalently, with $\pi_{0}(\mathbb A(\downarrow))$-indexed families whose component at an orbit $[s]$ is
a small category equipped with an action of the isotropy monoid $\mathrm{Stab}(s)$. Reader
($\mathbb A$ discrete) and State ($\mathbb A$ collapsing to codiscrete) are the two special cases in
which every isotropy action is trivial, so the datum degenerates to a $\pi_{0}$-indexed family of
\emph{plain} categories --- an arbitrary family for Reader, a single category for State.
\prov{MacBeth, proved and exhaustively verified at $|S|=2$, degree-$1$, the argument uniform in
$(S,P,\downarrow)$ (\texttt{2026-08-11-update-monad-liftings-holonomy-full.md}). The update monad and
the data $(S,P,\downarrow)$ are Ahman--Uustalu \cite{AhmanUustalu13}; the $\mathbb A(\downarrow)$
framing, the classifier $\mathbf{Fun}(\mathbb A(\downarrow),\Cat)$, and the
isotropy${}={}$holonomy diagnosis are the author's. \textsf{[open]} at higher object-degree and beyond
the update-monad class; a novelty-check against Uustalu's \emph{container combinatorics} (TTCS 2017) is
owed and deferred.}
\end{Theorem}

\begin{proof}[Why it holds]
Necessity is the displayed law: instantiate $\mu\circ T\mu=\mu\circ\mu T$ on an
$\mathrm{Upd}^{3}$-element with outer update $f$ and constant middles and inners; letting the constants
range over $P$ makes $p=f(s)$ and $q$ range over all of $P$, and the two bracketings agree on the
innermost object exactly when $\rho_{s,p\oplus q}=\rho_{s{\downarrow}p,q}\circ\rho_{s,p}$. That is
functoriality of $\rho$ on $\mathbb A(\downarrow)$. Sufficiency is the converse construction: given any
$F\colon\mathbb A(\downarrow)\to\Cat$, the source-routed aggregator
$A_{f}(Q)=\coprod_{s}\coprod_{c\in\ob\mathcal C_{s}}Q_{s}^{\,\mathrm{out}(c)}$ with $\varepsilon$ the
identities and $\delta$ the composition-with-$\rho$ satisfies the unit and associativity laws \emph{iff}
$\rho$ is functorial. Both directions are the content of the census engine, whose law-checks are
exhaustive over the finite shape spaces at $|S|=2$.
\end{proof}

Functoriality is a weak requirement, and that is the whole point: a functor out of $\mathbb A(\downarrow)$
may send the loops at a state --- the isotropy --- to genuine automorphisms of the fibre category. When
it does, the lifting carries holonomy, and no count of components can see it. The smallest monoid already
does this.

\begin{Example}[$\mathbb Z/2$ acting trivially: four liftings where $\pi_{0}$ predicts one]\label{ex:z2-holonomy}
Take $S=\{0,1\}$ and $P=\mathbb Z/2$ acting \emph{trivially}, $s{\downarrow}p=s$. Then
$\mathbb A(\downarrow)=B\mathbb Z/2\ \sqcup\ B\mathbb Z/2$: two components, each a one-object category
whose endomorphism monoid is the \emph{group} $\mathbb Z/2$, since $\mathrm{Stab}(0)=\mathrm{Stab}(1)=
\mathbb Z/2$. Its $\pi_{0}$ is $2$ --- identical to a two-leaf Reader, whose discrete action category
forces the \emph{only} transport on each component to be the identity. Yet here the census is decisive:
on a two-object fibre the unit forces $\rho_{s,o}=\id$, and
functoriality forces the non-trivial loop $\rho_{s,1}$ to square to the identity, i.e.\ into the isotropy
$\mathbb Z/2$ acting on $\mathcal C_{s}$ --- leaving the identity or the swap, \emph{independently at each
of the two states}:
\[
   (\rho_{0,1},\rho_{1,1})\in\{(\id,\id),\,(\id,\mathrm{swap}),\,(\mathrm{swap},\id),\,(\mathrm{swap},\mathrm{swap})\}.
\]
All four satisfy the monad laws on every one of the $16384$ shapes checked, and an intertwiner
computation shows they are \emph{pairwise non-isomorphic} as liftings --- e.g.\ $(\id,\id)$ and
$(\mathrm{swap},\mathrm{swap})$ cannot agree, since $\varphi_{s}\circ\id=\mathrm{swap}\circ\varphi_{s}$
forces $\mathrm{swap}=\id$. So $\mathbb Z/2$-acting-trivially has $\pi_{0}=2$ and \textbf{four}
non-isomorphic liftings: one $\mathbb Z/2$ of holonomy per orbit. The count $\pi_{0}$ said ``one lifting
per component''; the answer is ``one \emph{representation of the isotropy} per component,'' and the
representation is real.
\prov{MacBeth, exhaustively verified (\texttt{2026-08-11-update-monad-liftings-holonomy-full.md} \S3):
four law-satisfying transports, pairwise non-isomorphic, over all $16384$ $\mathrm{Upd}^{3}$ shapes on
the two-object fibre.}
\end{Example}

With the counterexample in hand, the two poles resolve into a single, sharp dichotomy.

\begin{teachbox}[Trivial two ways: why Reader and State are holonomy-free]
Holonomy-freeness --- the collapse to $\pi_{0}$-indexed \emph{plain} categories --- holds exactly when
every component of $\mathbb A(\downarrow)$ is holonomy-trivial, i.e.\ every functor out of it is constant
up to canonical iso. There are two disjoint ways for a monoid to force this, and Reader and State are one
each.
\begin{itemize}[nosep,leftmargin=1.4em]
  \item \textbf{Reader: the action category is discrete.} $P=1$, so the only arrow is the identity, the
        isotropy is trivial, and there are no cross-state arrows at all. There is no transport
        \emph{to be} nontrivial: holonomy is absent because there is nothing to transport along.
        $\mathbf{Fun}(S_{\mathrm{disc}},\Cat)$ is an $S$-indexed family of categories, $\pi_{0}=|S|$.
  \item \textbf{State: reset elements erase the label.} The overwrite monoid has, for each $m$, a
        \emph{reset} $w_{m}$ with $s{\downarrow}w_{m}=m$ for \emph{all} $s$. In the associativity
        instance a reset middle forces the endpoint to $m$ regardless of the source, degenerating the
        deepest-object law to \emph{endpoint-locality}: $\rho$ depends on the update only through its
        endpoint. That is exactly a functor out of the \emph{codiscrete} category on $S$ --- a trivial
        coherent system of isomorphisms --- which fuses all fibres into one category and trivialises the
        transport. Here holonomy is not absent; it is \emph{erased}. $\pi_{0}=1$.
\end{itemize}
So the ``invisible'' store multiplication of \S\ref{sec:state-liftings} was not empty --- it was a reset
performing an erasure, and the erasure is what we mistook for absence. A generic monoid supplies neither
escape. (A \emph{third} trivialisation exists but is neither pole: let $\mathbb Z/2$ act \emph{freely},
$s{\downarrow}1=1-s$; then $\mathbb A(\downarrow)$ is codiscrete with \emph{trivial} isotropy, giving a
single lifting --- holonomy-free by triviality of the stabiliser, not by erasure. Discrete, reset-collapse,
and free are three roads to the same cliff; the generic monoid takes none of them.) Holonomy-freeness is
the special case. The rule is holonomy.
\end{teachbox}

\begin{teachbox}[An outside view: the same collapse in belief propagation]
The collapse mechanism --- \emph{trivial holonomy of a transport around cycles forces locally-consistent
data to glue into one global object} --- is not special to containers. It recurs, provably and
independently, in categorical belief propagation: ter Horst, Mahadevan and Zambrano \cite{CBP} prove
(their Thm.~6.14) that locally consistent beliefs on a factor graph glue to a single global descent datum
\emph{exactly} when the cycle-transport holonomy is trivial. There is no shared vocabulary --- their
setting is sheaves on factor graphs, not polynomial monad liftings --- and this is emphatically \emph{not}
a container result. It is a structural resonance: the ``trivial holonomy $\Rightarrow$ collapse to a global
section'' skeleton appears to be a genuine cross-domain pattern, and State's collapse is one more instance
of it. Its converse --- \emph{nontrivial} holonomy obstructs the global glue --- is exactly what
Example~\ref{ex:z2-holonomy} exhibits on the container side: four local pieces that refuse to become one.
\end{teachbox}

\begin{teachbox}[The holonomy is a group representation --- and composition remembers it]
Step back and name what the isotropy is doing. A functor $\mathbb A(\downarrow)\to\Cat$ restricted to the
loops at a state $s$ is an action of the stabiliser $\mathrm{Stab}(s)$ on the fibre category
$\mathcal C_{s}$ --- a \emph{representation of the isotropy group} on a category. Example~\ref{ex:z2-holonomy}
is the regular $\mathbb Z/2$-representation, twice. This is a \emph{second-order} compositional datum:
not the behaviour of the lifting but the symmetry it carries, and it survives composition. It is the exact
sibling of the class $[\omega]\in H^{2}$ that governs the directed / Zappa--Szép mode of composition
(Chapter~\ref{ch:zs}, Theorem~\ref{thm:h2}), where reentrancy is obstructed by a cohomology class rather
than a group representation. Two modes, two second-order invariants, one lesson: \emph{composition remembers a group}.
For orchestration this is the payoff. A threading agent whose update monoid has nontrivial isotropy does
not merely have a family of behaviours; it carries a representation, and orchestrating it composes that
representation into the whole. The two ``clean'' monads --- Reader and State --- are precisely the agents
whose representation happens to be trivial. Everything in between remembers.
\end{teachbox}

The section closes with two forward-glances --- one example still to be filled, one line of
constructions still to be built.

\begin{teachbox}[Higher-order trees, decoded from the free-monad formula]
Read a container $C=\cont S P$ as a \emph{signature}: shapes $S$ are operation symbols, positions
$P(s)$ are the arity of $s$. The free monad of \S\ref{sec:moncomon-free} then \emph{is} the type of
finite $C$-terms: Definition~\ref{def:freemonad} gives $S^{*}=\mu Y.(1+\sum_{s}(P(s)\to Y))$, the
well-founded trees whose nodes are operations and whose $P(s)$ children are the operands, with
variable leaves for the free generators. This is the whole content of ``the free monad is syntax'':
$C^{*}$ is exactly the abstract-syntax-tree type of the signature $C$, and the free-monad Lemma
(Theorem~\ref{thm:free-lemma}) says a map out of it is fixed by its action on the one-node terms.

A \emph{higher-order} signature is one whose operations take operations as arguments --- so that an
arity $P(s)$ is itself a \emph{function type} $A\to B$ rather than a bare set, and a node's children
are indexed by such a type. Feeding such a $C$ into the same formula, $C^{*}$ becomes the type of
higher-order trees: nodes whose subtree-families range over an exponential. Nothing new is proved;
the higher-order tree type is \emph{read off} the free-monad formula by taking the arities to be
function types. The natural such example is Ghani--Kurz's $n$-dimensional trees, staged to
foreshadow the opetopic tower.

\textbf{Signature TODO --- ask Neil / browse.} The \emph{exact} Ghani--Kurz $n$-dimensional-tree
signature $\cont S P$ --- which operation symbols, which function-typed arities --- is not yet in
this book's reading set, so the specific container is left unfilled here rather than guessed. Once
it is deep-read, instantiate $C=\cont S P$ above and check that $C^{*}$ reproduces their tree type
verbatim. Do not invent the signature before reading it.
\end{teachbox}

\begin{teachbox}[On the horizon: from the store comonad to Workers]
Neil's programme continues past the transfer with a specific comonad --- the \emph{store}
comonad of a state set $S$ --- and the arrows it defines. The next section,
\S\ref{sec:moncomon-workers}, develops these: the \emph{reader/store} modality
$\Delta S\otimes(-)$ on $\Cont$, its \emph{Kleisli} arrows $\Delta S\otimes p\to q$
(``$p$-to-$q$ transformations with access to state $S$''), and the \emph{category of Workers}
they assemble into --- a category \emph{graded} by $(\Set,\times)$, in which composing two
workers multiplies their state spaces to $S\times T$. One entry on Neil's list is left open
here: the \emph{oracle/continuation}, a coalgebra $S\to\Ext{[p,q]}(S)$ into the internal-hom
container $[p,q]$, read as a stateful container morphism answering $p$-queries with
$q$-responses; whether these organise into a monad on $\Cont$ is \textsf{[to be developed]}.
\end{teachbox}

\section{Processes with state: the store comonad and the category of Workers}
\label{sec:moncomon-workers}

The transfer produced comonads on $\Cont$ abstractly, from any base monad. This section works one
comonad concretely --- the \emph{store} comonad of a state set --- and follows it to its natural end:
a category whose morphisms are \emph{processes carrying state}, and whose composition \emph{multiplies}
the state. It is the point at which the machinery of this chapter touches the everyday object a
programmer calls a stateful component, and an economist calls a process with memory. Neil's prompt for
it was the reader monad: ``the basic definitions of the reader monad \ldots\ but different processes
have different state, so there is something more to be said.'' The ``something more'' is exactly the
writeback that turns the read-only reader into the read--write store, and --- pushed through composition
--- forces the state spaces to multiply.

\subsection{The state object and the store comonad}
\label{sec:workers-store}

\begin{Definition}[State object]\label{def:state-object}
For a set $S$, the \emph{state object} is the container $\Delta S:=\cont S{(s\mapsto S)}$: shape set
$S$, and \emph{every} position fibre equal to $S$. Its extension is
$\Ext{\Delta S}X=\sum_{s:S}X^{S}=S\times X^{S}$.
\prov{Definition}
\end{Definition}

$\Delta S$ is the container in which each shape sees all of $S$ as its positions. Read through the
object dictionary of Chapter~\ref{ch:comonoids}, that uniform fibre has a name.

\begin{Proposition}[$\Delta S$ is codiscrete; $\Ext{\Delta S}$ is the store comonad]
\label{prop:store}
Under the correspondence between directed containers and small categories
(Theorem~\ref{thm:objdict}), $\Delta S$ is the \emph{codiscrete} category on $S$ --- object set $S$,
a unique morphism for each ordered pair $(s,s')$ --- with directed structure
\[
  \mathsf o(s)=s,\qquad s\downarrow p=p,\qquad p\oplus p'=p'\ \ (\text{second projection}).
\]
The comonad it induces on $\Cont$ is the \emph{store} (costate) comonad
\[
  \Ext{\Delta S}X=S\times X^{S},\qquad
  \varepsilon(s,v)=v(s),\qquad
  \delta(s,v)=\bigl(s,\ \lambda p.\,(p,v)\bigr),
\]
with $v\colon S\to X$.
\prov{MacBeth (identification); store comonad due to Uustalu--Vene \cite{UustaluVene08};
Lean-verified in \texttt{StateComonad.lean}\protect\footnotemark}
\end{Proposition}
\footnotetext{The directed-container laws D1--D5 for $\Delta S$ and the store comonad's counit/comultiplication
laws are machine-checked in \texttt{StateComonad.lean}, which is present in the repository's \texttt{lean/}
tree but not committed as source to this book's tree; the weaker in-text tag is used per the Preface
convention.}

\begin{proof}[Proof sketch]
The codiscrete category has $\mathcal C(s,s')=1$, so
$P(s)=\coprod_{s'}\mathcal C(s,s')=\coprod_{s'}1=S$, which is $\Delta S$. Reading Theorem~\ref{thm:objdict}
off this category: the identity at $s$ is the arrow $s\to s$, so $\mathsf o(s)=s$; a position $p\in P(s)=S$
names its codomain, so $s\downarrow p=\cod(p)=p$; the unique composite $s\to s'\to s''$ has codomain
$s''$, so $p\oplus p'=p'$. Substituting these into the comonad induced by a directed container
(Theorem~\ref{thm:comonad}), $\varepsilon(s,v)=v(\mathsf o(s))=v(s)$ and
$\delta(s,v)=(s,\lambda p.(s\downarrow p,\lambda p'.v(p\oplus p')))=(s,\lambda p.(p,v))$, the store comonad.
\end{proof}

Two special cases pin the reader/store story. Feeding the tensor unit, $\Delta S\otimes\y=\Delta S$, so
the store comonad \emph{is} the modality $\Delta S\otimes(-)$ evaluated at $\y$. And the contravariant part
of $\Ext{\Delta S}$ --- the exponential $X^{S}=(S\to X)$ --- is precisely the \emph{reader} functor, a monad
with $\eta\,x=\lambda s.\,x$ and $\mu\,g=\lambda s.\,g\,s\,s$. Thus
\[
  \Ext{\Delta S}\;=\;\underbrace{S\times{}}_{\text{current state}}\ \underbrace{(-)^{S}}_{\text{reader}},
\]
the reader with a current-state factor bolted on. That extra $S\times$ --- the shape, i.e.\ the state a
morphism can \emph{write back} --- is Neil's ``something more''. The rigidity is genuine: perturbing a single
fibre of $\Delta S$ away from $S$ destroys the constant-fibre profile, and the store counit and
comultiplication no longer typecheck (a negative control in the companion computations). The store structure
lives exactly in the ``all fibres $=S$'' condition.

\subsection{Workers, and why the state multiplies}
\label{sec:workers-cat}

\begin{Definition}[Worker]\label{def:worker}
A \emph{Worker} from $p$ to $q$ \emph{with state $S$} is a container morphism $w\colon\Delta S\otimes p\to q$.
\prov{Definition}
\end{Definition}

Writing $q=\cont CD$ and $\Delta S\otimes p=\cont{S\times A}{\bigl((s,a)\mapsto S\times B(a)\bigr)}$, such a
$w$ unpacks into a forward map and a \emph{split} backward map:
\begin{itemize}[nosep]
  \item forward $f\colon S\times A\to C$ --- from state $s$ and input shape $a$, the output shape $f(s,a)$;
  \item backward $f^{\sh}\colon D(f(s,a))\to S\times B(a)$, i.e.\ a pair
    $\bigl(f^{\sh}_{1}(s,a)\colon D(f(s,a))\to S,\ f^{\sh}_{2}(s,a)\colon D(f(s,a))\to B(a)\bigr)$: the first
    component \emph{writes back} a new state at each output position, the second is the ordinary
    position back-map.
\end{itemize}
So a Worker reads a state, transforms a $p$-interface into a $q$-interface, and writes a state back at each
response --- the read--write process the store comonad was built to carry.

The engine of composition is a single strict identity of containers.

\begin{Lemma}[The state multiplies]\label{lem:delta-mult}
$\Delta S\otimes\Delta T=\Delta(S\times T)$ (strict equality of containers), and $\Delta 1=\y$. Hence
$S\mapsto\Delta S$ is a strict monoidal map $(\Set,\times,1)\to(\Cont,\otimes,\y)$ on objects.
\prov{MacBeth, Lean-verified in \texttt{Workers.lean}/\texttt{StateComonad.lean}\protect\footnotemark}
\end{Lemma}
\footnotetext{$\Delta S\otimes\Delta T=\Delta(S\times T)$ holds by \texttt{rfl} in \texttt{StateComonad.lean};
as with Proposition~\ref{prop:store}, the file is in the \texttt{lean/} tree but not committed to this book's
tree.}

\begin{proof}
$\Delta S\otimes\Delta T=\cont{S\times T}{\bigl((s,t)\mapsto S\times T\bigr)}=\Delta(S\times T)$: the shape
sets and every position fibre coincide. $\Delta 1=\cont 1{(*\mapsto 1)}=\y$.
\end{proof}

\begin{Remark}[Why $\otimes$, and not $\times$]\label{rem:why-otimes}
Lemma~\ref{lem:delta-mult} is exactly why the tensor is the \emph{Dirichlet} $\otimes$, whose positions
multiply. Under the categorical product, $\Delta S\times\Delta T=\cont{S\times T}{\bigl((s,t)\mapsto S+T\bigr)}$
has fibres of size $|S|+|T|$, not $|S\times T|$ --- so $\Delta S\times\Delta T\neq\Delta(S\times T)$, and the
state would \emph{not} multiply. Only the tensor whose positions multiply makes the context multiply; the
choice of $\otimes$ is forced by the phenomenon it is meant to model.
\end{Remark}

Given $w\colon\Delta S\otimes p\to q$ and $w'\colon\Delta T\otimes q\to r$, Lemma~\ref{lem:delta-mult} and
functoriality of $\otimes$ compose them into a single Worker:
\[
  \Delta(T\times S)\otimes p
  \;=\;\Delta T\otimes(\Delta S\otimes p)
  \xrightarrow{\ \Delta T\otimes w\ }\Delta T\otimes q
  \xrightarrow{\ w'\ } r .
\]
A state-$S$ worker followed by a state-$T$ worker is a worker whose state is the \emph{product} of the two:
composition \emph{accumulates} state. Applying the outer worker first, the coordinates land the state in
$T\times S$; via the symmetry $\Delta T\otimes\Delta S\cong\Delta S\otimes\Delta T$ (both $=\Delta(S\times T)$,
Lemma~\ref{lem:delta-mult}) we record the accumulated grade in the reading order $S\times T$, matching the
grade multiplication of Theorem~\ref{thm:workers} below. In coordinates, with $r=\cont EG$ and
$e:=g(t,f(s,a))$, the composite has forward $((t,s),a)\mapsto g(t,f(s,a))$ and, backward at a response
$d'\in G(e)$, the state component
$\bigl(g^{\sh}_{1}(t,f(s,a),d'),\,f^{\sh}_{1}(s,a,g^{\sh}_{2}(t,f(s,a),d'))\bigr)\in T\times S$ and the
position component $f^{\sh}_{2}(s,a,g^{\sh}_{2}(t,f(s,a),d'))\in B(a)$. The identity is
$\id_{p}\colon\Delta 1\otimes p=p\to p$, with state $1$.

\begin{Theorem}[The category of Workers]\label{thm:workers}
There is a category \emph{graded by $(\Set,\times)$}, \Workers{}, with objects the containers,
graded hom-sets
\[
  \Workers{}_{S}(p,q):=\Cont(\Delta S\otimes p,\ q),
\]
identity in grade $1$, and composition
$\Workers{}_{S}(p,q)\times\Workers{}_{T}(q,r)\to\Workers{}_{S\times T}(p,r)$ as above,
associative and unital up to the associator and unitors of $(\Set,\times)$. Equivalently,
\Workers{} is the coKleisli category of the $(\Set,\times)$-graded comonad $S\mapsto\Delta S\otimes(-)$
on $\Cont$, whose comultiplication is the strict equality $\Delta(S\times T)\otimes-=\Delta S\otimes\Delta T\otimes-$
of Lemma~\ref{lem:delta-mult}.
\prov{MacBeth, Lean-verified in \texttt{Workers.lean}\protect\footnotemark}
\end{Theorem}
\footnotetext{\texttt{Workers.lean} defines the composition and checks the three graded-category laws
(unit-left, unit-right, associativity), sorry-free; the finite cases of the underlying identities were
independently exhausted by direct computation. As above, the file is not committed to this book's tree.}

\begin{proof}[Proof sketch]
Composition lands in grade $S\times T$ by Lemma~\ref{lem:delta-mult}. Unitality: $\id_{q}\circ w$ has grade
$1\times S$ and, on unpacking, is $w$ under the unitor $1\times S\cong S$; symmetrically on the right.
Associativity: both bracketings of $w_3\circ w_2\circ w_1$ have underlying morphism
$w_3\circ(\Delta U\otimes w_2)\circ(\Delta U\otimes\Delta T\otimes w_1)$ in $\Cont$, equal by associativity of
composition and functoriality of $\otimes$; the grades $(U\times T)\times S$ and $U\times(T\times S)$ agree
by the associator of $(\Set,\times)$. The coKleisli reading holds because the comultiplication is the strict
identity of Lemma~\ref{lem:delta-mult}. The graded-category laws are the content of \texttt{Workers.lean}.
\end{proof}

\begin{Remark}[$S$ varies: the Para reading, and the nearest neighbour]\label{rem:workers-para}
Collapsing the grade into a single hom-set, $\textsc{Para}(p,q):=\sum_{S:\Set}\Cont(\Delta S\otimes p,q)$
--- a Worker \emph{together with} its choice of state --- composed by tensoring the states, is
Gavranovi\'c's \emph{Para} construction \cite{Gavranovic} for the action $S\cdot p=\Delta S\otimes p$ of
$(\Set,\times)$ on $\Cont$. This identification is honest only over the core groupoid $\mathrm{Core}(\Set)$,
the level at which $\Delta$ is functorial: a bare function $S\to S'$ gives no canonical morphism
$\Delta S\to\Delta S'$ (the backward part would need $S'\to S$), so over all of $\Set$ we have a
\emph{graded} category, not a strict actegory. The Para identification and its exact match to Gavranovi\'c's
axioms are \emph{computed / further work}, not proved.

The nearest neighbour in the literature is Capucci--Myers' \emph{contextads} \cite{CapucciMyers}. Their
Example~3.24 --- every $M$-graded comonad is a \emph{trivially-fibred contextad}, i.e.\ a colax action
$\mathcal C\times M\to\mathcal C$ --- identifies \Workers{} precisely: with $M=(\Set,\times)$ it is the
trivially-fibred, colax-action corner of their framework. Their Theorem~A.4 ($\mathrm{Kl}(T)\cong\mathrm{Ctx}(\odot)$,
which they leave unproven) sits in the \emph{opposite}, dependency-essential corner: there the grade is a
dependent family $X\to S$ over the fixed shapes of one monad, multiplying via that monad's own $\mathsf{seq}$,
whereas ours is an external grade ranging over all of $(\Set,\times)$ and multiplying cartesianly. We settle
\emph{no} fragment of their A.4; \Workers{} and A.4 are the two ends of the same axis.
\prov{Para: MacBeth, computed / further work; citation \cite{Gavranovic} at browse-summary level
(deep-read to-do). Neighbour placement vs \cite{CapucciMyers} (deep-read): MacBeth.}
\end{Remark}

\paragraph{For the grant: two axes of composition.} \Workers{} isolates one of two orthogonal axes
along which compositional systems combine. Here the context \emph{multiplies}, $S\times T$; the grade
accumulates monotonically and there is \emph{no obstruction} --- every pair of Workers composes, always. This
is the \emph{state} axis. It is dual to the \emph{directed} axis of Chapter~\ref{ch:zs}, where composing two
directed systems is a Zappa--Sz\'ep problem that \emph{may fail to have a solution}, the failure measured by
a class $[\omega]\in H^{2}$ (Theorem~\ref{thm:h2}). State composition is unobstructed but accumulating;
directed composition is non-accumulating but obstructed. A realistic compositional system --- an
orchestration of stateful agents, a supply chain of processes with memory --- lives on both axes at once, and
the grant's compositional-correctness story is the joint account: track the state grade for what the system
\emph{remembers}, track the $H^{2}$ class for whether it \emph{composes at all}. A \emph{third} axis --- the
effect--coeffect one, obstructed by branching --- joins these two once the arrow category of the next chapter
is in hand; the three are assembled into a single table in \S\ref{sec:threemodes}.

\begin{teachbox}[Forward pointer: a Worker as an effect--coeffect arrow]
A Worker $\Delta S\otimes p\to q$ reads a coeffect (the incoming state, a \emph{comonadic} context) and,
through its writeback, produces an effect (the outgoing state). It is thus one instance of a general
\emph{effect--coeffect arrow}, in which a comonad supplies the context and a monad the output, welded by a
distributive law relating the two feeds of \S\ref{sec:moncomon-entwine} --- read now as the compositor of a
Freyd/arrow category rather than of algebras. That arrow category, and the branching condition under which
it exists, are developed as the \emph{third} mode of composition in \S\ref{sec:threemodes}
(Theorem~\ref{thm:arrows}): the arrows form a category iff the effect monad is non-branching. One
sub-question is left open here: whether the store comonad's Workers \emph{embed into}, or \emph{generate},
that arrow category, and how the state grade of \S\ref{sec:workers-store} interacts with an effect grade,
is \textsf{[in development]} --- recorded as a pointer, with no result claimed.
\end{teachbox}

\section{Syntax and behaviour}
\label{sec:moncomon-outlook}

Step back and the two constructions line up as a single duality. The free monad is
\emph{syntax}: the finite terms a signature generates, composed by substitution, universal
among monads --- the moves a system \emph{can make}. The cofree comonad is \emph{behaviour}:
the possibly-infinite execution trees a signature can unfold, decomposed by taking subtrees,
cofree among comonads --- the observations a system \emph{can offer}. Both stay inside
$\Cont$; both are computed by a (co)inductive tree construction; and the position data --- leaves
for terms, nodes for behaviours --- is exactly what the fixed-point recursion produces.

The asymmetry that made the cofree side the easy side (Theorem~\ref{thm:char}: connected-limit
preservation) is the same asymmetry that makes the cofree side land back on the equivalence
chain. The free monad is a \emph{monad} and stops there; monads for $\triangleleft$ are not, in
general, categories. But the cofree \emph{comonad} is a comonoid, and every comonoid for
$\triangleleft$ is a directed container --- a category. This is why behaviour, and not syntax,
is where the geometry of Chapter~\ref{ch:comonoids} reappears: the space of behaviours of any
container is organised by its subtree category.

\paragraph{One op, two faces.} Both machines of this chapter are the fibrewise $(-)\op$ doing its
work, and it is worth seeing that they apply it to \emph{different targets}. A container morphism
runs backward on positions --- $\Cont=\int(\cod)\op$, the teachbox after
Proposition~\ref{prop:transfer} --- so every position fibre carries a built-in $(-)\op$. The
\emph{transfer} applies that op to the fibre \emph{object}: it turns the base monad $M$ into the
fibrewise comonad $(M\op)_{*}$, so a monad on the base becomes a comonad on $\Cont$
(Proposition~\ref{prop:transfer}). The \emph{free/cofree} pair applies the same op to the
\emph{recursion scheme}: dualising an initial algebra to a final coalgebra, well-founded
substitution-syntax $\mu$ to possibly-infinite behaviour-syntax $\nu$, sends the free monad to the
cofree comonad. One op, two faces --- on the fibre object it exchanges monad and comonad; on the
recursion scheme it exchanges syntax and behaviour --- and the entwining
(Theorem~\ref{thm:entwine}) is what happens when the first face meets a second copy of $M$: the two
container feeds of one monad, welded by $M$'s oplax product-comparison.

\paragraph{For the grant.} Free monads and cofree comonads are the syntax/behaviour pair of a
compositional system: its available moves and its observable unfolding. A verified account of
their composition would let one certify, for instance, that an agent-orchestration combinator
preserves the update (\emph{Put}) structure of its parts (Corollary~\ref{cor:lens}), by
tracking it through the free monad of the combined signature. The free-monad Lemma
(Theorem~\ref{thm:free-lemma}) is the tool that makes such a combinator \emph{definable} by its
action on generators alone --- the universal property doing the work a bespoke construction
would otherwise have to.

\paragraph{What remains open.} Two loose ends are worth naming. First, the cofree side awaits
its Lean certificate: the $M$-type carrier needs a coinductive framework
(Theorem~\ref{thm:cofree-dircont}, footnote), which is an infrastructure step, not a
mathematical one. Second, the free-monad Lemma is proved in coordinates but only partly
machine-checked (Theorem~\ref{thm:free-lemma}); completing the multiplication-homomorphism law
in Lean is the next formalisation target. Neither gap touches the constructions themselves,
which are explicit and, in the finite cases of Examples~\ref{ex:free} and \ref{ex:cofree},
directly computable.

\chapter{Composing systems: Zappa--Sz\'ep and distributive laws}
\label{ch:zs}

Chapter~\ref{ch:comonoids} settled what a single directed container is and what
maps it has. This chapter asks how \emph{two} of them compose into one --- the
question with the economic consequences (mis-composed supply chains lose liability
traceability; a wrong contract composite is an exploit). We package this as a
\emph{Zappa--Sz\'ep} (two-sided semidirect / matched-pair) product of small
categories, the categorical generalisation of the group/monoid Zappa--Sz\'ep
product.

\section{The pairwise criterion}
\label{sec:zscriterion}

For objects $a,b$ of a small category, $\Hom(a,b)$ is canonically a right
$\End(a)$-set (precomposition) and a left $\End(b)$-set (postcomposition). A
Zappa--Sz\'ep decomposition $K=\mathcal C\bowtie\mathcal D$ of $K$ over a wide
subcategory $\mathcal D$ asks for a \emph{unique factorisation} of each morphism
through a transversal and an endomorphism.

\begin{Theorem}[Pairwise Zappa--Sz\'ep criterion]\label{thm:zscriterion}
Let $\mathcal D$ be a wide subcategory of $K$. Then $K=\mathcal C\bowtie\mathcal D$
for a complementary wide subcategory $\mathcal C$ if and only if
\begin{itemize}[nosep]
  \item \textup{(L)} each hom-presheaf $\Hom_K(-,b)$ is a \emph{free} right
        $\mathcal D$-set (equivalently: the source-oriented factorisation
        $f=t\circ e$ is unique for every $f$), and
  \item \textup{(G)} the free bases close into a wide subcategory $\mathcal C$.
\end{itemize}
\prov{MacBeth}%
\footnote{This is the corrected form of the SEED conjecture ``every small category
is a Zappa--Sz\'ep product of a poset skeleton and its vertex monoids,'' which was
\emph{refuted} computationally (smallest counterexample: two objects, one arrow
$s\colon a\to b$, $\End(a)=\End(b)=\mathbb Z/2$ fixing $s$, so
$2\nmid 1=|\Hom(a,b)|$). The freeness half (L) is written up self-contained in
\texttt{papers/pairwise-zappa-szep.tex} (a built PDF is in the repo; the
\texttt{.tex} source is not yet committed to this tree). A brute-force enumeration
of $\approx 26$ small categories matches the criterion exactly.}
\end{Theorem}

\begin{Theorem}[Cocycle axioms $\Longleftrightarrow$ associativity]\label{thm:zsassoc}
The Zappa--Sz\'ep cocycle axioms ZS1--ZS4 (relating the two actions of a matched
pair) hold if and only if the product on $\mathcal C\bowtie\mathcal D$ is
associative.
\prov{MacBeth}%
\footnote{MacBeth has a Lean formalisation of the monoid case (ZS1--ZS4
$\Longleftrightarrow$ associativity of the Zappa--Sz\'ep product) in
\texttt{ZappaSzep.lean}, built and reported zero-\texttt{sorry}; that source is
\emph{not yet committed} to this repository's \texttt{lean/} tree, so the tag is
\textsf{[MacBeth]} rather than \textsf{Lean-verified} here.}
\end{Theorem}

\section{(G) as a strict factorisation system, and its $H^2$ obstruction}

\begin{Proposition}[(G) is a strict factorisation system]\label{prop:sfs}
Condition (G) --- that the free bases close into a wide subcategory --- is exactly
the statement that the decomposition extends to a \emph{strict factorisation
system} on $K$. Distributive laws of categories correspond to strict factorisation
systems.
\prov{Cited: Rosebrugh--Wood (distributive laws $\Longleftrightarrow$ SFS)}
\end{Proposition}

\begin{Theorem}[The (G) obstruction is an $H^2$ class]\label{thm:h2}
Under a trivial-action hypothesis, (G) is classified by a second cohomology class:
$\textup{(G)}\Longleftrightarrow[\omega]=0$, where $[\omega]$ lives in the
\emph{Baues--Wirsching} cohomology $H^2(\mathrm{Sk}_{\mathcal C};\mathcal D)$ of the
$\mathcal C$-skeleton with a source-only (presheaf) natural system of coefficients.
This is the split case of the Baues--Wirsching linear-extension classification.
\prov{MacBeth (identification), resting on Cited: Baues--Wirsching 1985}%
\footnote{The cohomological \emph{identification} --- that the obstruction to (G)
is a Baues--Wirsching $H^2$ class and that ``(G)$\Leftrightarrow[\omega]=0$'' is the
split case of the BW linear-extension theorem --- is MacBeth's; the underlying
classification theorem is Baues--Wirsching, \emph{Cohomology of small categories},
JPAA \textbf{38} (1985) 187--211 (modern account: arXiv:2508.00727). The rigid
twist of the counterexample above is the $\mathbb Z/2$ generator; the nonabelian
boundary is the Kac/Rosebrugh--Wood law.}
\end{Theorem}

\section{The classical anchor}

\begin{Proposition}[Classical Zappa--Sz\'ep recovered]\label{prop:monoidanchor}
Specialising Theorem~\ref{thm:zscriterion} to a one-object category (a monoid)
recovers the classical Zappa--Sz\'ep product of monoids, and (L) becomes the
familiar condition that the factorisation map is a bijection.
\prov{MacBeth}%
\footnote{MacBeth flags an earlier \emph{over}claim for correction: the
``Ahman--Uustalu update monad as a degenerate Zappa--Sz\'ep product'' reading is
\emph{not} a special case of the two-category matched pair --- it is a monoid acting
on a set, not two composable monoids --- and should not be presented as such.}
\end{Proposition}

\section{Three modes of composition}
\label{sec:threemodes}

The book has now built, in three different chapters, three different answers to a
single question --- \emph{when do two compositional systems compose into one, and
what obstructs it?} --- and it is worth setting them side by side, because they do
not agree, and the disagreement is the content. The Zappa--Sz\'ep product of this
chapter welds two \emph{directed} systems (two small categories,
\S\ref{sec:zscriterion}); the category of Workers of \S\ref{sec:moncomon-workers}
welds two \emph{stateful} processes; and the entwining of
\S\ref{sec:moncomon-entwine} welds an \emph{effect} to a \emph{coeffect}. Each is a
genuine mode of composition, each has a precise criterion for when the weld holds,
and the three criteria are three different kinds of mathematics. This is the
grant's compositional-correctness question, and the honest answer is not one
obstruction but a table of them.

\begin{center}
\renewcommand{\arraystretch}{1.35}
\begin{tabular}{@{}llll@{}}
\toprule
Mode & Object welded & Obstruction & Always exists?\\
\midrule
\textbf{Directed} (ZS) & $\mathcal C\bowtie\mathcal D$, two categories
  & $[\omega]\in H^{2}(\mathrm{Sk}_{\mathcal C};\mathcal D)$ & \emph{No} --- may fail\\
\textbf{State} (Workers) & $\Delta S\otimes p\to q$, stateful maps
  & \emph{none} & \emph{Yes} --- grade accumulates\\
\textbf{Effect--coeffect} & arrows $G_{M}p\to T_{M}q$
  & branching ($M$ arity $\ge 2$) & \emph{No} --- iff $M$ non-branching\\
\bottomrule
\end{tabular}
\end{center}

A cohomology class, nothing at all, and a Boolean condition on an arity: three
obstructions of three different types, for three different ways of gluing. The
temptation is to collapse the rows into one master theorem --- ``compositional
systems compose iff a single invariant vanishes'' --- and it should be resisted,
because it is false: $H^{2}$, the empty obstruction, and the branching flag do not
unify into one class. The synthesis that is true, and that the grant rests on, is
weaker and better: \emph{heterogeneous systems weld along an axis, and whether the
weld holds is a computable invariant of the pieces --- but which invariant depends
on the axis.} We take the three rows in turn.

\paragraph{The directed axis --- an obstruction that can genuinely fail.}
Composing two directed containers is the subject of this chapter: it is a
Zappa--Sz\'ep problem, solvable exactly when the free factorisation bases close
into a subcategory (condition (G), Theorem~\ref{thm:zscriterion}), and that closure
is classified by a cohomology class $[\omega]\in H^{2}(\mathrm{Sk}_{\mathcal C};\mathcal D)$
in Baues--Wirsching cohomology (Theorem~\ref{thm:h2}). The class can be nonzero ---
the two-object category with $\mathbb Z/2$ vertex groups fixing a single crossing
arrow is the minimal witness (footnote to Theorem~\ref{thm:zscriterion}) --- so
this mode of composition \emph{can fail}, and when it fails the failure is a
specific, computable cocycle. This is the mode with economic teeth: a mis-composed
supply chain, a wrongly matched pair of contracts, is exactly a nonzero $[\omega]$.

\paragraph{The state axis --- no obstruction, but the context multiplies.}
Composing two Workers, by contrast, \emph{always} succeeds. A Worker is a container
map $\Delta S\otimes p\to q$ carrying a state space $S$
(\S\ref{sec:moncomon-workers}); two of them compose because $\Delta S\otimes\Delta T
=\Delta(S\times T)$ strictly (Lemma in \S\ref{sec:workers-store}), and the composite
simply carries the product state $S\times T$. There is no closure condition to
check and no class to vanish; the price of unconditional composition is that the
grade \emph{accumulates}, monotonically, forever. This is the $(\Set,\times)$-graded
category of \S\ref{sec:moncomon-workers}: composition is total, the obstruction set
is empty, and what a realistic system pays instead is memory.

\paragraph{The effect--coeffect axis --- obstructed by branching.}
The third mode is the one the entwining of \S\ref{sec:moncomon-entwine} was built
for, now read as a composition rule rather than a lifting of algebras. Fix a base
monad $M$ and recall its two container feeds: the effect monad
$T_{M}\cont SP=\cont{MS}{P^{\star}}$ of Ahman--Bauer and the coeffect comonad
$G_{M}\cont SP=\cont S{M\circ P}$ of the transfer. Section~\ref{sec:moncomon-entwine}
assembled these into the \emph{arrow profunctor} $\Arr_{M}(p,q)=\Cont(G_{M}p,\,T_{M}q)$
(Definition~\ref{def:arrowprof}): an effect--coeffect arrow $p\rightsquigarrow q$ reads
its input through the coeffect $G_{M}$ (a comonadic context) and delivers its output
through the effect $T_{M}$. The profunctor exists for every $M$; the question of
\emph{composition} is the question whether $\Arr_{M}$ is the hom-profunctor of a
category, and its biKleisli composite threads one arrow's effect past the next arrow's
coeffect through the \emph{reverse} entwining $\kappa\colon G_{M}T_{M}\Rightarrow T_{M}G_{M}$
--- the orientation opposite to the law $\lambda$ of Theorem~\ref{thm:entwine}. Whether
the arrows compose is therefore whether $\kappa$ is a distributive law, and
\S\ref{sec:moncomon-entwine} already told us the answer: the reverse orientation holds
precisely when $M$ does not branch.

\begin{Theorem}[Effect--coeffect arrows, and their positive class]\label{thm:arrows}
Let $M$ be a cartesian ($\prod$-cointerpretation) $\Set$-monad, with effect monad
$T_{M}$ and coeffect comonad $G_{M}$ as above. The arrow profunctor
$\Arr_{M}(p,q)=\Cont(G_{M}p,\,T_{M}q)$ (Definition~\ref{def:arrowprof}), with the
biKleisli composite whose compositor is the reverse entwining
$\kappa\colon G_{M}T_{M}\Rightarrow T_{M}G_{M}$, is the hom-profunctor of a
category --- indeed a Hughes arrow, equivalently a Freyd category, over
$(\Cont,\times)$ --- \emph{if and only if $M$ is non-branching} (every shape has
arity $\le 1$). The sole obstructed axiom is associativity (the multiplication
axiom $E2'$ of the distributive law). Moreover, a cartesian $M$ is non-branching if
and only if it is a \emph{writer-with-absorbing-exceptions} monad
\[
  M X\ \cong\ E + A\times X,\qquad (A,\cdot,e)\ \text{a monoid},\quad
  (E,\odot)\ \text{a left $A$-set},
\]
with writer unit $x\mapsto(e,x)$ and the log $A$ acting on the exceptions $E$; and
for \emph{every} monad in this class $\kappa$ satisfies all four axioms $E1'$--$E4'$,
so the arrow category exists throughout it.
\prov{MacBeth}%
\footnote{Proved in \texttt{proofs/2026-07-29-effect-coeffect-arrows.md} (the
non-branching criterion, Theorem~A), \texttt{...-arrows-first.md} (the Hughes
arrow / Freyd characterisation, Theorem~B), and
\texttt{proofs/2026-07-30-affine-classification.md} (the positive class
$E+A\times(-)$ and the fact that $E2'$ holds across it). The effect monad $T_{M}$ is
Ahman--Bauer \cite{AhmanBauer24} Thm~6.3; the coeffect comonad $G_{M}$ and the
compositor $\kappa$ are the transfer (Proposition~\ref{prop:transfer}) and the
reverse of Theorem~\ref{thm:entwine}. The unit-law fragment of the arrow category is
machine-checked (\texttt{BiKleisli.lean}, not committed to this tree); the
associativity axiom $E2'$ is proved in coordinates and machine-verified for the
$\Set$-monad examples, not yet Lean-certified in general, so the tag is
\textsf{[MacBeth]}.}
\end{Theorem}

The mechanism behind the last clause is worth one line, because it is why the
positive class is so clean. Non-branching means every element of every $M$-object
has at most one leaf, so every product $\prod_{b}$ that appears in $T_{M}$,
$T_{M}T_{M}$, $G_{M}T_{M}$, $\dots$ ranges over \emph{at most one} factor. The lax
comparison $\kappa$ is then, shape by shape, either the identity (one leaf: rewrap
the single factor) or the monad unit $\eta^{M}\colon 1\to M1$ (no leaf: the empty
product). The two-or-more-leaf comparison --- the \emph{union of products versus
product of unions} discrepancy that Example~\ref{ex:entwine-branching} used to break
the reverse law for $\mathsf{Pf}$ --- never forms, and the associativity axiom $E2'$
degenerates to the associativity of the monoid $A$, which of course holds. Branching
is precisely the appearance of a second leaf, and a second leaf is precisely what
lets the multiplication $\mu^{M}$ merge two positions that the cartesian comparison
cannot keep correlated.

\paragraph{The deepest point: one entwined structure, two faces, and which
direction you commute.} The effect--coeffect obstruction is not really about
\emph{whether} the two feeds of $M$ interact --- they always do --- but about
\emph{which way}. The single monad $M$, feeding $\Cont$ through $T_{M}$ and $G_{M}$,
carries two comparison maps between its composites $T_{M}G_{M}$ and $G_{M}T_{M}$, and
they behave oppositely.
\begin{itemize}[nosep]
  \item The \textbf{bialgebra face} is $\lambda\colon T_{M}G_{M}\Rightarrow G_{M}T_{M}$,
        built from the \emph{oplax} product-comparison $\str$ (Theorem~\ref{thm:entwine}).
        It is an entwining for \emph{every} $M$ --- including the branching ones ---
        and it is what lifts $G_{M}$ to a comonad on $T_{M}$-algebras and $T_{M}$ to a
        monad on $G_{M}$-coalgebras. This is the unconditional half: the affirmative
        answer to the Plotkin--Turi question of whether one monad's effect and coeffect
        feeds admit a bialgebraic semantics. They always do.
  \item The \textbf{arrow face} is $\kappa\colon G_{M}T_{M}\Rightarrow T_{M}G_{M}$, the
        \emph{reverse}, \emph{lax} comparison (Theorem~\ref{thm:arrows}). It is an
        entwining --- and hence turns the arrow profunctor $\Arr_{M}$ into a genuine
        \emph{category} --- \emph{only} when $M$ is non-branching.
\end{itemize}
Same monad, same two feeds, same $\Cont$: the unification of effects and coeffects
\emph{exists} unconditionally as a bialgebra, and its \emph{categorical packaging as
an arrow calculus} is what branching obstructs. The obstruction is not to the
interaction but to its orientation --- to commuting the effect past the coeffect in
the direction that composition of arrows demands. That the same variance which
forced $G_{M}$ to be a comonad (positions run contravariantly, \S\ref{ch:cont})
also forces the \emph{single} legal direction of the interaction is, to this author,
the most satisfying fact in the chapter.

\begin{Remark}[A neighbour with a different engine: interaction laws]\label{rem:kru}
The reader who knows Katsumata, Rivas and Uustalu's \emph{interaction laws of monads
and comonads} \cite{KRU20} should be told how that framework differs from this one,
because it answers a nearby question by other means. Their interaction law is a
\emph{pairing} $\psi\colon TX\times DY\to X\times Y$ --- run an effect against a
coeffect, extract a value and a final state --- realised as a monoid object in a Chu
construction over $([\mathcal C,\mathcal C],\text{Day})$; it is emphatically
\emph{not} a distributive law, and indeed distributive laws are \emph{inputs} to
their machinery (used to glue pairings on composite monads), never outputs. Our
$\kappa$ is a distributive law and nothing else, so the two accounts are orthogonal
in kind --- a pairing versus a compositor --- rather than rival; what they share is
striking only at the level of the obstruction, for their degeneracy theorems
(their Thms~1--3) locate the same branching wall from the side of extensive
categories that our arity criterion locates from the side of polynomial functors.%
\footnote{On \cite{KRU20} at deep-read (full text,
\texttt{reading/2026-07-29-katsumata-rivas-uustalu-1912-13477.md}); the pairing/
compositor distinction and the ``distributive laws are inputs'' reading are verified
against their \S3 and \S6.}
\end{Remark}

\paragraph{For the grant: Path~5.} These three modes are the spine of the grant's
Path~5 --- \emph{a container-native theory of when compositional systems compose}.
An orchestration of stateful agents, a supply chain of processes with memory, an
ontology merge: each is a system that must be glued to another, and each gluing runs
along one or more of these axes at once. The account this book can offer is exactly
the table above: to certify that two systems compose, identify the axis, then
compute the corresponding invariant --- the Baues--Wirsching class for a directed
weld, the accumulated grade for a stateful one, the arity of the effect monad for an
effect--coeffect one. None of the three invariants is a heuristic; each is forced by
the categorical structure and (for the first two, and the finite cases of the third)
machine-checkable. The grant's claim is not that composition is always possible ---
the directed and effect--coeffect rows say plainly that it is not --- but that
\emph{whether} it is possible is, in every mode, a decidable property of the parts.
That is a stronger and more honest promise than a single master obstruction would
have been, and it is what three chapters of container theory have earned the right
to state.

\section{When two of the modes meet: emergent holonomy}
\label{sec:emergent-holonomy}

The table put \emph{State} and \emph{Directed} in different rows, and it was right to do
so: their obstructions are a nothing and a cohomology class, and no honest master theorem
folds them together. But a table is a still photograph, and it can hide what happens
when you make the systems \emph{move}. Compose two \emph{stateful} agents --- two
objects from the State row --- and watch what appears. What appears is a
Zappa--Sz\'ep product: an object from the Directed row. The two modes were never as
far apart as their rows; they meet on a single object, and on that object something is
waiting that neither agent alone contains.

\paragraph{The promise.} Composing two holonomy-free threading agents can
\emph{synthesise} holonomy that neither possesses; a single integer --- the number of
states at which the two agents' reachable-state orbits cross --- measures exactly how
much; and when that integer is $1$ everywhere, a degree-$2$ cohomology class decides
whether the two agents' symmetries stay unentangled or fuse. Orchestration remembers
more than its parts, and we can say precisely how much more.

\paragraph{One agent, recalled.} Chapter~\ref{ch:comonoids}'s climax
(\S\ref{sec:update-liftings-holonomy}) classified the liftings of a \emph{single}
update monad $\mathrm{Upd}_{(S,P,\downarrow)}$ by the functors
$\mathbb A(\downarrow)\to\Cat$ (Theorem~\ref{thm:update-classification}). The one datum
that is not a count is the \emph{holonomy}: the restriction of such a functor to the
loops at a state $s$ --- the isotropy monoid $\mathrm{Stab}(s)$ --- is a representation
of $\mathrm{Stab}(s)$ on the fibre category $\mathcal C_{s}$. Reader and State are the
degenerate agents whose isotropy is trivial; a generic agent carries a real
representation, and orchestrating it must compose that representation into the whole.
This section is about \emph{how}.

\subsection*{Two agents make a Zappa--Sz\'ep product}

Take two update monads that share the state set $S$:
$\mathrm{Upd}_{(S,P,\downarrow)}$ and $\mathrm{Upd}_{(S,P',\downarrow')}$, with update
monoids $P$, $P'$ threading $S$ by $\downarrow$ and $\downarrow'$. To compose them into
one monad is to give a distributive law between them, and a distributive law is exactly
a \emph{matched pair}: mutual actions of $P$ and $P'$ that make the carrier $P\times P'$
a monoid --- the Zappa--Sz\'ep product $P\bowtie P'$ (\S\ref{sec:zscriterion};
distributive laws are strict factorisation systems, Rosebrugh--Wood \cite{RW}), with
\[
   (p,p')\cdot(q,q')\ =\ \bigl(\,p\,(p'\triangleright q)\,,\ (p'\triangleleft q)\,q'\,\bigr),
   \qquad\text{unit }(o,o').
\]
This is a Zappa--Sz\'ep product of two \emph{monoids}, the classical anchor of
Proposition~\ref{prop:monoidanchor} --- two composable agents, not the mis-reading (a
single update monad as a degenerate matched pair) that its footnote warns against. The
matched pair also threads the shared state: $s\downarrow_{\bowtie}(p,p'):=(s\downarrow
p)\downarrow' p'$ is a right action of $P\bowtie P'$ on $S$, and the composite functor
is again an update monad --- the one built from $P\bowtie P'$ acting on $S$. So the
classification of \S\ref{sec:update-liftings-holonomy} applies verbatim to the
composite, and it applies through the Zappa--Sz\'ep product.

\begin{Theorem}[The classifier composes by the Zappa--Sz\'ep product]\label{thm:classifier-composes}
The composite of $\mathrm{Upd}_{(S,P,\downarrow)}$ and $\mathrm{Upd}_{(S,P',\downarrow')}$
is the update monad of $P\bowtie P'$ acting on $S$ by $\downarrow_{\bowtie}$. Its
degree-$1$ proof-relevant liftings are therefore classified by
\[
   \mathbf{Fun}\bigl(\mathbb A(\downarrow_{\bowtie}),\Cat\bigr)
   \ =\
   \mathbf{Fun}\bigl(\mathbb A(\downarrow)\bowtie\mathbb A(\downarrow'),\Cat\bigr),
\]
where $\mathbb A(\downarrow)\bowtie\mathbb A(\downarrow')$ is the Zappa--Sz\'ep product
of the two action categories: objects $S$, an arrow $s\xrightarrow{(p,p')}(s\downarrow
p)\downarrow' p'$ for each matched pair, composed by the Zappa--Sz\'ep multiplication.
\prov{MacBeth (\texttt{2026-08-12-holonomy-composition-zs-bridge.md}, part (a)), on
Theorem~\ref{thm:update-classification} and the update monad of Ahman--Uustalu
\cite{AhmanUustalu13}; the matched-pair composition is the monoid case of
\S\ref{sec:zscriterion} (\cite{RW}). \textsf{[open]} at higher object-degree, inherited.}
\end{Theorem}

The classifier of the Ch.~\ref{ch:comonoids} climax is \emph{monoidal under
orchestration}: the classifying object of the composite is built functorially from the
two factor action categories by the same Zappa--Sz\'ep product that welds two directed
systems in \S\ref{sec:zscriterion}. This is the meeting of the two rows, stated as an
equality of classifiers. Now for what lives at the meeting.

\subsection*{The naive guess, and why it is wrong}

Each agent has, at a state $s$, its own isotropy: $\mathrm{Stab}_{P}(s)$ and
$\mathrm{Stab}_{P'}(s)$, and hence its own holonomy. The composite has an isotropy too
--- $\mathrm{Stab}_{P\bowtie P'}(s)$, the loops at $s$ in the product action category.
The natural guess, the one the table invites, is that the composite holonomy is just
\emph{the two together}:
\[
   \text{(guess)}\qquad
   \mathrm{Stab}_{P\bowtie P'}(s)\ \overset{?}{\cong}\
   \mathrm{Stab}_{P}(s)\bowtie\mathrm{Stab}_{P'}(s).
\]
Half of it is true and free: if $p$ fixes $s$ and $p'$ fixes $s$ then so does $(p,p')$,
so the right side always \emph{includes} into the left. The equality is the claim that
every loop of the composite factors into a loop of each agent --- that the composite
remembers nothing the parts did not already carry.

It is false. And its failure is the phenomenon.

\begin{Example}[$S_{3}$ on three points: holonomy from nowhere]\label{ex:s3-emergent}
Let $G=S_{3}$ act on $S=\{1,2,3\}$, and factor it exactly as $G=P\cdot P'$ with
$P=A_{3}=\langle(123)\rangle$ the rotations and $P'=\langle(12)\rangle$ a single
transposition. Every permutation is uniquely a rotation times a chosen transposition
($|A_{3}|\cdot|{\langle(12)\rangle}|=6=|S_{3}|$ and $A_{3}\cap\langle(12)\rangle=\{e\}$),
so this is a genuine Zappa--Sz\'ep factorisation, and it acts on $S$. Look at $s=1$.
\begin{itemize}[nosep,leftmargin=1.4em]
  \item \textbf{Each agent is holonomy-free here.} No rotation fixes $1$, so
        $\mathrm{Stab}_{P}(1)=\{e\}$; the transposition $(12)$ moves $1$, so
        $\mathrm{Stab}_{P'}(1)=\{e\}$. Both factor holonomies at $1$ are trivial. The
        guess predicts the composite holonomy is trivial too.
  \item \textbf{The composite is not.} In $S_{3}$ the stabiliser of $1$ is
        $\mathrm{Stab}_{G}(1)=\langle(23)\rangle\cong\mathbb Z/2$ --- and $(23)$
        factors as $(132)\cdot(12)$, a rotation times the transposition, \emph{neither
        of which fixes $1$}. Read it as a loop: start at $1$, apply the $P'$-leg $(12)$
        to reach $2$, then apply the $P$-leg $(132)$ to return to $1$.
        \[
           1\ \xrightarrow{\ (12)\ }\ 2\ \xrightarrow{\ (132)\ }\ 1 .
        \]
        A dispatch and a return. Neither step is a stabiliser; their composite is. The
        composite carries a genuine $\mathbb Z/2$ of holonomy on the fibre
        $\mathcal C_{1}$ --- a nontrivial fibre automorphism --- that \emph{neither
        agent possesses}.
\end{itemize}
This is reentrancy created by orchestration. The equality-guess would have hidden it
exactly; its failure is the discovery.
\prov{MacBeth (\texttt{2026-08-12-holonomy-composition-zs-bridge.md}, part (b),
refuting the equality; the containment holds always). The $S_{3}$ witness is
Lean-verified, axiom-free (\texttt{EmergentHolonomy.lean}), not committed to this tree.
Exhaustively swept: the inclusion held in all $448$ point-checks over the exact
factorisations of $S_{3},S_{4},A_{4},D_{4},\mathbb Z/2{\times}\mathbb Z/2$ and was
\emph{proper} in $268$ --- including the abelian $\mathbb Z/2{\times}\mathbb Z/2$, so
emergence is not a nonabelian artefact.}
\end{Example}

\subsection*{The measure of emergence: where the orbits cross}

How much holonomy does orchestration synthesise at $s$? The answer is a single positive
integer with a picture attached. Each agent has, from $s$, a \emph{reachable-state
orbit}: $s\cdot P$ under the first, $s\cdot P'$ under the second. In the $S_{3}$ example
$s\cdot P=\{1,2,3\}$ (rotations reach everything) while $s\cdot P'=\{1,2\}$; the two
orbits cross at $\{1,2\}$, two points. That $2$ is exactly the invariant $h(1)$ we are
about to define, and here --- where both factor isotropies are trivial --- it is also the
order of the emergent holonomy group.

\begin{Theorem}[Emergent holonomy $=$ meeting points]\label{thm:meeting-points}
For an exact factorisation $G=P\cdot P'$ acting on $S$ --- the internal realisation of the
composite threading, $s\cdot g=(s\downarrow p)\downarrow' p'$ for $g=pp'$ --- write
$U=\mathrm{Stab}_{G}(s)$, $A=\mathrm{Stab}_{P}(s)$, $B=\mathrm{Stab}_{P'}(s)$. The composite ``dispatch-and-return''
loops at $s$, taken up to the two factor symmetries, are in canonical bijection with the
crossing points of the two reachable-state orbits; and their number
\[
   h(s)\ :=\ |A\backslash U/B|
   \ =\ \frac{|\mathrm{Stab}_{G}(s)|}{|\mathrm{Stab}_{P}(s)|\cdot|\mathrm{Stab}_{P'}(s)|}
   \ =\ \bigl|(s\cdot P)\cap(s\cdot P')\bigr|
\]
is a positive integer, equal by all three formulas. Moreover
\[
   h(s)=1
   \quad\Longleftrightarrow\quad
   (s\cdot P)\cap(s\cdot P')=\{s\}
   \quad\Longleftrightarrow\quad
   \mathrm{Stab}_{G}(s)=\mathrm{Stab}_{P}(s)\cdot\mathrm{Stab}_{P'}(s):
\]
the two orbits meet \emph{only} at $s$ exactly when the composite holonomy is honest ---
every composite loop splits into one loop of each agent, and the guess above holds after
all.
\prov{MacBeth (\texttt{2026-08-13-emergent-holonomy-meeting-points.md}, Thm.~6,
Cor.~7); verified with $0$ mismatches across $9$ groups incl.\ $A_{5}$, $2594$
point-checks.}
\end{Theorem}

The middle formula deserves a second look, because \emph{a priori it is not even an
integer}: it divides the size of a subgroup by the sizes of two others that need not
have anything to do with each other. The next box says why it always comes out whole.

\begin{teachbox}[Why the naive ratio is a whole number]
The one fact behind Theorem~\ref{thm:meeting-points} is a property of exact
factorisations alone: for \emph{every} $g\in G$,
\[
   P\cap gP'g^{-1}=\{e\}.
\]
No action, no point $s$ --- just that $P,P'$ are complementary subgroups
($P\cap P'=\{e\}$, $|P||P'|=|G|$). It follows that inside $U=\mathrm{Stab}_{G}(s)$ every
double coset $AgB$ has the \emph{same} size, $|A|\cdot|B|$; so their number is exactly
$|U|/(|A||B|)$, an integer, and it counts the loops-up-to-factor-symmetry. The geometric
identity $=|(s\cdot P)\cap(s\cdot P')|$ is the payoff: the waypoint $t$ of a loop
$s\to t\to s$ is a crossing point of the two orbits, and distinct emergent loops pin
distinct crossings. One honest double coset (waypoint $s$) is the aligned part where both
legs fix $s$; the other $h(s)-1$ are the synthesised reentrant loops.
\prov{Lemma~1 of \texttt{2026-08-13-emergent-holonomy-meeting-points.md}; the
disjointness checked $41\,064$ times, $0$ violations.}
\end{teachbox}

So the emergent holonomy is not a mystery to be feared but a number to be computed:
$h(s)-1$ is the count of reentrant loops that orchestration manufactures at $s$, one for
each state where the two agents' reachable orbits cross beyond the obvious. An auditor
of a composed system --- a supply chain, a pair of contracts, two orchestrated agents ---
reads it straight off the two state actions. No cohomology is needed to \emph{detect}
emergence; only, as we now see, to classify the case where there is none.

\subsection*{The clean case, and the class that certifies it}

When is the composite a clean product --- when do the two agents' holonomies sit on
independent registers of the fibre, ignoring each other? By
Theorem~\ref{thm:meeting-points} the precondition is $h(s)=1$: the orbits meet only at
$s$, so the composite isotropy $E:=\mathrm{Stab}_{G}(s)$ is genuinely the Zappa--Sz\'ep
product $A\bowtie B$ of the two factor isotropies. But being a product of the two groups
is not yet being a product of their \emph{representations}. That is the last obstruction,
and it is one degree up.

\begin{Theorem}[The split obstruction is a degree-$2$ class]\label{thm:emergent-h2}
Suppose $s$ is aligned ($h(s)=1$) with $A$ abelian and normal in $E$. Then $E$ is an
extension $1\to A\to E\to B\to 1$, classified by a class $[\omega]\in H^{2}(B;A)$; and
when the conjugation action is trivial (in particular for the $\mathbb Z/2$ witness),
\[
   [\omega]=0
   \ \Longleftrightarrow\ E\cong A\times B
   \ \Longleftrightarrow\
   \text{every composite holonomy is \emph{unentangled}}
   \ (\rho(A),\rho(B)\text{ commute in }\mathrm{Aut}(\mathcal C_{s})).
\]
A nonzero $[\omega]$ is the \emph{abstract} non-split extension (e.g.\ $\mathbb Z/4$), a
composite holonomy whose two agent-parts cannot be disentangled.
\prov{MacBeth (\texttt{2026-08-12-holonomy-composition-zs-bridge.md}, part (c$'$)),
resting on the Eilenberg--Mac Lane / Baues--Wirsching extension classification
\cite{BW85}. Scope: aligned, abelian, normal $A$ --- the honest domain; the nonabelian
case is deliberately not claimed. \emph{Realisability caveat}
(\texttt{2026-08-14-two-omega-sites-isotropy-restriction.md}, Thm.~2.5): within this very
hypothesis an aligned point-stabiliser with $A$ normal is a semidirect product
$E=A\rtimes B$ and hence \emph{split}, so $[\omega]=0$ always; the nonzero class classifies
an abstract extension but is not realised by an aligned Zappa--Sz\'ep point-stabiliser
(see the teachbox below).}
\end{Theorem}

The smallest instance is the whole story in miniature. Take both isotropies to be
$\mathbb Z/2$, acting trivially, so $H^{2}(\mathbb Z/2;\mathbb Z/2)\cong\mathbb F_{2}=
\{0,\varepsilon\}$:
\begin{center}
\renewcommand{\arraystretch}{1.3}
\begin{tabular}{@{}cllc@{}}
\toprule
$[\omega]$ & extension $E$ classified by $[\omega]$ & the two agents' holonomies & composite\\
\midrule
$0$ & $\mathbb Z/2\times\mathbb Z/2$ & two $\mathbb Z/2$ representations & \emph{unentangled} (they commute)\\
$\varepsilon$ & $\mathbb Z/4$ (cyclic, non-split) & two $\mathbb Z/2$ representations & \emph{entangled} order-$4$ automorphism\\
\bottomrule
\end{tabular}
\end{center}
The two agents' holonomies are the same in both rows --- two copies of $\mathbb Z/2$.
What differs is whether they combine into two independent commuting symmetries of the
fibre or braid into a single order-$4$ one, and that difference is a single bit,
$[\omega]$. This is the clean resolution of a degree gap that dogged the whole story: the
factor holonomies are degree-$1$ data (representations), and \emph{whether they assemble
without entanglement} is decided by a degree-$2$ class --- never by an equality of the
two.

\begin{teachbox}[Same $\mathbb Z/2$, two different sites --- and why the naive identification fails]
Twice now an $H^{2}$ class with a $\mathbb Z/2$ generator has decided whether a
Zappa--Sz\'ep composite is clean, and both times the generator is the \emph{same} token
bit $\varepsilon$ read off the \emph{same} composition data. The first was the handoff
obstruction of Theorem~\ref{thm:h2}, $[\omega_{h}]\in H^{2}(\mathrm{Sk}_{\mathcal C};
\mathcal D)$ --- Baues--Wirsching cohomology of the \emph{handoff category}, whose objects
are the roles a token passes through, deciding whether two directed systems compose at
all. The second is the stabiliser obstruction just above,
$[\omega_{\mathrm{st}}]\in H^{2}(B;A)$ --- ordinary group cohomology of a single
\emph{point-stabiliser}, whose one object is a group, deciding whether an already-formed
composite holonomy splits. Same cyclic group, same generator, same underlying bit. Surely
they are one class seen twice?

They are not, and the reason is worth the box. Baues--Wirsching cohomology
\emph{generalises} group cohomology: a one-object category \emph{is} a group $G$, and there
$H^{2}_{\mathrm{BW}}$ is exactly $H^{2}(G;-)$. For a general category $\mathcal C$ every
object $x$ still carries its automorphism group $\mathrm{Aut}(x)$, and \emph{restriction to
isotropy} is a natural map
\[
   i_{x}^{*}\colon H^{2}_{\mathrm{BW}}(\mathcal C;\mathcal D)\longrightarrow
   H^{2}\bigl(\mathrm{Aut}(x);\mathcal D(x)\bigr)
\]
that reads off the part of a category class living on the \emph{loops} at $x$. It is the
only map of the right shape --- degree $2$, correct variance --- that could turn the
handoff class (a category class) into the stabiliser class (a group class). So compute it.
The handoff category $\mathrm{Sk}_{\mathcal C}$ is \emph{automorphism-rigid}: its objects
are the roles $\mathrm{Sup}\to\mathrm{Wk}\rightrightarrows\mathrm{Rt}$, and not one of them
has a non-identity self-map --- the token that could have been a loop was pushed into the
\emph{coefficient} $\mathcal D(\mathrm{Sup})=\mathbb Z/2$, not into an arrow. Every isotropy
group is trivial, $H^{2}(1;-)=0$, so $i_{x}^{*}$ is the \emph{zero map at every object}. It
annihilates $[\omega_{h}]$ --- indeed every class on $\mathrm{Sk}_{\mathcal C}$ --- because
the handoff generator is supported on the length-$2$ chain
$\mathrm{Sup}\to\mathrm{Wk}\to\mathrm{Rt}$ through three \emph{distinct} objects: an
off-diagonal, nerve-level class, exactly the part isotropy restriction throws away. The
stabiliser class, by contrast, is a pure loop class at one object. The two $\mathbb F_{2}$'s
coincide numerically, \emph{not} by any restriction map.

Nor does a change of viewpoint rescue the identification, and the collapse is two-sided.
The classes live over \emph{different base categories} --- the handoff category
$\mathrm{Sk}_{\mathcal C}$ (objects: roles) and the state-action category
$\mathbb A(\downarrow_{\bowtie})$ (objects: states) --- so relating them would need a common
refinement bridging the two sites. No comparison functor
$F\colon\mathbb A(\downarrow_{\bowtie})\to\mathrm{Sk}_{\mathcal C}$ provides one: a functor
sends loops to loops, $\mathrm{Sk}_{\mathcal C}$ has none but identities, so every
pulled-back class is again isotropy-invisible ($i_{s}^{*}F^{*}[\omega_{h}]=0$ for all $F$).
In the cases one can check by hand --- one-object group targets, comparison functors in
either direction --- the bridge is provably obstructed; the fully general many-object
cospan is the one honest open edge. And the stabiliser site collapses from its own side:
an aligned point-stabiliser with $A$ normal is a semidirect product $E=A\rtimes B$, hence
\emph{split}, so $[\omega_{\mathrm{st}}]=0$ wherever it is even defined. The nonzero
$\mathbb Z/4$ row of the table above is thus the \emph{abstract} extension classified by
$\varepsilon$, not something aligned Zappa--Sz\'ep geometry realises --- the stabiliser
class is forced to $0$ exactly where the naive identification wants it nonzero.

The discipline is one this book has had to learn before: two degree-$2$ obstructions with
the same cyclic group and the same generator need \emph{not} be the same class. Check the
support. The handoff bit and the stabiliser bit are two images of one Zappa--Sz\'ep datum
$\varepsilon$ under two \emph{incomparable} maps out of category cohomology --- an
off-diagonal (nerve) coordinate and a pure-isotropy coordinate --- and $i_{x}^{*}$ is
precisely the projection that keeps the second and kills the first. It is a rhyme, not an
identity; and knowing \emph{why} it rhymes is better than mistaking it for a repetition.
\prov{MacBeth (\texttt{2026-08-14-two-omega-sites-isotropy-restriction.md}, Parts~1, 1.5,
2), resting on the isotropy-restriction map of Baues--Wirsching \cite{BW85}. Open edge: the
fully general many-object cospan $\mathrm{Sk}_{\mathcal C}\to\mathcal C\leftarrow
\mathbb A(\downarrow_{\bowtie})$.}
\end{teachbox}

\paragraph{What this does to the table.} It does not collapse it. The three obstructions
of \S\ref{sec:threemodes} --- a Baues--Wirsching class, nothing at all, a branching flag
--- remain three different kinds of mathematics, and there is still no master theorem
that makes them one. What this section shows is subtler and, to the author, better:
\emph{two of the rows are not disjoint problems but one problem seen from two sides.}
Composing two agents from the State row \emph{is} a Zappa--Sz\'ep composition from the
Directed row, and on that one object both invariants live at once --- the isotropy
holonomy of the state mode and the $[\omega]$ of the directed mode, meeting at the
point-stabiliser. The table's rows are axes, not walls; a single composition can run
along two of them, and when it does, the modes talk to each other.

\paragraph{For the grant.} Here is the earned sentence, the one the three composition
chapters were built to reach: \emph{unprotected orchestration can synthesise holonomy
the parts lack, and a degree-$2$ cohomology class certifies when the composite is a
clean product of the parts.} It comes with an operational reading a practitioner can
use without any cohomology at all --- to find where a composed system carries reentrancy
its components do not, count the states where the two components' reachable-state orbits
cross. Where they cross only trivially, the system is holonomy-clean and the finer
question of entanglement is a single computable bit. Where they cross more, orchestration
has manufactured symmetry, and the audit has found it. Fair is foul: two reentrancy-free
agents, composed, can make reentrancy --- and the mathematics says exactly where, and
exactly how much.

\chapter{Functors over the category of containers (Phase 2 --- outline)}
\label{ch:phase2}

This chapter is a stub sketching Neil's proposed Phase 2 programme. The earlier chapters
studied containers and the maps \emph{between} them, and Chapter~\ref{ch:moncomon} built the
two canonical endofunctors --- the free monad and the cofree comonad. Phase 2 studies
\emph{endofunctors of $\Cont$} in general: operations $\Cont\to\Cont$ (equivalently
$\Poly\to\Poly$) and the structures they carry.

\paragraph{What a functor $\Cont\to\Cont$ is.} It assigns to each container a
container, functorially in container morphisms. The programme is to identify the
natural such operations and the (co)monad structures they support.

\paragraph{Candidate examples (to be developed, not yet results).}
\begin{itemize}[nosep]
  \item \emph{The derivative / one-hole-context functor} $\partial$: sends
        $\cont SP$ to the container of ``a shape with one position singled out''.
        For lists, $\partial(\sum_n X^n)$ is the zipper. The combinatorial
        derivative is classical for species; its container-level naturality and
        its interaction with $\triangleleft$ are the targets here.
  \item \emph{The free monad and cofree comonad functors} on $(\Poly,\triangleleft)$:
        the syntax/behaviour pair. These are no longer stubs --- they are developed as
        results in their own right in Chapter~\ref{ch:moncomon}, and are recalled here only
        as the leading examples of the Phase 2 programme.
\end{itemize}

\appendix
\chapter{Reading \& citation tracker}
\label{app:tracker}

The status column records how thoroughly MacBeth has engaged each source: ``read
in full'' (worked through, definitions and key proofs internalised), ``skimmed''
(read for the relevant statement), or ``to-read''.

\begin{center}
\renewcommand{\arraystretch}{1.25}
\begin{tabular}{@{}p{5.4cm}p{3.1cm}p{2.4cm}p{2.6cm}@{}}
\toprule
\textbf{Paper} & \textbf{arXiv / ref} & \textbf{Status} & \textbf{Used in} \\
\midrule
Abbott--Altenkirch--Ghani, \emph{Containers: constructing strictly positive types} (2005)
  & TCS 342 & read in full & Ch.\,1--3 \\
Abbott--Altenkirch--Ghani, \emph{Categories of containers} (2003)
  & FoSSaCS, LNCS 2620 & Prop.\,4.3/Ex.\,4.4 verified & Ch.\,3 \\
Gambino--Kock, \emph{Polynomial functors and polynomial monads} (2009/13)
  & arXiv:0906.4931 & read in full & Ch.\,3 \\
Carboni--Johnstone, \emph{Connected limits, familial representability, Artin glueing} (1995)
  & MSCS 5 & read in full & Ch.\,3 \\
Diers, \emph{Cat\'egories localisables} (1977)
  & Th\`ese, Paris VI & skimmed & Ch.\,3 \\
Ahman--Chapman--Uustalu, \emph{When is a container a comonad?} (2014)
  & LMCS 10(3) & read in full & Ch.\,5 \\
Ahman--Uustalu, \emph{Directed containers as categories} (2016)
  & arXiv:1604.01187 & read in full & Ch.\,5 \\
Ahman--Bauer, \emph{Comodule representations of second-order functionals} (2024)
  & arXiv:2409.17664 & read in full & Ch.\,7 \\
Le--Siqueira--Kremnitzer, \emph{Free PLTL algebras \& hyperdoctrines} (2025)
  & Topos Inst.\ blog & read in full & Ch.\,7 \\
Spivak, \emph{Functorial aggregation / Polynomial functors} (2021)
  & arXiv:2111.10968 & read in full & Ch.\,4--5 \\
Niu--Spivak, \emph{Polynomial Functors: A Mathematical Theory of Interaction} (2023)
  & arXiv:2312.00990 & read in full & Ch.\,6 \\
Spivak, \emph{A reference for categorical structures on \textup{Poly}} (2022)
  & arXiv:2202.00534 & read in full & Ch.\,6 \\
Shapiro--Spivak, \emph{Structures on categories of polynomials} (2023)
  & arXiv:2305.00167 & read in full & Ch.\,5 \\
Shapiro--Spivak, \emph{A polynomial construction of nerves} (2024)
  & arXiv:2405.13157 & skimmed & Ch.\,5 \\
Clarke, \emph{Internal lenses as functors and cofunctors} (2020)
  & arXiv:2009.06835 & read in full & Ch.\,5 \\
Clarke, \emph{Delta lenses as coalgebras for a comonad} (2022)
  & arXiv:2108.00390 & read in full & Ch.\,5 \\
Rosebrugh--Wood, \emph{Distributive laws and factorization} (JPAA 175, 2002)
  & JPAA 175 & read in full & Ch.\,6 \\
Baues--Wirsching, \emph{Cohomology of small categories} (1985)
  & JPAA 38, 187--211 & read in full & Ch.\,6 \\
\emph{(modern BW account)} (2025)
  & arXiv:2508.00727 & skimmed & Ch.\,6 \\
Aguiar, \emph{Internal categories and quantum groups} (1997)
  & PhD, Cornell & skimmed & Ch.\,5 \\
Uustalu--Vene, \emph{Comonadic notions of computation} (2008)
  & ENTCS 203(5) & read for store comonad & Ch.\,7 \\
Capucci--Myers, \emph{Contextads as wreaths} (2024)
  & arXiv:2410.21889 & read in full & Ch.\,7 \\
Capucci--Gavranovi\'c--Hedges--Rischel, \emph{Categorical cybernetics / Para} (2021)
  & arXiv:2105.06332 & to-read (deep-read TODO) & Ch.\,7 \\
Spivak, \emph{Categories by Kan extension} (2025)
  & arXiv:2503.21974 & skimmed (abstract) & Ch.\,7 \\
\bottomrule
\end{tabular}
\end{center}

\begin{thebibliography}{99}

\bibitem{AAG05} M.~Abbott, T.~Altenkirch, N.~Ghani.
\emph{Containers: constructing strictly positive types}.
Theoretical Computer Science \textbf{342} (2005) 3--27.

\bibitem{AAGcat} M.~Abbott, T.~Altenkirch, N.~Ghani.
\emph{Categories of containers}.
Proc.\ FoSSaCS 2003, LNCS \textbf{2620}, 23--38.

\bibitem{GK} N.~Gambino, J.~Kock.
\emph{Polynomial functors and polynomial monads}.
\href{https://arxiv.org/abs/0906.4931}{arXiv:0906.4931};
Math.\ Proc.\ Cambridge Philos.\ Soc.\ \textbf{154} (2013) 153--192.

\bibitem{CJ} A.~Carboni, P.~Johnstone.
\emph{Connected limits, familial representability and Artin glueing}.
Mathematical Structures in Computer Science \textbf{5} (1995) 441--459.
(Corrigenda: \emph{MSCS} \textbf{14} (2004) 185--187.)

\bibitem{Diers} Y.~Diers.
\emph{Cat\'egories localisables}.
Th\`ese de doctorat, Universit\'e Paris VI, 1977.

\bibitem{ACU14} D.~Ahman, J.~Chapman, T.~Uustalu.
\emph{When is a container a comonad?}
Logical Methods in Computer Science \textbf{10}(3), 2014.

\bibitem{AU16} D.~Ahman, T.~Uustalu.
\emph{Directed containers as categories}.
Proc.\ MSFP 2016, EPTCS \textbf{207}, 89--98.
\href{https://arxiv.org/abs/1604.01187}{arXiv:1604.01187}.

\bibitem{AhmanUustalu13} D.~Ahman, T.~Uustalu.
\emph{Update monads: cointerpreting directed containers}.
In \emph{19th Int.\ Conf.\ on Types for Proofs and Programs (TYPES 2013)},
LIPIcs \textbf{26}, Schloss Dagstuhl, 2014, pp.\ 1--23.

\bibitem{CBP} E.~ter Horst, S.~Mahadevan, J.~D.~Zambrano.
\emph{Categorical belief propagation: sheaf-theoretic inference via descent and holonomy}.
\href{https://arxiv.org/abs/2601.04456}{arXiv:2601.04456} (2026).

\bibitem{AhmanBauer24} D.~Ahman, A.~Bauer.
\emph{Comodule representations of second-order functionals}.
\href{https://arxiv.org/abs/2409.17664}{arXiv:2409.17664};
J.\ Log.\ Algebraic Methods Program.\ \textbf{146} (2025).

\bibitem{ToposPLTL} Q.~Le, J.~Siqueira, K.~Kremnitzer.
\emph{Free PLTL algebras and a coalgebraic extension of hyperdoctrines}.
Topos Institute Blog, 26 September 2025.
\href{https://topos.institute/blog/2025-09-26-free-pltl-algebras-and-hyperdoctrines/}{topos.institute}.

\bibitem{Spivak21} D.~I.~Spivak.
\emph{Functorial aggregation} (polynomial functors).
\href{https://arxiv.org/abs/2111.10968}{arXiv:2111.10968}, 2021.

\bibitem{NiuSpivak23} N.~Niu, D.~I.~Spivak.
\emph{Polynomial Functors: A Mathematical Theory of Interaction}.
\href{https://arxiv.org/abs/2312.00990}{arXiv:2312.00990}, 2023.

\bibitem{SpivakRef} D.~I.~Spivak.
\emph{A reference for categorical structures on \textup{Poly}}.
\href{https://arxiv.org/abs/2202.00534}{arXiv:2202.00534}, 2022.

\bibitem{DJN} B.~Dorta, M.~Jarvis, N.~Niu.
\emph{Monoidal structures on generalized polynomial categories}.
\href{https://arxiv.org/abs/2305.05655}{arXiv:2305.05655}, 2023.

\bibitem{DUV} G.~De~Pascalis, T.~Uustalu, N.~Veltr\`i.
\emph{Monoid structures on indexed containers}.
\href{https://arxiv.org/abs/2509.25879}{arXiv:2509.25879}, 2025.

\bibitem{SS23} B.~Shapiro, D.~I.~Spivak.
\emph{Structures on categories of polynomials}.
\href{https://arxiv.org/abs/2305.00167}{arXiv:2305.00167}, 2023.

\bibitem{SS2405} B.~Shapiro, D.~I.~Spivak.
\emph{A polynomial construction of nerves for higher categories}.
\href{https://arxiv.org/abs/2405.13157}{arXiv:2405.13157}, 2024.

\bibitem{Clarke20} B.~Clarke.
\emph{Internal lenses as functors and cofunctors}.
Proc.\ ACT 2019, EPTCS \textbf{323}, 183--195.
\href{https://arxiv.org/abs/2009.06835}{arXiv:2009.06835}.

\bibitem{Clarke22} B.~Clarke.
\emph{Delta lenses as coalgebras for a comonad}.
\href{https://arxiv.org/abs/2108.00390}{arXiv:2108.00390}, 2022.

\bibitem{Aguiar} M.~Aguiar.
\emph{Internal categories and quantum groups}.
PhD thesis, Cornell University, 1997.

\bibitem{RW} R.~Rosebrugh, R.~J.~Wood.
\emph{Distributive laws and factorization}.
Journal of Pure and Applied Algebra \textbf{175} (2002) 327--353.

\bibitem{BW85} H.-J.~Baues, G.~Wirsching.
\emph{Cohomology of small categories}.
Journal of Pure and Applied Algebra \textbf{38} (1985) 187--211.

\bibitem{BWmodern} \emph{(Modern account of Baues--Wirsching cohomology)}.
\href{https://arxiv.org/abs/2508.00727}{arXiv:2508.00727}, 2025.

\bibitem{UustaluVene08} T.~Uustalu, V.~Vene.
\emph{Comonadic notions of computation}.
Electronic Notes in Theoretical Computer Science \textbf{203}(5) (2008) 263--284.

\bibitem{CapucciMyers} M.~Capucci, D.~J.~Myers.
\emph{Contextads as wreaths; Kleisli, Para, and Span constructions as wreath products}.
\href{https://arxiv.org/abs/2410.21889}{arXiv:2410.21889}, 2024.

\bibitem{Gavranovic} M.~Capucci, B.~Gavranovi\'c, J.~Hedges, E.~F.~Rischel.
\emph{Towards foundations of categorical cybernetics}.
Proc.\ ACT 2021, EPTCS \textbf{372}, 235--248;
see also B.~Gavranovi\'c, \emph{Fundamental components of deep learning}, PhD thesis, 2024.

\bibitem{SpivakKanExt} D.~I.~Spivak.
\emph{Categories by Kan extension}.
\href{https://arxiv.org/abs/2503.21974}{arXiv:2503.21974}, 2025.

\bibitem{Beck69} J.~Beck.
\emph{Distributive laws}.
In: Seminar on Triples and Categorical Homology Theory, Lecture Notes in Mathematics \textbf{80},
Springer (1969) 119--140.

\bibitem{PowerWatanabe02} J.~Power, H.~Watanabe.
\emph{Combining a monad and a comonad}.
Theoretical Computer Science \textbf{280} (2002) 137--162.

\bibitem{BrzezinskiMajid98} T.~Brzezi\'nski, S.~Majid.
\emph{Coalgebra bundles}.
Communications in Mathematical Physics \textbf{191} (1998) 467--492.

\bibitem{KRU20} S.~Katsumata, E.~Rivas, T.~Uustalu.
\emph{Interaction laws of monads and comonads}.
Proc.\ FoSSaCS 2020, LNCS \textbf{12077}, 356--375;
\href{https://arxiv.org/abs/1912.13477}{arXiv:1912.13477}.

\bibitem{vonGlehn18} T.~von Glehn.
\emph{Polynomials and models of type theory}.
PhD thesis, University of Cambridge, 2015;
see also Theory and Applications of Categories \textbf{33} (2018).

\bibitem{Streicher} T.~Streicher.
\emph{Fibered Categories \`a la Jean B\'enabou}.
Lecture notes; \href{https://arxiv.org/abs/1801.02927}{arXiv:1801.02927} (orig.\ 1999).

\bibitem{Jacobs99} B.~Jacobs.
\emph{Categorical Logic and Type Theory}.
Studies in Logic and the Foundations of Mathematics \textbf{141}, Elsevier, 1999.

\bibitem{Street72} R.~Street.
\emph{The formal theory of monads}.
Journal of Pure and Applied Algebra \textbf{2} (1972) 149--168.

\bibitem{Hermida93} C.~Hermida.
\emph{Fibrations, logical predicates and indeterminates}.
PhD thesis, University of Edinburgh, 1993.

\bibitem{HermidaJacobs98} C.~Hermida, B.~Jacobs.
\emph{Structural induction and coinduction in a fibrational setting}.
Information and Computation \textbf{145}(2) (1998) 107--152. DOI 10.1006/inco.1998.2725.

\bibitem{OrestisAgda} Orestis.
\emph{Effects} --- an Agda formalisation of monad and comonad liftings to containers.
Unpublished; personal communication via N.~Ghani, 2026. Module \texttt{Containers.Effects}
(files \texttt{CoLift.agda}, \texttt{Lift.agda}, \texttt{BiLift.agda}).

\bibitem{Katsumata13} S.~Katsumata.
\emph{Relating computational effects by $\top\top$-lifting}.
Information and Computation \textbf{222} (2013) 228--246.

\bibitem{PradicPrice} Pradic, Price.
\emph{Weihrauch problems as containers}.
\href{https://arxiv.org/abs/2501.17250}{arXiv:2501.17250}, 2025.

\bibitem{Fairbanks25} A.~D.~Fairbanks.
Answer to \emph{``Examples of non-polynomial comonads on $\mathrm{Set}$?''}
MathOverflow, question~457580, answer posted 12~January 2025.
\href{https://mathoverflow.net/a/485788}{mathoverflow.net/a/485788}.

\end{thebibliography}

\end{document}
